% Trigonométrie : formules, équations et inéquations — Mathématiques 1ère E — Cours signé OnBuch+
% Programme officiel MINESEC. Compiler avec : tectonic N03-trigonometrie.tex
\newcommand{\DOCMATIERE}{Mathématiques}
\newcommand{\DOCNIVEAU}{1ère E}
\newcommand{\DOCMODULE}{Relations et opérations fondamentales dans l'ensemble des nombres réels}
\newcommand{\DOCLECON}{N03}
\newcommand{\DOCTITRE}{Trigonométrie : formules, équations et inéquations}
% OnBuch+ — préambule « cours signés » ENRICHI (compilé avec tectonic / XeLaTeX)
% Système visuel « Pop » d'OnBuch+ : fond crème (jamais blanc), bordures encre
% épaisses, ombres « dures » décalées, titres Archivo Black en capitales.
% Version MATHÉMATIQUES : graphes (pgfplots/tikz), boîtes pédagogiques
% (définition, propriété, méthode, attention, le savais-tu, exercice résolu,
% exercice, TP, bilan), page de garde et signature OnBuch+ ancrée en bas.
%
% Macros attendues dans main.tex AVANT \input{preamble} :
%   \DOCMATIERE  \DOCNIVEAU  \DOCMODULE  \DOCLECON (numéro)  \DOCTITRE
\documentclass[11pt]{article}

\usepackage{fontspec}
\usepackage[french]{babel}
\usepackage[a4paper,margin=2cm,top=3.4cm,bottom=2.8cm,headheight=1.4cm,footskip=1.3cm]{geometry}
\usepackage{amsmath,amssymb,amsfonts}
\usepackage{mathtools} % \xmapsto, \coloneqq... souvent utilisés par les modèles
\usepackage{cancel} % simplifications barrées
% \abs{z} (module, valeur absolue) et \norm{u} (norme), utilisés par les modèles
\providecommand{\abs}{}\let\abs\relax\DeclarePairedDelimiter{\abs}{\lvert}{\rvert}
\providecommand{\norm}{}\let\norm\relax\DeclarePairedDelimiter{\norm}{\lVert}{\rVert}
\usepackage[table]{xcolor} % [table] : fournit aussi \rowcolors (alternance de lignes)
\usepackage{pagecolor}
\usepackage{tikz}
\usetikzlibrary{arrows.meta,positioning,calc,shapes.geometric,shapes.misc,shapes.arrows,decorations.pathmorphing,decorations.pathreplacing,decorations.markings,patterns,shadows,mindmap,trees,fit,backgrounds,matrix,intersections,angles,quotes}
\usepackage{pgfplots}
\usepackage{pgf-pie} % diagrammes circulaires \pie (statistiques)
\pgfplotsset{compat=1.18}
\usepgfplotslibrary{fillbetween,groupplots} % aires entre courbes (calcul intégral)
\usepackage{enumitem}
\usepackage{fancyhdr}
% [most] charge toutes les bibliothèques tcolorbox (listings, minted, raster,
% documentation...) et gonfle inutilement la mémoire TeX sur un document long
% (~300 boîtes) : on ne charge que ce qui est réellement utilisé.
\usepackage[skins,breakable]{tcolorbox}
\usepackage{graphicx}
\usepackage{colortbl}
\usepackage{booktabs}
\usepackage{tabularx}
\usepackage{diagbox} % case d'en-tête diagonale des tableaux à double entrée
% Colonnes extensibles centrée / gauche / droite (C, L, R), souvent utilisées par les modèles
\newcolumntype{C}{>{\centering\arraybackslash}X}
\newcolumntype{L}{>{\raggedright\arraybackslash}X}
\newcolumntype{R}{>{\raggedleft\arraybackslash}X}
\usepackage{array}
\usepackage{multicol}
\usepackage{siunitx}
% parse-numbers=false : accepte aussi \SI{2,5 \times 10^{-3}}{...}
\sisetup{locale=FR,output-decimal-marker={,},group-minimum-digits=4,parse-numbers=false}
\DeclareSIUnit\ohm{\ensuremath{\symup{\Omega}}} % U+2126 absent de la police
\usepackage{qrcode}
% \degree : commande fréquemment utilisée par les modèles pour un angle ou une
% température en dehors de \SI{}{} (non fournie par les paquets chargés).
\providecommand{\degree}{^{\circ}}
% \qed (fin de démonstration), utilisé par les modèles sans amsthm
\providecommand{\qed}{\hfill$\square$}
% \barre : double barre des valeurs interdites dans un tableau de variations (tkz-tab)
\providecommand{\barre}{\ensuremath{\|}}
% environnement proof (amsthm non chargé) utilisé par les modèles
\ifdefined\proof\else\newenvironment{proof}[1][Démonstration]{\par\noindent\textit{#1.}\ }{\qed\par}\fi
\usepackage{lastpage}
\usepackage{hyperref}
\hypersetup{hidelinks}

% ============================================================
% Palette « Pop » OnBuch+ — mêmes hex que la classe P (plus_ui.dart)
% ============================================================
\definecolor{poporange}{HTML}{F59321}
\definecolor{popdark}{HTML}{E07A0C}
\definecolor{popcream}{HTML}{FFF6E8}
\definecolor{popink}{HTML}{1C1714}
\definecolor{poppurple}{HTML}{7A5AE0}
\definecolor{popgreen}{HTML}{2E9E6A}
\definecolor{poppink}{HTML}{E0457B}
\definecolor{popblue}{HTML}{2F7DE1}
\definecolor{popgold}{HTML}{C98A12}
\definecolor{popmuted}{HTML}{8A7F76}
\definecolor{popline}{HTML}{EADFD2}
\definecolor{popgray}{HTML}{8A7F76}
\colorlet{popgrey}{popgray}
% Alias tolérants (le modèle nomme parfois les couleurs différemment)
\colorlet{popwhite}{white}
\colorlet{popblack}{popink}
\colorlet{popred}{poppink}
\colorlet{popyellow}{popgold}
% Teintes claires pour fonds de tableaux / zones
\colorlet{popblueL}{popblue!10!white}
\colorlet{poporangeL}{poporange!14!white}
\colorlet{popgreenL}{popgreen!12!white}
\colorlet{poppurpleL}{poppurple!12!white}
\colorlet{poppinkL}{poppink!12!white}
\colorlet{popgoldL}{popgold!14!white}

\pagecolor{popcream}

% ============================================================
% Typographie
% ============================================================
\defaultfontfeatures{Ligatures=TeX}
\setmainfont{PlusJakartaSans-Medium.ttf}[Path=../../../fonts/,BoldFont=PlusJakartaSans-Bold.ttf]
% Maths en sans-serif (Fira Math) avec chiffres et lettres droites de Plus
% Jakarta, pour que les formules se fondent dans le texte.
\usepackage[math-style=ISO,bold-style=ISO]{unicode-math}
\setmathfont{FiraMath-Regular.otf}[Path=../../../fonts/]
\setmathfont{PlusJakartaSans-Medium.ttf}[Path=../../../fonts/,range={up/num,up/{Latin,latin},\mathup}]
\usepackage{icomma} % virgule décimale sans espace en mode math
% Caractères mathématiques Unicode absents des polices de texte (Plus Jakarta,
% Archivo Black) : rendus par leur équivalent mathématique, en texte comme en
% formule — sinon ils s'affichent en carré vide (ex. « Arithmétique dans ℤ »).
\usepackage{newunicodechar}
\newunicodechar{ℝ}{\ensuremath{\mathbb{R}}}
\newunicodechar{ℤ}{\ensuremath{\mathbb{Z}}}
\newunicodechar{ℂ}{\ensuremath{\mathbb{C}}}
\newunicodechar{ℕ}{\ensuremath{\mathbb{N}}}
\newunicodechar{ℚ}{\ensuremath{\mathbb{Q}}}
\newunicodechar{𝔻}{\ensuremath{\mathbb{D}}}
\newunicodechar{α}{\ensuremath{\alpha}}
\newunicodechar{β}{\ensuremath{\beta}}
\newunicodechar{γ}{\ensuremath{\gamma}}
\newunicodechar{δ}{\ensuremath{\delta}}
\newunicodechar{ε}{\ensuremath{\varepsilon}}
\newunicodechar{θ}{\ensuremath{\theta}}
\newunicodechar{λ}{\ensuremath{\lambda}}
\newunicodechar{μ}{\ensuremath{\mu}}
\newunicodechar{σ}{\ensuremath{\sigma}}
\newunicodechar{τ}{\ensuremath{\tau}}
\newunicodechar{φ}{\ensuremath{\varphi}}
\newunicodechar{ψ}{\ensuremath{\psi}}
\newunicodechar{ω}{\ensuremath{\omega}}
\newunicodechar{Σ}{\ensuremath{\Sigma}}
% majuscules produites par \MakeUppercase dans les titres (α-aminés -> Α)
\newunicodechar{Α}{\ensuremath{\alpha}}
\newunicodechar{Β}{\ensuremath{\beta}}
\newunicodechar{Γ}{\ensuremath{\Gamma}}
\newunicodechar{Δ}{\ensuremath{\Delta}}
\newunicodechar{Ε}{\ensuremath{\varepsilon}}
\newunicodechar{Θ}{\ensuremath{\Theta}}
\newunicodechar{Λ}{\ensuremath{\Lambda}}
\newunicodechar{Μ}{\ensuremath{\mu}}
\newunicodechar{Τ}{\ensuremath{\tau}}
\newunicodechar{Φ}{\ensuremath{\Phi}}
\newunicodechar{Ψ}{\ensuremath{\Psi}}
\newunicodechar{Ω}{\ensuremath{\Omega}}
\newunicodechar{↦}{\ensuremath{\mapsto}}
\newunicodechar{∈}{\ensuremath{\in}}
\newunicodechar{∉}{\ensuremath{\notin}}
\newunicodechar{∧}{\ensuremath{\wedge}}
\newunicodechar{⇔}{\ensuremath{\Leftrightarrow}}
\newunicodechar{⟺}{\ensuremath{\Longleftrightarrow}}
\newunicodechar{⇒}{\ensuremath{\Rightarrow}}
\newunicodechar{⟹}{\ensuremath{\Longrightarrow}}
\newunicodechar{∘}{\ensuremath{\circ}}
\newunicodechar{∪}{\ensuremath{\cup}}
\newunicodechar{∩}{\ensuremath{\cap}}
\newunicodechar{≡}{\ensuremath{\equiv}}
\newunicodechar{⊂}{\ensuremath{\subset}}
\newunicodechar{⊥}{\ensuremath{\perp}}
\newunicodechar{∃}{\ensuremath{\exists}}
\newunicodechar{∀}{\ensuremath{\forall}}
\newunicodechar{∛}{\ensuremath{\sqrt[3]{\,}}}
\newunicodechar{∜}{\ensuremath{\sqrt[4]{\,}}}
\newunicodechar{⃗}{\ensuremath{{}^{\to}}}
\newunicodechar{ⁿ}{\ifmmode^{n}\else\textsuperscript{\ensuremath{n}}\fi}
\newunicodechar{ˣ}{\ifmmode^{x}\else\textsuperscript{\ensuremath{x}}\fi}
\newunicodechar{ᵇ}{\ifmmode^{b}\else\textsuperscript{\ensuremath{b}}\fi}
\newunicodechar{ʳ}{\ifmmode^{r}\else\textsuperscript{\ensuremath{r}}\fi}
\newunicodechar{ᵘ}{\ifmmode^{u}\else\textsuperscript{\ensuremath{u}}\fi}
\newunicodechar{ʸ}{\ifmmode^{y}\else\textsuperscript{\ensuremath{y}}\fi}
\newunicodechar{ᵃ}{\ifmmode^{a}\else\textsuperscript{\ensuremath{a}}\fi}
\newunicodechar{ᵉ}{\ifmmode^{e}\else\textsuperscript{\ensuremath{e}}\fi}
\newunicodechar{ᵏ}{\ifmmode^{k}\else\textsuperscript{\ensuremath{k}}\fi}
\newunicodechar{ᵖ}{\ifmmode^{p}\else\textsuperscript{\ensuremath{p}}\fi}
\newunicodechar{ᴺ}{\ifmmode^{N}\else\textsuperscript{\ensuremath{N}}\fi}
\newunicodechar{ᶜ}{\ifmmode^{c}\else\textsuperscript{\ensuremath{c}}\fi}
\newunicodechar{ʰ}{\ifmmode^{h}\else\textsuperscript{\ensuremath{h}}\fi}
\newunicodechar{ᵠ}{\ifmmode^{q}\else\textsuperscript{\ensuremath{q}}\fi}
\newunicodechar{ₙ}{\ifmmode_{n}\else\textsubscript{\ensuremath{n}}\fi}
\newunicodechar{ₐ}{\ifmmode_{a}\else\textsubscript{\ensuremath{a}}\fi}
\newunicodechar{ᵢ}{\ifmmode_{i}\else\textsubscript{\ensuremath{i}}\fi}
\newunicodechar{ᵣ}{\ifmmode_{r}\else\textsubscript{\ensuremath{r}}\fi}
\newunicodechar{ₖ}{\ifmmode_{k}\else\textsubscript{\ensuremath{k}}\fi}
\newunicodechar{ᵦ}{\ifmmode_{\beta}\else\textsubscript{\ensuremath{\beta}}\fi}
\newunicodechar{ₓ}{\ifmmode_{x}\else\textsubscript{\ensuremath{x}}\fi}
\newfontfamily\popdisplay{ArchivoBlack-Regular.ttf}[Path=../../../fonts/]
\newfontfamily\poplogo{SpaceGrotesk-Bold.ttf}[Path=../../../fonts/]
\newfontfamily\popsemi{PlusJakartaSans-SemiBold.ttf}[Path=../../../fonts/]
\newfontfamily\popxbold{PlusJakartaSans-ExtraBold.ttf}[Path=../../../fonts/]
\newfontfamily\popxboldls{PlusJakartaSans-ExtraBold.ttf}[Path=../../../fonts/,LetterSpace=3.0]

\setlist[enumerate]{leftmargin=1.7em,itemsep=5pt,topsep=4pt}
\setlist[itemize]{leftmargin=1.7em,itemsep=4pt,topsep=4pt,label={\color{poporange}\textbullet}}
\setlength{\parindent}{0pt}
\setlength{\parskip}{5pt}
\renewcommand{\arraystretch}{1.25}

% pgfplots : style Pop par défaut
\pgfplotsset{
  every axis/.append style={axis line style={popink,line width=1pt},tick style={popink},
    grid style={popline,line width=0.5pt},label style={font=\small},tick label style={font=\footnotesize}},
  popcurve/.style={poporange,line width=1.8pt,no markers,smooth},
}
% Même style disponible dans un \draw TikZ simple (hors pgfplots)
\tikzset{popcurve/.style={poporange,line width=1.8pt}}

% ============================================================
% Pill / squiggle
% ============================================================
\newcommand{\poppill}[3]{%
  \tikz[baseline=-0.55ex]\node[draw=popink,line width=1.2pt,rounded corners=2.6pt,%
    fill=#2,inner xsep=6.5pt,inner ysep=3pt]{%
    \popxboldls\fontsize{7}{7}\selectfont\color{#3}\MakeUppercase{#1}};%
}
\newcommand{\popsquiggle}[1]{%
  \tikz[baseline=-0.4ex]{
    \draw[#1,line width=2.6pt,line cap=round]
      plot [smooth] coordinates {(0,0.09) (0.34,0.01) (0.68,0.21) (1.02,0.01) (1.36,0.10)};
  }%
}
% Mot-clé surligné (marqueur orange sous le texte)
\newcommand{\cle}[1]{{\popxbold\color{popdark}#1}}

% ============================================================
% En-tête / pied
% ============================================================
\pagestyle{fancy}
\fancyhf{}
\renewcommand{\headrulewidth}{0pt}
\renewcommand{\footrulewidth}{0pt}
\lhead{\raisebox{-0.15ex}{\poplogo\fontsize{13}{13}\selectfont\color{popink}OnBuch\color{poporange}+}}
\rhead{\poppill{\DOCMATIERE\ \textbullet\ \DOCNIVEAU\ \textbullet\ Leçon \DOCLECON}{white}{popink}}
\lfoot{\popxbold\fontsize{7.6}{7.6}\selectfont\color{popmuted}\MakeUppercase{Cours signé OnBuch+}\ \textbullet\ {\popsemi\DOCTITRE}}
\rfoot{\tikz[baseline=-0.6ex]\node[rounded corners=8.5pt,fill=popink,minimum height=17pt,minimum width=17pt,inner xsep=5pt,inner sep=0]{\popxbold\fontsize{8}{8}\selectfont\color{popcream}\thepage\,/\,\pageref*{LastPage}};}

% ============================================================
% Titres
% ============================================================
\newcommand{\chaptitre}[2]{%
  \begin{tcolorbox}[enhanced,colback=poporange,colframe=popink,boxrule=2.6pt,arc=11pt,
      drop shadow=popink,left=15pt,right=15pt,top=12pt,bottom=11pt,before skip=3pt,after skip=18pt]
    \poppill{#2}{popink}{popcream}\par\vspace{9pt}
    {\raggedright\hyphenpenalty=10000\exhyphenpenalty=10000\popdisplay\fontsize{22}{24}\selectfont\color{popcream}\MakeUppercase{#1}\par}
  \end{tcolorbox}
}
% Section numérotée : pastille orange avec le numéro + titre Archivo
\newcounter{popsec}
\newcommand{\coursec}[1]{%
  \stepcounter{popsec}\setcounter{popsub}{0}%
  \par\vspace{16pt}\noindent
  \begin{minipage}{\linewidth}
  \tikz[baseline=-0.7ex]\node[draw=popink,line width=1.6pt,fill=poporange,rounded corners=4pt,minimum size=22pt,inner sep=0]{\popdisplay\fontsize{12}{12}\selectfont\color{popcream}\thepopsec};%
  \hspace{8pt}{\popdisplay\fontsize{15}{17}\selectfont\color{popink}\MakeUppercase{#1}}\par
  \vspace{4pt}\noindent{\color{poporange}\rule{36pt}{3.6pt}}
  \end{minipage}\par\nopagebreak\vspace{8pt}
}
\newcounter{popsub}
\newcommand{\courssub}[1]{%
  \stepcounter{popsub}%
  \par\vspace{9pt}\noindent{\popxbold\fontsize{11.5}{13}\selectfont\color{popink}{\color{poporange}\thepopsec.\thepopsub}\hspace{6pt}#1}\par\nopagebreak\vspace{4pt}
}

% ============================================================
% Boîtes pédagogiques — même recette : fond blanc, bordure épaisse colorée,
% ombre dure encre, badge plein. Titre optionnel [#1].
% ============================================================
\tcbset{popbase/.style={breakable,enhanced,colback=white,boxrule=2.3pt,arc=9pt,
  shadow={3pt}{-3pt}{0pt}{popink},left=14pt,right=14pt,top=11pt,bottom=11pt,before skip=11pt,after skip=15pt,
  fontupper=\popsemi\fontsize{10.6}{14.5}\selectfont\color{popink},
  fontlower=\popsemi\fontsize{10.6}{14.5}\selectfont\color{popink}}}
% Coche dessinée (le glyphe ✓ n'existe pas dans les polices du document)
\newcommand{\popcheck}[1]{\tikz[baseline=-0.5ex]\draw[#1,line width=1.6pt,line cap=round,line join=round](-0.13,0.0)--(-0.04,-0.09)--(0.13,0.1);}
\newcommand{\popboxhead}[3]{\poppill{#1}{#2}{white}\ifx\relax#3\relax\else\hspace{6pt}{\popxbold\fontsize{10.6}{12}\selectfont\color{popink}#3}\fi\par\vspace{7pt}}

\newtcolorbox{aretenir}[1][]{popbase,colframe=popgreen,before upper={\popboxhead{À retenir}{popgreen}{#1}}}
\newtcolorbox{exemplebox}[1][]{popbase,colframe=poporange,before upper={\popboxhead{Exemple}{poporange}{#1}}}
\newtcolorbox{definition}[1][]{popbase,colframe=popblue,before upper={\popboxhead{Définition}{popblue}{#1}}}
\newtcolorbox{propriete}[1][]{popbase,colframe=poppurple,before upper={\popboxhead{Propriété}{poppurple}{#1}}}
\newtcolorbox{methode}[1][]{popbase,colframe=popgold,before upper={\popboxhead{Méthode}{popgold}{#1}}}
\newtcolorbox{attention}[1][]{popbase,colframe=poppink,before upper={\popboxhead{Attention}{poppink}{#1}}}
\newtcolorbox{experience}[1][]{popbase,colframe=popblue,colback=popblueL,before upper={\popboxhead{Expérience}{popblue}{#1}}}
% « Le savais-tu ? » : carte violette pleine (contexte, culture, Cameroun)
\newtcolorbox{savaistu}[1][]{popbase,colback=poppurple,colframe=popink,
  fontupper=\popsemi\fontsize{10.4}{14.2}\selectfont\color{white},
  before upper={\poppill{Le savais-tu\,?}{white}{poppurple}\ifx\relax#1\relax\else\hspace{6pt}{\popxbold\color{white}#1}\fi\par\vspace{7pt}}}
% Exercice résolu : énoncé en haut, \tcblower puis la solution sur fond crème
\newcounter{popexo}\newcounter{popexr}
\newtcolorbox{exoresolu}[1][]{popbase,colframe=poporange,colbacklower=popcream,
  segmentation style={popink,line width=1.2pt,dashed},
  before upper={\stepcounter{popexr}\popboxhead{Exercice résolu \thepopexr}{poporange}{#1}},
  before lower={\poppill{Solution}{popink}{popcream}\par\vspace{6pt}}}
% Exercice (non résolu en place) — la correction est dans la partie Corrigés
\newtcolorbox{exercice}[1][]{popbase,colframe=popink,
  before upper={\stepcounter{popexo}\popboxhead{Exercice \thepopexo}{popink}{#1}}}
% Corrigé
\newtcolorbox{corrige}[1][]{popbase,colframe=popgreen,colback=popgreenL,
  before upper={\popboxhead{Corrigé}{popgreen}{#1}}}
% Objectifs / prérequis / bilan
\newtcolorbox{objectifs}{popbase,colframe=popink,colback=poporangeL,
  before upper={\popboxhead{Objectifs de la leçon}{popink}{}}}
\newtcolorbox{prerequis}{popbase,colframe=popink,colback=popblueL,
  before upper={\popboxhead{Prérequis}{popblue}{}}}
\newtcolorbox{bilan}{popbase,colframe=popink,colback=white,boxrule=2.8pt,
  before upper={\popboxhead{Fiche bilan}{poporange}{L'essentiel à connaître pour l'examen}}}
% Figure encadrée avec légende
\newtcolorbox{popfigure}[1][]{enhanced,breakable=false,colback=white,colframe=popline,boxrule=1.4pt,arc=8pt,
  left=10pt,right=10pt,top=10pt,bottom=8pt,before skip=10pt,after skip=12pt,halign=center,
  lower separated=false,halign lower=center,
  fontlower=\popsemi\fontsize{9}{11.5}\selectfont\color{popmuted},
  #1}
\newcommand{\legende}[1]{\par\vspace{6pt}{\popsemi\fontsize{9}{11.5}\selectfont\color{popmuted}#1}}

% ============================================================
% Signature OnBuch+
% ============================================================
\newcommand{\poplogoinline}[1]{{\poplogo\fontsize{#1}{#1}\selectfont\color{popink}OnBuch\color{poporange}+}}
\newcommand{\signatureonbuch}{%
  \begin{tcolorbox}[enhanced,colback=popink,colframe=popink,boxrule=2.2pt,arc=10pt,
      drop shadow=poporange,left=14pt,right=14pt,top=10pt,bottom=10pt,before skip=0pt,after skip=0pt]
    \tikz[baseline=-0.7ex]\node[circle,fill=popgreen,draw=popcream,line width=1.3pt,minimum size=22pt,inner sep=0]{\popcheck{white}};%
    \hspace{9pt}\begin{minipage}[c]{0.62\linewidth}
      {\popxbold\fontsize{12}{13}\selectfont\color{popcream}Signé\ \textbullet\ {\poplogo OnBuch\color{poporange}+}}\par\vspace{2pt}
      {\raggedright\popsemi\fontsize{8}{10.5}\selectfont\color{popline}\MakeUppercase{Équipe pédagogique OnBuch+}\\\MakeUppercase{Conforme au programme officiel MINESEC}\par}
    \end{minipage}\hfill
    {\poplogo\fontsize{22}{22}\selectfont\color{popcream}OnBuch\color{poporange}+}
  \end{tcolorbox}%
}

% ============================================================
% Page de garde
% #1 titre · #2 matière · #3 niveau · #4 module · #5 n° leçon · #6 durée ·
% #7 description / objectifs (texte ou itemize)
% ============================================================
\newcommand{\pagegarde}[7]{%
  \thispagestyle{empty}
  \newgeometry{a4paper,margin=1.7cm,top=1.6cm,bottom=1.4cm}%
  \begin{tikzpicture}[remember picture,overlay]
    \fill[poporange!18] ([xshift=-3.2cm,yshift=-3.6cm]current page.north east) circle (4.6cm);
    \fill[poppurple!14] ([xshift=2.4cm,yshift=7.6cm]current page.south west) circle (3.8cm);
  \end{tikzpicture}%
  {\poplogo\fontsize{30}{30}\selectfont\color{popink}OnBuch\color{poporange}+}\hfill
  \raisebox{0.6ex}{\poppill{Cours signé \textbullet\ Premium}{popink}{popcream}}\par
  \vspace{1pt}\popsquiggle{poppurple}\par
  \vspace{8pt}
  \begin{tcolorbox}[enhanced,colback=poporange,colframe=popink,boxrule=2.8pt,arc=13pt,
      drop shadow=popink,left=18pt,right=18pt,top=12pt,bottom=13pt,after skip=10pt]
    \poppill{#2 \textbullet\ #3}{popink}{popcream}\hspace{5pt}\poppill{#4}{white}{popink}\par\vspace{8pt}
    {\popdisplay\fontsize{14}{14}\selectfont\color{popink}LEÇON #5}\par\vspace{4pt}
    {\raggedright\hyphenpenalty=10000\exhyphenpenalty=10000\popdisplay\fontsize{25}{27}\selectfont\color{popcream}\MakeUppercase{#1}\par}
    \vspace{8pt}
    {\popxbold\fontsize{9.5}{9.5}\selectfont\color{popink}\MakeUppercase{Durée conseillée : #6}}
  \end{tcolorbox}
  \begin{tcolorbox}[enhanced,colback=white,colframe=popink,boxrule=2.2pt,arc=10pt,
      drop shadow=popink,left=16pt,right=16pt,top=10pt,bottom=10pt,after skip=8pt]
    \poppill{Dans cette leçon}{poppurple}{white}\par\vspace{5pt}
    \popsemi\fontsize{9.3}{12.6}\selectfont\color{popink}#7
  \end{tcolorbox}
  \noindent\begin{minipage}[t]{0.487\linewidth}
    \begin{tcolorbox}[enhanced,colback=popgreen,colframe=popink,boxrule=2.2pt,arc=10pt,
        drop shadow=popink,left=11pt,right=11pt,top=10pt,bottom=10pt,height=3.1cm,valign=center]
      \tcbox[colback=white,colframe=popink,boxrule=1.3pt,arc=5pt,boxsep=0pt,nobeforeafter,
        left=3pt,right=3pt,top=3pt,bottom=3pt,valign=center]{\qrcode[height=1.9cm]{https://play.google.com/store/apps/details?id=com.luvvix.onbuch&hl=fr}}%
      \hspace{8pt}\begin{minipage}[c]{0.5\linewidth}
        {\popdisplay\fontsize{11}{12.5}\selectfont\color{white}\MakeUppercase{Télécharge OnBuch}}\par\vspace{4pt}
        {\raggedright\popsemi\fontsize{8.4}{11}\selectfont\color{white}Gratuit \textbullet\ Google Play\\Tuteur IA, annales, concours, résultats\par}
      \end{minipage}
    \end{tcolorbox}
  \end{minipage}\hfill
  \begin{minipage}[t]{0.487\linewidth}
    \begin{tcolorbox}[enhanced,colback=poppurple,colframe=popink,boxrule=2.2pt,arc=10pt,
        drop shadow=popink,left=11pt,right=11pt,top=10pt,bottom=10pt,height=3.1cm,valign=center]
      \tikz[baseline=-0.7ex]\node[circle,fill=white,draw=popink,line width=1.3pt,minimum size=1.6cm,inner sep=0]{\popdisplay\fontsize{24}{24}\selectfont\color{poppurple}+};%
      \hspace{8pt}\begin{minipage}[c]{0.6\linewidth}
        {\popdisplay\fontsize{11}{12.5}\selectfont\color{white}\MakeUppercase{Passe à OnBuch+}}\par\vspace{4pt}
        {\raggedright\popsemi\fontsize{8.4}{11}\selectfont\color{white}Tous les cours signés, Tuteur IA illimité, fiches PDF — onglet OnBuch+ de l'app\par}
      \end{minipage}
    \end{tcolorbox}
  \end{minipage}
  \vspace{8pt}
  \popcontenu
  \vfill
  \signatureonbuch
  \restoregeometry
  \newpage
}

% Rangée « Ce cours contient » (page de garde)
\newcommand{\poptile}[3]{% #1 couleur #2 gros texte #3 légende
  \begin{tcolorbox}[enhanced,colback=#1,colframe=popink,boxrule=2pt,arc=9pt,shadow={2.5pt}{-2.5pt}{0pt}{popink},
      left=6pt,right=6pt,top=9pt,bottom=9pt,halign=center,valign=center,height=1.75cm,width=\linewidth,nobeforeafter]
    {\popdisplay\fontsize{12.5}{13}\selectfont\color{white}\MakeUppercase{#2}}\par\vspace{3pt}
    {\centering\popsemi\fontsize{7.8}{9.5}\selectfont\color{white}#3\par}
  \end{tcolorbox}}
\newcommand{\popcontenu}{%
  \par\noindent{\popxboldls\fontsize{7.5}{7.5}\selectfont\color{popmuted}CE COURS SIGNÉ CONTIENT}\par\vspace{7pt}
  \noindent\begin{minipage}[t]{0.235\linewidth}\poptile{popblue}{Cours}{complet, illustré, pas à pas}\end{minipage}\hfill
  \begin{minipage}[t]{0.235\linewidth}\poptile{poporange}{Méthodes}{et exercices résolus}\end{minipage}\hfill
  \begin{minipage}[t]{0.235\linewidth}\poptile{poppink}{Exercices}{type examen + corrigés}\end{minipage}\hfill
  \begin{minipage}[t]{0.235\linewidth}\poptile{popgreen}{Fiche bilan}{l'essentiel à retenir}\end{minipage}\par}

% Sommaire
\newtcolorbox{sommaire}{enhanced,colback=white,colframe=popink,boxrule=2.2pt,arc=10pt,
  drop shadow=popink,left=18pt,right=18pt,top=13pt,bottom=13pt,before skip=6pt,after skip=20pt,
  before upper={\poppill{Sommaire}{poppurple}{white}\par\vspace{9pt}\popsemi\fontsize{10.6}{14.5}\selectfont\color{popink}}}

% Signature de fin de cours — poussée tout en bas de la dernière page
\newcommand{\findecours}{%
  \par\vfill\nopagebreak
  \begin{center}\popsquiggle{poporange}\end{center}\vspace{4pt}
  \signatureonbuch
}
\AtBeginDocument{\def\llbracket{\mathopen{[\mkern-3.5mu[}}\def\rrbracket{\mathclose{]\mkern-3.5mu]}}} % intervalles d'entiers (glyphe ⟦ absent de la police)
% Glyphes absents de Fira Math : redéfinis à partir de glyphes disponibles
\AtBeginDocument{%
  \def\setminus{\mathbin{\backslash}}\let\smallsetminus\setminus
  \def\vdots{\mathord{\vbox{\baselineskip3.2pt\lineskiplimit0pt\kern4pt\hbox{.}\hbox{.}\hbox{.}}}}%
  \def\ddots{\mathinner{\mkern1mu\raise6.5pt\hbox{.}\mkern2mu\raise3.6pt\hbox{.}\mkern2mu\raise0.7pt\hbox{.}\mkern1mu}}%
  \def\triangle{\mathord{\scalebox{0.9}{$\symup{\Delta}$}}}%
  \def\nabla{\mathord{\raisebox{\depth}{\scalebox{1}[-1]{$\symup{\Delta}$}}}}%
  \def\checkmark{\popcheck{popdark}}%
  \def\star{\mathbin{\ast}}\def\bigstar{\mathord{\ast}}%
  \def\diamond{\mathbin{\scalebox{0.6}{$\Diamond$}}}%
  \def\bigcup{\mathop{\mathchoice{\raisebox{-0.3em}{\scalebox{1.7}{$\cup$}}}{\scalebox{1.25}{$\cup$}}{\cup}{\cup}}\displaylimits}%
  \def\bigcap{\mathop{\mathchoice{\raisebox{-0.3em}{\scalebox{1.7}{$\cap$}}}{\scalebox{1.25}{$\cap$}}{\cap}{\cap}}\displaylimits}%
  \def\lesseqqgtr{\mathrel{\lessgtr}}\def\bigoplus{\mathop{\oplus}\displaylimits}%
}
\providecommand{\ssi}{si et seulement si} % abréviation fréquente
\AtBeginDocument{\providecommand{\wideparen}{\overparen}} % arc de cercle

%% FIGFIX-BEGIN v4 — correctifs globaux des figures (docs/onbuch-generation/tools/figfix)
\usepackage{adjustbox}\usepackage{etoolbox}
% toute figure TikZ/pgfplots plus large que la ligne est réduite à la largeur disponible
\BeforeBeginEnvironment{tikzpicture}{\begin{adjustbox}{max width=\linewidth}}
\AfterEndEnvironment{tikzpicture}{\end{adjustbox}}
% légendes pgfplots lisibles par-dessus les courbes
\pgfplotsset{every axis/.append style={legend style={fill=white,fill opacity=0.92,text opacity=1,draw=black!15,inner sep=3pt}}}
% étiquettes d'arêtes (syntaxe edge["..."]) avec fond blanc
\tikzset{every edge quotes/.append style={fill=white,inner sep=1.2pt}}
%% FIGFIX-END
\begin{document}

\pagegarde{Trigonométrie : formules, équations et inéquations}{Mathématiques}{1ère E}{Relations et opérations fondamentales dans l'ensemble des nombres réels}{N03}{12 h}{Cette leçon complète l'étude de la trigonométrie en Première : après les formules d'addition et de duplication, l'élève apprend à transformer les expressions trigonométriques, puis à résoudre dans ℝ ou sur un intervalle les équations et inéquations faisant intervenir cos X, sin X, tan X et a cos X + b sin X, en s'appuyant sur le cercle trigonométrique et les courbes représentatives.
\vspace{4pt}
\begin{multicols}{2}\small
\begin{itemize}
\item À la fin de cette leçon, je sais exprimer cos(a+b), cos(a−b), sin(a+b), sin(a−b) à l'aide de cos a, cos b, sin a et sin b
\item À la fin de cette leçon, je sais exprimer tan(a+b) et tan(a−b) en fonction de tan a et tan b
\item À la fin de cette leçon, je sais calculer des valeurs exactes de cosinus, sinus et tangente de réels bien choisis (π/12, 5π/12, ...)
\item À la fin de cette leçon, je sais écrire a cos X + b sin X sous la forme A cos(X + φ)
\item À la fin de cette leçon, je sais justifier graphiquement l'existence et le nombre de solutions des équations cos X = a, sin X = a, tan X = a sur un intervalle borné
\item À la fin de cette leçon, je sais résoudre dans un intervalle quelconque de ℝ les équations cos X = a, sin X = a, tan X = a et a cos X + b sin X = c
\end{itemize}\end{multicols}}

\begin{sommaire}
\begin{multicols}{2}
\begin{enumerate}
\item Formules d'addition et de duplication
\item Transformation de a cos X + b sin X
\item Équations cos X = a, sin X = a, tan X = a
\item Équations du type a cos X + b sin X = c
\item Inéquations trigonométriques de base
\item Inéquations avec a cos X + b sin X et synthèse
\item Activité d'intégration
\item Méthodes et exercices résolus
\item Exercices
\item Corrigés des exercices
\item Fiche bilan
\end{enumerate}
\end{multicols}
\end{sommaire}

\begin{objectifs}
\begin{itemize}
\item À la fin de cette leçon, je sais exprimer cos(a+b), cos(a−b), sin(a+b), sin(a−b) à l'aide de cos a, cos b, sin a et sin b
\item À la fin de cette leçon, je sais exprimer tan(a+b) et tan(a−b) en fonction de tan a et tan b
\item À la fin de cette leçon, je sais calculer des valeurs exactes de cosinus, sinus et tangente de réels bien choisis (π/12, 5π/12, ...)
\item À la fin de cette leçon, je sais écrire a cos X + b sin X sous la forme A cos(X + φ)
\item À la fin de cette leçon, je sais justifier graphiquement l'existence et le nombre de solutions des équations cos X = a, sin X = a, tan X = a sur un intervalle borné
\item À la fin de cette leçon, je sais résoudre dans un intervalle quelconque de ℝ les équations cos X = a, sin X = a, tan X = a et a cos X + b sin X = c
\item À la fin de cette leçon, je sais délimiter sur le cercle trigonométrique les points images vérifiant une inégalité sur cos X, sin X ou tan X
\item À la fin de cette leçon, je sais résoudre une inéquation trigonométrique sur un intervalle borné puis sur un intervalle quelconque de ℝ
\item À la fin de cette leçon, je sais résoudre algébriquement et graphiquement une inéquation où a cos X + b sin X est comparé à un réel c
\end{itemize}
\end{objectifs}

\begin{prerequis}
\begin{itemize}
\item Cercle trigonométrique, enroulement de la droite réelle, mesures d'angles en radians (leçon N01)
\item Cosinus, sinus et tangente d'un réel ; valeurs remarquables (π/6, π/4, π/3)
\item Angles associés : cos(−x), sin(−x), cos(π−x), sin(π+x), cos(π/2 − x), etc.
\item Résolution d'équations et d'inéquations du premier et du second degré dans ℝ
\item Lecture graphique : courbes des fonctions sinus, cosinus et tangente
\end{itemize}
\end{prerequis}

\newpage

\coursec{Formules d'addition et de duplication}

À Douala, un technicien règle l'inclinaison d'un panneau solaire : l'énergie captée dépend du \cle{cosinus} de l'angle d'incidence des rayons du soleil. Pour garantir un rendement minimal, il doit déterminer pour quelles heures de la journée cet angle vérifie une inéquation trigonométrique. De même, un ingénieur télécoms à Yaoundé superpose deux signaux sinusoïdaux et doit écrire leur somme sous la forme d'une seule onde pour trouver les instants où le signal s'annule. Dans les deux cas, il faut savoir \emph{transformer} des expressions comme $\cos(a-b)$ ou $\cos x + \sin x$ en expressions plus simples.

\textbf{Problématique.} Comment exprimer le cosinus, le sinus et la tangente d'une somme ou d'une différence d'angles à l'aide des cosinus et sinus de chaque angle ? Comment en déduire des valeurs exactes nouvelles ?

\courssub{Rappels indispensables}

Avant d'additionner des angles, rappelons les deux outils de base de toute la trigonométrie.

\begin{aretenir}[Identité fondamentale]
Pour tout réel $x$ :
\[ \cos^2 x + \sin^2 x = 1. \]
C'est une conséquence directe du théorème de Pythagore appliqué au cercle trigonométrique (de rayon $1$) : le point image $M(\cos x\,;\,\sin x)$ vérifie $OM^2 = \cos^2 x + \sin^2 x = 1$.
\end{aretenir}

\begin{aretenir}[Valeurs remarquables]
\begin{center}
\begin{tabular}{|c|c|c|c|c|c|}
\hline
\rowcolor{popblueL}
$x$ & $0$ & $\dfrac{\pi}{6}$ & $\dfrac{\pi}{4}$ & $\dfrac{\pi}{3}$ & $\dfrac{\pi}{2}$ \\
\hline
$\cos x$ & $1$ & $\dfrac{\sqrt{3}}{2}$ & $\dfrac{\sqrt{2}}{2}$ & $\dfrac{1}{2}$ & $0$ \\
\hline
$\sin x$ & $0$ & $\dfrac{1}{2}$ & $\dfrac{\sqrt{2}}{2}$ & $\dfrac{\sqrt{3}}{2}$ & $1$ \\
\hline
$\tan x$ & $0$ & $\dfrac{\sqrt{3}}{3}$ & $1$ & $\sqrt{3}$ & non défini \\
\hline
\end{tabular}
\end{center}
On retient aussi : $\cos(-x) = \cos x$ (parité), $\sin(-x) = -\sin x$ (imparité), et les formules de l'angle complémentaire :
\[ \cos\left(\tfrac{\pi}{2} - x\right) = \sin x \qquad \text{et} \qquad \sin\left(\tfrac{\pi}{2} - x\right) = \cos x. \]
\end{aretenir}

\courssub{La formule fondamentale : $\cos(a-b)$}

\textbf{Intuition.} Sur le cercle trigonométrique, l'angle $\widehat{AOB}$ entre les points images de $a$ et de $b$ ne dépend que de la \emph{différence} $a - b$. Or cet angle peut se mesurer de deux façons : géométriquement (c'est $a-b$) et algébriquement (avec les coordonnées de $A$ et $B$). En comparant les deux, on obtient une formule.

\begin{propriete}[Formule fondamentale — démonstration exigible]
Pour tous réels $a$ et $b$ :
\[ \cos(a - b) = \cos a \cos b + \sin a \sin b. \]
\end{propriete}

\begin{experience}[Démonstration par le produit scalaire]
Plaçons sur le cercle trigonométrique de centre $O$ les points $A$ et $B$, images respectives des réels $a$ et $b$ :
\[ A(\cos a\,;\,\sin a) \qquad \text{et} \qquad B(\cos b\,;\,\sin b). \]
Posons $\vec{u} = \overrightarrow{OA}$ et $\vec{v} = \overrightarrow{OB}$. Comme le cercle a pour rayon $1$, on a $\|\vec{u}\| = \|\vec{v}\| = 1$ : ce sont des vecteurs \cle{unitaires}.

\textbf{Étape 1 : calcul géométrique du produit scalaire.} L'angle $(\vec{u}\,;\,\vec{v})$ a pour mesure $b - a$ (ou $a - b$, le cosinus étant pair, cela ne change rien). Donc :
\[ \vec{u} \cdot \vec{v} = \|\vec{u}\| \times \|\vec{v}\| \times \cos(a - b) = 1 \times 1 \times \cos(a-b) = \cos(a-b). \]

\textbf{Étape 2 : calcul analytique du produit scalaire.} Avec les coordonnées dans un repère orthonormé :
\[ \vec{u} \cdot \vec{v} = x_A x_B + y_A y_B = \cos a \cos b + \sin a \sin b. \]

\textbf{Étape 3 : conclusion.} Les deux expressions du même produit scalaire sont égales :
\[ \cos(a-b) = \cos a \cos b + \sin a \sin b. \qquad \blacksquare \]
\end{experience}

\begin{popfigure}
\begin{center}
\begin{tikzpicture}[scale=2.6]
\draw[popline,->] (-1.4,0) -- (1.5,0) node[below left] {$\cos$};
\draw[popline,->] (0,-1.3) -- (0,1.4) node[below left] {$\sin$};
\draw[popblue,thick] (0,0) circle (1);
\coordinate (O) at (0,0);
\coordinate (A) at (0.5,0.866);
\coordinate (B) at (0.94,0.342);
\draw[popink,thick,->] (O) -- (A) node[above right] {$A(\cos a\,;\,\sin a)$};
\draw[poporange,thick,->] (O) -- (B) node[right] {$B(\cos b\,;\,\sin b)$};
\fill[popink] (A) circle (0.02);
\fill[poporange] (B) circle (0.02);
\draw[popink,->] (0.45,0) arc (0:60:0.45) node[midway,right] {$a$};
\draw[poporange,->] (0.62,0) arc (0:20:0.62) node[midway,right] {$b$};
\draw[poppurple,->] ([shift=(20:0.85)]O) arc (20:60:0.85) node[midway,above right] {$a-b$};
\draw[popmuted,dashed] (A) -- (B);
\node[popdark] at (0,-1.2) {$\vec{u}\cdot\vec{v} = \cos(a-b) = \cos a\cos b + \sin a\sin b$};
\end{tikzpicture}
\end{center}
\legende{Cercle trigonométrique : les vecteurs $\overrightarrow{OA}$ et $\overrightarrow{OB}$ sont unitaires, leur produit scalaire se calcule de deux façons.}
\end{popfigure}

\courssub{Formules d'addition pour le cosinus et le sinus}

Toutes les autres formules se \emph{déduisent} de la formule fondamentale. C'est un point essentiel de la leçon : il ne faut apprendre par cœur qu'une seule formule, et savoir retrouver les autres.

\begin{experience}[Déduction guidée de $\cos(a+b)$]
L'idée est simple : une somme est une différence avec l'opposé. Écrivons $a + b = a - (-b)$ et appliquons la formule fondamentale avec $-b$ à la place de $b$ :
\begin{align*}
\cos(a + b) &= \cos\bigl(a - (-b)\bigr) \\
&= \cos a \cos(-b) + \sin a \sin(-b) \\
&= \cos a \cos b - \sin a \sin b,
\end{align*}
car $\cos(-b) = \cos b$ (parité) et $\sin(-b) = -\sin b$ (imparité). $\blacksquare$
\end{experience}

\begin{experience}[Déduction guidée de $\sin(a+b)$ et $\sin(a-b)$]
On utilise le passage du sinus au cosinus par l'angle complémentaire : $\sin x = \cos\left(\dfrac{\pi}{2} - x\right)$. Avec $x = a+b$ :
\begin{align*}
\sin(a+b) &= \cos\left(\tfrac{\pi}{2} - a - b\right) = \cos\left(\left(\tfrac{\pi}{2} - a\right) - b\right) \\
&= \cos\left(\tfrac{\pi}{2} - a\right)\cos b + \sin\left(\tfrac{\pi}{2} - a\right)\sin b \\
&= \sin a \cos b + \cos a \sin b.
\end{align*}
En remplaçant alors $b$ par $-b$, on obtient :
\[ \sin(a-b) = \sin a \cos b - \cos a \sin b. \qquad \blacksquare \]
\end{experience}

\begin{propriete}[Formules d'addition]
Pour tous réels $a$ et $b$ :
\begin{align*}
\cos(a+b) &= \cos a \cos b - \sin a \sin b, &
\cos(a-b) &= \cos a \cos b + \sin a \sin b, \\
\sin(a+b) &= \sin a \cos b + \cos a \sin b, &
\sin(a-b) &= \sin a \cos b - \cos a \sin b.
\end{align*}
\end{propriete}

\begin{attention}[Le piège classique : $\cos(a+b) \neq \cos a + \cos b$]
Le cosinus \textbf{n'est pas additif} ! Contre-exemple avec $a = b = \dfrac{\pi}{3}$ :
\[ \cos\left(\tfrac{\pi}{3} + \tfrac{\pi}{3}\right) = \cos\tfrac{2\pi}{3} = -\tfrac{1}{2}, \qquad \text{alors que} \qquad \cos\tfrac{\pi}{3} + \cos\tfrac{\pi}{3} = \tfrac{1}{2} + \tfrac{1}{2} = 1. \]
Les deux résultats sont différents. De même $\sin(a+b) \neq \sin a + \sin b$. On doit \emph{toujours} développer avec les formules d'addition.
\end{attention}

\courssub{Formules d'addition pour la tangente}

\begin{propriete}[Tangente d'une somme et d'une différence]
Pour tous réels $a$ et $b$ tels que $\tan a$, $\tan b$ et le dénominateur existent :
\[ \tan(a+b) = \frac{\tan a + \tan b}{1 - \tan a \tan b}, \qquad \tan(a-b) = \frac{\tan a - \tan b}{1 + \tan a \tan b}. \]
\textbf{Conditions d'existence :} il faut $a \not\equiv \dfrac{\pi}{2}\ [\pi]$, $b \not\equiv \dfrac{\pi}{2}\ [\pi]$, et $a + b \not\equiv \dfrac{\pi}{2}\ [\pi]$ (respectivement $a-b$), ce qui garantit aussi $\tan a \tan b \neq 1$ (respectivement $\neq -1$).
\end{propriete}

\begin{experience}[Démonstration guidée (exercice dirigé)]
Supposons toutes les quantités définies. Partons de la définition $\tan(a+b) = \dfrac{\sin(a+b)}{\cos(a+b)}$ :
\[ \tan(a+b) = \frac{\sin a \cos b + \cos a \sin b}{\cos a \cos b - \sin a \sin b}. \]
L'astuce consiste à diviser le numérateur \emph{et} le dénominateur par $\cos a \cos b$ (non nul car $\tan a$ et $\tan b$ existent) :
\[ \tan(a+b) = \frac{\dfrac{\sin a}{\cos a} + \dfrac{\sin b}{\cos b}}{1 - \dfrac{\sin a \sin b}{\cos a \cos b}} = \frac{\tan a + \tan b}{1 - \tan a \tan b}. \]
La formule de $\tan(a-b)$ s'obtient en remplaçant $b$ par $-b$ et en utilisant $\tan(-b) = -\tan b$. $\blacksquare$
\end{experience}

\courssub{Formules de duplication et de linéarisation}

Que se passe-t-il lorsque $b = a$ dans les formules d'addition ? On obtient les formules de \cle{duplication}, qui relient les lignes trigonométriques de $2a$ à celles de $a$.

\begin{propriete}[Formules de duplication]
Pour tout réel $a$ :
\[ \cos 2a = \cos^2 a - \sin^2 a = 2\cos^2 a - 1 = 1 - 2\sin^2 a, \qquad \sin 2a = 2\sin a \cos a, \]
et, lorsque $\tan a$ et $\tan 2a$ existent :
\[ \tan 2a = \frac{2\tan a}{1 - \tan^2 a}. \]
\end{propriete}

\begin{experience}[Démonstration]
Dans $\cos(a+b) = \cos a\cos b - \sin a \sin b$, posons $b = a$ : $\cos 2a = \cos^2 a - \sin^2 a$. En remplaçant $\sin^2 a$ par $1 - \cos^2 a$, on obtient $\cos 2a = 2\cos^2 a - 1$ ; en remplaçant $\cos^2 a$ par $1 - \sin^2 a$, on obtient $\cos 2a = 1 - 2\sin^2 a$. De même, $\sin 2a = \sin a \cos a + \cos a \sin a = 2\sin a \cos a$, et $\tan 2a = \dfrac{\tan a + \tan a}{1 - \tan a \tan a} = \dfrac{2\tan a}{1 - \tan^2 a}$. $\blacksquare$
\end{experience}

En \emph{isolant} $\cos^2 a$ et $\sin^2 a$ dans les formules de duplication, on obtient les formules de \cle{linéarisation}, très utiles pour faire descendre le degré d'une expression (notamment en physique et en calcul d'intégrales en Terminale).

\begin{propriete}[Formules de linéarisation]
Pour tout réel $a$ :
\[ \cos^2 a = \frac{1 + \cos 2a}{2}, \qquad \sin^2 a = \frac{1 - \cos 2a}{2}. \]
\end{propriete}

\begin{aretenir}[Tableau récapitulatif des formules]
\begin{center}
\begin{tabularx}{\linewidth}{|l|X|}
\hline
\rowcolor{popblueL}
\textbf{Type} & \textbf{Formules} \\
\hline
Addition (cos) & $\cos(a\pm b) = \cos a\cos b \mp \sin a \sin b$ \\
\hline
Addition (sin) & $\sin(a\pm b) = \sin a\cos b \pm \cos a\sin b$ \\
\hline
Addition (tan) & $\tan(a\pm b) = \dfrac{\tan a \pm \tan b}{1 \mp \tan a \tan b}$ \\
\hline
Duplication & $\cos 2a = 2\cos^2 a - 1 = 1 - 2\sin^2 a$ ; \quad $\sin 2a = 2\sin a\cos a$ ; \quad $\tan 2a = \dfrac{2\tan a}{1-\tan^2 a}$ \\
\hline
Linéarisation & $\cos^2 a = \dfrac{1+\cos 2a}{2}$ ; \quad $\sin^2 a = \dfrac{1-\cos 2a}{2}$ \\
\hline
\end{tabularx}
\end{center}
Attention aux signes croisés : dans $\cos(a+b)$ c'est un \textbf{moins}, dans $\sin(a+b)$ c'est un \textbf{plus}.
\end{aretenir}

\courssub{Application : calcul de valeurs exactes nouvelles}

Toute la puissance des formules d'addition apparaît ici : en décomposant un angle « exotique » comme somme ou différence d'angles remarquables, on calcule des valeurs \emph{exactes} inaccessibles directement.

\begin{methode}[Calculer une valeur exacte par décomposition]
\begin{enumerate}
\item Chercher deux angles remarquables ($\dfrac{\pi}{6}, \dfrac{\pi}{4}, \dfrac{\pi}{3}$\ldots) dont la somme ou la différence donne l'angle visé. Exemple : $\dfrac{\pi}{12} = \dfrac{\pi}{3} - \dfrac{\pi}{4}$ car $\dfrac{4\pi}{12} - \dfrac{3\pi}{12} = \dfrac{\pi}{12}$.
\item Écrire la formule d'addition correspondante.
\item Substituer les valeurs remarquables connues.
\item Simplifier soigneusement l'expression (radicaux, dénominateur commun).
\end{enumerate}
\end{methode}

\begin{exoresolu}[Calcul exact de $\cos\dfrac{\pi}{12}$, $\sin\dfrac{\pi}{12}$ et $\tan\dfrac{\pi}{12}$]
Énoncé. En remarquant que $\dfrac{\pi}{12} = \dfrac{\pi}{3} - \dfrac{\pi}{4}$, calculer les valeurs exactes de $\cos\dfrac{\pi}{12}$, $\sin\dfrac{\pi}{12}$ et $\tan\dfrac{\pi}{12}$.
\tcblower
\textbf{Solution détaillée.} On applique les formules de la différence avec $a = \dfrac{\pi}{3}$ et $b = \dfrac{\pi}{4}$.

\textbf{1. Cosinus :}
\begin{align*}
\cos\frac{\pi}{12} &= \cos\left(\frac{\pi}{3} - \frac{\pi}{4}\right) = \cos\frac{\pi}{3}\cos\frac{\pi}{4} + \sin\frac{\pi}{3}\sin\frac{\pi}{4} \\
&= \frac{1}{2}\times\frac{\sqrt{2}}{2} + \frac{\sqrt{3}}{2}\times\frac{\sqrt{2}}{2} = \frac{\sqrt{2}}{4} + \frac{\sqrt{6}}{4} = \boxed{\frac{\sqrt{6}+\sqrt{2}}{4}}.
\end{align*}

\textbf{2. Sinus :}
\begin{align*}
\sin\frac{\pi}{12} &= \sin\left(\frac{\pi}{3} - \frac{\pi}{4}\right) = \sin\frac{\pi}{3}\cos\frac{\pi}{4} - \cos\frac{\pi}{3}\sin\frac{\pi}{4} \\
&= \frac{\sqrt{3}}{2}\times\frac{\sqrt{2}}{2} - \frac{1}{2}\times\frac{\sqrt{2}}{2} = \boxed{\frac{\sqrt{6}-\sqrt{2}}{4}}.
\end{align*}

\textbf{3. Tangente :} on peut faire le quotient des deux résultats précédents :
\[ \tan\frac{\pi}{12} = \frac{\sqrt{6}-\sqrt{2}}{\sqrt{6}+\sqrt{2}} = \frac{(\sqrt{6}-\sqrt{2})^2}{(\sqrt{6})^2-(\sqrt{2})^2} = \frac{8 - 2\sqrt{12}}{4} = 2 - \sqrt{3}. \]

\textbf{Vérification numérique :} $\dfrac{\sqrt{6}+\sqrt{2}}{4} \approx \dfrac{2{,}449 + 1{,}414}{4} \approx 0{,}966$, et $\cos 15° \approx 0{,}966$ à la calculatrice. Cohérent !
\end{exoresolu}

\begin{exemplebox}[Deux autres valeurs exactes]
\textbf{a)} Comme $\dfrac{5\pi}{12} = \dfrac{\pi}{2} - \dfrac{\pi}{12}$, les formules de l'angle complémentaire donnent immédiatement :
\[ \cos\frac{5\pi}{12} = \sin\frac{\pi}{12} = \frac{\sqrt{6}-\sqrt{2}}{4} \qquad \text{et} \qquad \sin\frac{5\pi}{12} = \cos\frac{\pi}{12} = \frac{\sqrt{6}+\sqrt{2}}{4}. \]
\textbf{b)} Avec la duplication : $\cos\dfrac{\pi}{12} = 2\cos^2\dfrac{\pi}{24} - 1$, d'où l'on peut tirer $\cos\dfrac{\pi}{24}$, et ainsi de suite : les formules « engendrent » toute une famille de valeurs exactes.
\end{exemplebox}

\begin{attention}[Rédaction et conditions d'existence]
\begin{itemize}
\item Dans les formules de la tangente, \textbf{toujours vérifier} que les tangentes existent et que le dénominateur est non nul avant d'écrire la formule. Écrire $\tan(a+b) = \dfrac{\tan a + \tan b}{1 - \tan a \tan b}$ pour $a = b = \dfrac{\pi}{4}$ donne un dénominateur nul : normal, car $\tan\dfrac{\pi}{2}$ n'existe pas !
\item Ne pas oublier les parenthèses : $\cos(a-b)$ n'est pas $\cos a - b$.
\item Dans une vérification à la calculatrice, penser à se mettre en mode \textbf{radian}.
\end{itemize}
\end{attention}

\begin{savaistu}[Des formules qui font de la musique]
Les formules d'addition expliquent le phénomène de \emph{battements} en acoustique. Quand deux sons de fréquences proches se superposent — deux cordes de guitare presque accordées, ou deux signaux radio chez un opérateur télécoms —, la somme des deux ondes se transforme (grâce à $\cos a + \cos b = 2\cos\dfrac{a+b}{2}\cos\dfrac{a-b}{2}$, conséquence directe de nos formules) en une onde dont l'amplitude « pulse » lentement : c'est ce wah-wah-wah que perçoit l'oreille. Les accordeurs d'instruments et les ingénieurs du son à Douala ou Yaoundé exploitent ce phénomène chaque jour !
\end{savaistu}

\begin{exercice}[Pour t'entraîner]
\begin{enumerate}
\item En écrivant $\dfrac{7\pi}{12} = \dfrac{\pi}{3} + \dfrac{\pi}{4}$, calculer les valeurs exactes de $\cos\dfrac{7\pi}{12}$ et $\sin\dfrac{7\pi}{12}$.
\item Sachant que $\cos a = \dfrac{3}{5}$ avec $a \in \left[0\,;\,\dfrac{\pi}{2}\right]$, calculer $\sin a$, puis $\cos 2a$ et $\sin 2a$.
\item Montrer que pour tout réel $x$ : $\cos\left(x + \dfrac{\pi}{4}\right) = \dfrac{\sqrt{2}}{2}\bigl(\cos x - \sin x\bigr)$.
\item En utilisant la linéarisation, exprimer $\cos^2\dfrac{\pi}{8}$ en fonction de $\cos\dfrac{\pi}{4}$, puis retrouver la valeur exacte de $\cos\dfrac{\pi}{8}$.
\end{enumerate}
\end{exercice}

\coursec{Transformation de $a\cos X + b\sin X$}

Dans la section précédente, nous avons établi les formules d'addition et de duplication. Nous allons maintenant en dégager une conséquence spectaculaire, au cœur du programme de Première : \emph{toute} expression de la forme $a\cos X + b\sin X$ peut se ramener à une \emph{seule} fonction cosinus, décalée et amplifiée. C'est cette transformation qui nous permettra, dans les sections suivantes, de résoudre les équations et inéquations du type $a\cos X + b\sin X = c$.

\courssub{Le constat : une somme de sinusoïdes est une sinusoïde}

\begin{experience}[Une observation qui change tout]
Considérons la fonction $f(X) = \cos X + \sqrt{3}\sin X$. À première vue, c'est une somme de deux fonctions : difficile d'en prévoir l'allure. Calculons quelques valeurs :
\begin{align*}
f(0) &= 1 + 0 = 1, &
f\left(\tfrac{\pi}{3}\right) &= \tfrac{1}{2} + \sqrt{3}\times\tfrac{\sqrt{3}}{2} = \tfrac{1}{2}+\tfrac{3}{2} = 2, &
f(\pi) &= -1 + 0 = -1.
\end{align*}
La valeur $2$ atteinte en $\frac{\pi}{3}$ est remarquable : c'est exactement $\sqrt{1^2 + (\sqrt{3})^2} = \sqrt{4}$. Or, d'après la formule d'addition $\cos(X-\varphi) = \cos X\cos\varphi + \sin X\sin\varphi$,
\[ 2\cos\left(X - \tfrac{\pi}{3}\right) = 2\left(\cos X \cos\tfrac{\pi}{3} + \sin X \sin\tfrac{\pi}{3}\right) = 2\left(\tfrac{1}{2}\cos X + \tfrac{\sqrt{3}}{2}\sin X\right) = \cos X + \sqrt{3}\sin X = f(X). \]
La somme $\cos X + \sqrt{3}\sin X$ n'est donc rien d'autre qu'\emph{une seule} fonction cosinus, d'amplitude $2$, déphasée de $\frac{\pi}{3}$. La figure ci-dessous confirme visuellement cette identité : les deux courbes sont confondues.
\end{experience}

\begin{popfigure}
\begin{center}
\begin{tikzpicture}
\begin{axis}[
  width=0.95\linewidth, height=6.2cm,
  axis lines=middle,
  xlabel={$X$}, ylabel={$y$},
  xmin=-3.6, xmax=6.6, ymin=-2.6, ymax=2.6,
  xtick={-3.1416,-1.5708,0,1.5708,3.1416,4.7124,6.2832},
  xticklabels={$-\pi$,$-\frac{\pi}{2}$,$0$,$\frac{\pi}{2}$,$\pi$,$\frac{3\pi}{2}$,$2\pi$},
  ytick={-2,-1,1,2},
  legend style={at={(0.02,0.02)},anchor=south west,font=\small},
  grid=both, grid style={popline!40},
]
\addplot[popcurve,domain=-3.4:6.4,samples=250]{cos(deg(x)) + sqrt(max(0,3))*sin(deg(x))};
\addlegendentry{$f(X)=\cos X+\sqrt{3}\sin X$}
\addplot[poporange,dashed,very thick,domain=-3.4:6.4,samples=250]{2*cos(deg(x - 1.0472))};
\addlegendentry{$2\cos\left(X-\frac{\pi}{3}\right)$}
\end{axis}
\end{tikzpicture}
\end{center}
\legende{Les deux courbes sont parfaitement superposées : la somme $\cos X + \sqrt{3}\sin X$ est une unique sinusoïde d'amplitude $2$.}
\end{popfigure}

Ce phénomène est général : additionner un cosinus et un sinus de \emph{même} variable, c'est toujours obtenir une sinusoïde unique. Le théorème suivant le dit précisément.

\courssub{Le théorème fondamental}

\begin{propriete}[Théorème — forme $A\cos(X+\varphi)$]
Soient $a$ et $b$ deux réels tels que $(a,b) \neq (0,0)$. Il existe un réel $A > 0$ et un réel $\varphi$ tels que, pour tout réel $X$ :
\[ a\cos X + b\sin X = A\cos(X + \varphi), \qquad \text{avec } A = \sqrt{a^2+b^2}, \]
où $\varphi$ est un réel vérifiant
\[ \cos\varphi = \frac{a}{A} \qquad \text{et} \qquad \sin\varphi = -\frac{b}{A}. \]
Le réel $A$ s'appelle l'\cle{amplitude} et $\varphi$ la \cle{phase} (ou déphasage). La démonstration ci-dessous est formatrice et exigible en classe.
\end{propriete}

\begin{experience}[Démonstration guidée pas à pas]
Posons $A = \sqrt{a^2+b^2}$. Comme $(a,b)\neq(0,0)$, on a $A > 0$, donc on peut factoriser :
\[ a\cos X + b\sin X = A\left(\frac{a}{A}\cos X + \frac{b}{A}\sin X\right). \]
\textbf{Étape 1 : un point du cercle trigonométrique.} Calculons :
\[ \left(\frac{a}{A}\right)^2 + \left(-\frac{b}{A}\right)^2 = \frac{a^2+b^2}{A^2} = \frac{A^2}{A^2} = 1. \]
Le point de coordonnées $\left(\frac{a}{A}\,;\, -\frac{b}{A}\right)$ appartient donc au cercle trigonométrique. Par définition du cosinus et du sinus, il existe un réel $\varphi$ tel que
\[ \cos\varphi = \frac{a}{A} \qquad \text{et} \qquad \sin\varphi = -\frac{b}{A}. \]
\textbf{Étape 2 : reconnaître une formule d'addition.} En remplaçant $\frac{a}{A}$ par $\cos\varphi$ et $\frac{b}{A}$ par $-\sin\varphi$ :
\begin{align*}
a\cos X + b\sin X &= A\left(\cos\varphi\cos X - \sin\varphi\sin X\right) \\
&= A\cos(X + \varphi),
\end{align*}
d'après la formule $\cos(X+\varphi) = \cos X\cos\varphi - \sin X\sin\varphi$ vue à la section précédente. Le théorème est démontré. $\square$
\end{experience}

\begin{popfigure}
\begin{center}
\begin{tikzpicture}[scale=2.6]
  % axes
  \draw[->,popdark] (-1.35,0) -- (1.45,0) node[below right] {$\cos$};
  \draw[->,popdark] (0,-1.25) -- (0,1.35) node[above left] {$\sin$};
  % cercle
  \draw[popblue,thick] (0,0) circle (1);
  % point M = (a/A, -b/A), exemple cos phi = √3/2, sin phi = -1/2 -> phi = -30°
  \coordinate (M) at (0.866,-0.5);
  \fill[poporange] (M) circle (0.028);
  \node[poporange,below right=1pt] at (M) {$M\left(\dfrac{a}{A}\,;\,-\dfrac{b}{A}\right)$};
  \draw[poporange,thick] (0,0) -- (M);
  % projetés
  \draw[dashed,popmuted] (M) -- (0.866,0) node[above,popdark] {$\frac{a}{A}$};
  \draw[dashed,popmuted] (M) -- (0,-0.5) node[left,popdark] {$-\frac{b}{A}$};
  % angle phi
  \draw[popgreen,thick,->] (0.45,0) arc (0:-30:0.45);
  \node[popgreen] at (0.62,-0.18) {$\varphi$};
  \node[popblue] at (-0.75,1.08) {$\mathcal{C}$};
\end{tikzpicture}
\end{center}
\legende{Le point $M\left(\frac{a}{A}\,;\,-\frac{b}{A}\right)$ est sur le cercle trigonométrique : c'est l'image du réel $\varphi$. Ici $\cos\varphi = \frac{\sqrt{3}}{2}$ et $\sin\varphi = -\frac{1}{2}$, donc $\varphi = -\frac{\pi}{6}$.}
\end{popfigure}

\begin{aretenir}[Forme équivalente avec le sinus]
Comme $\cos\theta = \sin\left(\frac{\pi}{2} - \theta\right)$ pour tout réel $\theta$, on obtient immédiatement, en posant $\theta = X + \varphi$ :
\[ a\cos X + b\sin X = A\cos(X+\varphi) = A\sin\left(\frac{\pi}{2} - X - \varphi\right). \]
C'est la seconde forme exigée par le programme. Retenir surtout la première : c'est elle que l'on utilise en pratique.
\end{aretenir}

\courssub{Méthode pratique et détermination de $\varphi$}

\begin{methode}[Transformer $a\cos X + b\sin X$ en $A\cos(X+\varphi)$]
\begin{enumerate}
  \item \textbf{Calculer l'amplitude} : $A = \sqrt{a^2+b^2}$.
  \item \textbf{Factoriser} $A$ : $a\cos X + b\sin X = A\left(\dfrac{a}{A}\cos X + \dfrac{b}{A}\sin X\right)$.
  \item \textbf{Identifier} : poser $\cos\varphi = \dfrac{a}{A}$ et $\sin\varphi = -\dfrac{b}{A}$.
  \item \textbf{Placer $\varphi$ sur le cercle trigonométrique} : les signes de $\cos\varphi$ et de $\sin\varphi$ déterminent le quadrant, donc la valeur exacte de $\varphi$ (à $2\pi$ près).
  \item \textbf{Conclure} : $a\cos X + b\sin X = A\cos(X+\varphi)$, puis \emph{vérifier} en développant $A\cos(X+\varphi)$.
\end{enumerate}
\end{methode}

\begin{attention}[Le piège du signe de $\varphi$]
L'erreur la plus fréquente est de poser $\sin\varphi = \dfrac{b}{A}$ au lieu de $\sin\varphi = -\dfrac{b}{A}$. Souviens-toi : la formule d'addition fait apparaître un \emph{moins} :
\[ A\cos(X+\varphi) = A\cos X\cos\varphi \;{\color{popink}-}\; A\sin X\sin\varphi. \]
Pour retrouver $+\,b\sin X$, il faut donc $-A\sin\varphi = b$, c'est-à-dire $\sin\varphi = -\frac{b}{A}$. \textbf{Réflexe de sécurité} : après avoir trouvé $\varphi$, développe toujours $A\cos(X+\varphi)$ et vérifie que tu retombes sur $a\cos X + b\sin X$. Cette vérification prend vingt secondes et sauve des points au Probatoire.
\end{attention}

\begin{exemplebox}[Exemple chiffré : transformer $\sqrt{3}\cos X - \sin X$]
Ici $a = \sqrt{3}$ et $b = -1$.
\begin{enumerate}
  \item $A = \sqrt{(\sqrt{3})^2 + (-1)^2} = \sqrt{3+1} = 2$.
  \item $\sqrt{3}\cos X - \sin X = 2\left(\dfrac{\sqrt{3}}{2}\cos X - \dfrac{1}{2}\sin X\right)$.
  \item On pose $\cos\varphi = \dfrac{\sqrt{3}}{2}$ et $\sin\varphi = -\dfrac{b}{A} = \dfrac{1}{2}$.
  \item Le point $\left(\frac{\sqrt{3}}{2}\,;\,\frac{1}{2}\right)$ est l'image de $\varphi = \dfrac{\pi}{6}$.
  \item Donc $\sqrt{3}\cos X - \sin X = 2\cos\left(X + \dfrac{\pi}{6}\right)$.
\end{enumerate}
\textbf{Vérification} : $2\cos\left(X+\frac{\pi}{6}\right) = 2\left(\frac{\sqrt{3}}{2}\cos X - \frac{1}{2}\sin X\right) = \sqrt{3}\cos X - \sin X$. C'est correct.
\end{exemplebox}

\begin{exemplebox}[Deux autres exemples types à connaître]
\textbf{(a)} $a = b = 1$ : $A = \sqrt{2}$, $\cos\varphi = \frac{1}{\sqrt{2}} = \frac{\sqrt{2}}{2}$ et $\sin\varphi = -\frac{\sqrt{2}}{2}$, donc $\varphi = -\frac{\pi}{4}$ :
\[ \cos X + \sin X = \sqrt{2}\cos\left(X - \frac{\pi}{4}\right). \]
\textbf{(b)} $a = 1$, $b = -\sqrt{3}$ : $A = \sqrt{1+3} = 2$, $\cos\varphi = \frac{1}{2}$ et $\sin\varphi = \frac{\sqrt{3}}{2}$, donc $\varphi = \frac{\pi}{3}$ :
\[ \cos X - \sqrt{3}\sin X = 2\cos\left(X + \frac{\pi}{3}\right). \]
\end{exemplebox}

\courssub{Conséquence immédiate : bornes de l'expression}

\begin{propriete}[Encadrement de $a\cos X + b\sin X$]
Pour tout réel $X$, avec $A = \sqrt{a^2+b^2}$ :
\[ -A \leq a\cos X + b\sin X \leq A. \]
Autrement dit, l'expression $a\cos X + b\sin X$ est \cle{bornée} par $-\sqrt{a^2+b^2}$ et $\sqrt{a^2+b^2}$, et ces bornes sont effectivement atteintes.
\end{propriete}

En effet, $a\cos X + b\sin X = A\cos(X+\varphi)$, et un cosinus est toujours compris entre $-1$ et $1$ ; il suffit de multiplier par $A > 0$. Le maximum $A$ est atteint lorsque $X + \varphi = 0 \pmod{2\pi}$, le minimum $-A$ lorsque $X + \varphi = \pi \pmod{2\pi}$. Ce résultat sera décisif pour discuter l'existence de solutions de l'équation $a\cos X + b\sin X = c$ : si $|c| > \sqrt{a^2+b^2}$, il n'y a \emph{aucune} solution.

\begin{exoresolu}[Transformation complète et encadrement]
\textbf{Énoncé.} On considère $f(X) = \cos X + \sqrt{3}\sin X$ pour $X \in \mathbb{R}$.
\begin{enumerate}
  \item Écrire $f(X)$ sous la forme $A\cos(X+\varphi)$ avec $A > 0$.
  \item En déduire l'encadrement de $f(X)$ et la valeur de $X \in [0\,;\,2\pi]$ pour laquelle $f$ est maximale.
\end{enumerate}
\tcblower
\textbf{Solution détaillée.}
\textbf{1.} On a $a = 1$ et $b = \sqrt{3}$, donc $A = \sqrt{1+3} = 2$. On factorise :
\[ f(X) = 2\left(\frac{1}{2}\cos X + \frac{\sqrt{3}}{2}\sin X\right). \]
On pose $\cos\varphi = \frac{1}{2}$ et $\sin\varphi = -\frac{\sqrt{3}}{2}$. Le point $\left(\frac{1}{2}\,;\,-\frac{\sqrt{3}}{2}\right)$ est dans le quatrième quadrant : c'est l'image de $\varphi = -\frac{\pi}{3}$. D'où
\[ f(X) = 2\cos\left(X - \frac{\pi}{3}\right). \]
\emph{Vérification} : $2\cos\left(X-\frac{\pi}{3}\right) = 2\left(\frac{1}{2}\cos X + \frac{\sqrt{3}}{2}\sin X\right) = \cos X + \sqrt{3}\sin X$. Correct.

\textbf{2.} Comme $-1 \leq \cos\left(X-\frac{\pi}{3}\right) \leq 1$, on obtient $-2 \leq f(X) \leq 2$. Le maximum $2$ est atteint quand $X - \frac{\pi}{3} = 0 \pmod{2\pi}$, soit $X = \frac{\pi}{3}$ dans $[0\,;\,2\pi]$. On retrouve la valeur calculée dans l'activité d'introduction : $f\left(\frac{\pi}{3}\right) = 2$.
\end{exoresolu}

\begin{savaistu}[Additionner deux tensions alternatives]
En électricité — tu le retrouveras en Terminale C — la tension délivrée par ENEO est \emph{alternative sinusoïdale} : elle s'écrit $u(t) = U\sqrt{2}\cos(\omega t + \varphi)$. Lorsque deux tensions de \emph{même fréquence} se superposent, leur somme est de la forme $a\cos(\omega t) + b\sin(\omega t)$. La transformation de cette section prouve que le résultat est encore une tension sinusoïdale de même fréquence, dont on calcule l'amplitude $A = \sqrt{a^2+b^2}$ et la phase $\varphi$. C'est exactement le même calcul que celui que tu viens d'apprendre : les mathématiques de Première éclairent littéralement le réseau électrique du Cameroun !
\end{savaistu}

\begin{exercice}[À toi de jouer]
\begin{enumerate}
  \item Écrire sous la forme $A\cos(X+\varphi)$ : \quad (a) $\cos X - \sin X$ ; \quad (b) $\sqrt{3}\cos X + \sin X$ ; \quad (c) $-\cos X + \sqrt{3}\sin X$.
  \item Donner, sans autre calcul, le maximum et le minimum de $g(X) = 3\cos X - 4\sin X$ sur $\mathbb{R}$.
\end{enumerate}
\end{exercice}

\begin{corrige}[Exercice — Transformation de $a\cos X + b\sin X$]
\textbf{1.(a)} $A = \sqrt{2}$ ; $\cos\varphi = \frac{\sqrt{2}}{2}$, $\sin\varphi = \frac{\sqrt{2}}{2}$, donc $\varphi = \frac{\pi}{4}$ : $\cos X - \sin X = \sqrt{2}\cos\left(X + \frac{\pi}{4}\right)$.

\textbf{1.(b)} $A = 2$ ; $\cos\varphi = \frac{\sqrt{3}}{2}$, $\sin\varphi = -\frac{1}{2}$, donc $\varphi = -\frac{\pi}{6}$ : $\sqrt{3}\cos X + \sin X = 2\cos\left(X - \frac{\pi}{6}\right)$.

\textbf{1.(c)} $A = 2$ ; $\cos\varphi = -\frac{1}{2}$, $\sin\varphi = -\frac{\sqrt{3}}{2}$, donc $\varphi = -\frac{2\pi}{3}$ : $-\cos X + \sqrt{3}\sin X = 2\cos\left(X - \frac{2\pi}{3}\right)$.

\textbf{2.} $A = \sqrt{9+16} = 5$. Donc le maximum de $g$ est $5$ et le minimum est $-5$.
\end{corrige}

Nous disposons désormais de l'outil clé : toute expression $a\cos X + b\sin X$ se ramène à un seul cosinus. Dans les prochaines sections, nous apprendrons d'abord à résoudre les équations de base $\cos X = a$, $\sin X = a$, $\tan X = a$, puis nous combinerons les deux acquis pour résoudre les équations et inéquations du type $a\cos X + b\sin X = c$.

\coursec{Équations \texorpdfstring{$\cos X = a$, $\sin X = a$, $\tan X = a$}{cos X = a, sin X = a, tan X = a}}

Dans la section précédente, nous avons appris à transformer une expression du type $a\cos X + b\sin X$ en une seule fonction sinusoïdale $A\cos(X+\varphi)$. Pourquoi un tel effort ? Précisément parce qu'une fois ramenée à la forme $\cos(X+\varphi) = \dfrac{c}{A}$, une équation devient une \emph{équation trigonométrique de base}, du type $\cos X = a$, $\sin X = a$ ou $\tan X = a$. C'est la résolution de ces équations fondamentales que nous allons maintenant maîtriser. Toute la trigonométrie du Probatoire repose sur elles : retenons-les comme on retient les tables de multiplication.

\courssub{Étude graphique préalable : combien de solutions ?}

Avant tout calcul, posons-nous une question géométrique : résoudre $\cos X = a$, c'est chercher les abscisses des points d'intersection de la courbe $y = \cos x$ avec la droite horizontale $y = a$. De même pour $\sin X = a$ et $\tan X = a$. Cette lecture graphique permet de \emph{prévoir} le nombre de solutions avant de les calculer.

\begin{popfigure}
\begin{center}
\begin{tikzpicture}
\begin{axis}[
  width=0.95\linewidth, height=5.2cm,
  axis lines=middle, xlabel={$x$}, ylabel={$y$},
  xmin=-0.4, xmax=7.2, ymin=-1.6, ymax=1.9,
  xtick={1.5708,3.1416,4.7124,6.2832},
  xticklabels={$\frac{\pi}{2}$,$\pi$,$\frac{3\pi}{2}$,$2\pi$},
  ytick={-1,0.5,1}, yticklabels={$-1$,$\frac12$,$1$},
  legend style={at={(0.98,0.98)},anchor=north east,font=\small}
]
\addplot[popcurve,domain=0:6.2832,samples=200]{cos(deg(x))};
\addlegendentry{$y=\cos x$}
\addplot[poporange,thick,domain=0:6.2832]{0.5};
\addlegendentry{$y=\frac12$}
\addplot[popgreen,thick,domain=0:6.2832]{-1};
\addlegendentry{$y=-1$}
\addplot[poppurple,thick,dashed,domain=0:6.2832]{1.4};
\addlegendentry{$y=1{,}4$}
\end{axis}
\end{tikzpicture}
\end{center}
\legende{Intersections de la courbe $y=\cos x$ avec des droites $y=a$. Pour $a=1{,}4$ : aucune intersection. Pour $a=\frac12$ : deux intersections par période. Pour $a=-1$ : une intersection par période (les deux familles de solutions se confondent).}
\end{popfigure}

\begin{propriete}[Discussion suivant les valeurs de $a$ (admise, justifiée graphiquement)]
\begin{itemize}
\item Équation $\cos X = a$ (resp. $\sin X = a$) :
  \begin{itemize}
  \item si $|a| > 1$ : \cle{aucune solution} (car $-1 \le \cos X \le 1$ et $-1 \le \sin X \le 1$) ;
  \item si $|a| \le 1$ : une \cle{infinité de solutions} dans $\mathbb{R}$, et en général \textbf{deux solutions par intervalle de longueur $2\pi$} (une seule si $a = \pm 1$).
  \end{itemize}
\item Équation $\tan X = a$ : pour \textbf{tout} réel $a$, il existe une infinité de solutions, exactement \textbf{une par intervalle de longueur $\pi$} ne contenant pas de valeur interdite.
\end{itemize}
\end{propriete}

\begin{attention}[Condition d'existence pour la tangente]
$\tan X$ n'existe que si $\cos X \neq 0$, c'est-à-dire $X \neq \dfrac{\pi}{2} + k\pi$ ($k \in \mathbb{Z}$). Avant toute équation faisant intervenir $\tan X$, commence par écrire cette \cle{condition d'existence}. De même, une solution trouvée qui serait de la forme $\frac{\pi}{2}+k\pi$ devra être rejetée.
\end{attention}

\courssub{Théorèmes fondamentaux : les solutions générales}

Le cercle trigonométrique va nous donner les théorèmes, et la symétrie va nous les démontrer.

\begin{experience}[Démonstration géométrique : $\cos X = \cos \alpha$]
Sur le cercle trigonométrique, le cosinus d'un réel $X$ est l'\emph{abscisse} de son point image $M$. Dire que $\cos X = \cos \alpha$, c'est dire que le point image $M$ de $X$ a la même abscisse que le point image $A$ de $\alpha$. Or, sur un cercle, les points ayant une abscisse donnée sont au nombre de deux, \textbf{symétriques par rapport à l'axe des abscisses} : ce sont les points images de $\alpha$ et de $-\alpha$. Comme deux réels ont le même point image si et seulement s'ils diffèrent d'un multiple de $2\pi$, on obtient :
\[ X = \alpha + 2k\pi \quad \text{ou} \quad X = -\alpha + 2k\pi, \quad k \in \mathbb{Z}. \]
Le même raisonnement avec l'\emph{ordonnée} (sinus) donne des points symétriques par rapport à l'axe des ordonnées, images de $\alpha$ et de $\pi - \alpha$. Enfin, la tangente se lit sur l'axe vertical tangent au cercle en $I(1\,;\,0)$ : les points images de $\alpha$ et $\alpha + \pi$ donnent la \emph{même} tangente, d'où une périodicité de $\pi$.
\end{experience}

\begin{propriete}[Théorèmes à connaître par cœur — démonstration géométrique exigible]
Pour tous réels $X$ et $\alpha$ :
\begin{align*}
\cos X = \cos \alpha &\iff X = \alpha + 2k\pi \ \text{ou}\ X = -\alpha + 2k\pi,\ k\in\mathbb{Z} ;\\
\sin X = \sin \alpha &\iff X = \alpha + 2k\pi \ \text{ou}\ X = \pi - \alpha + 2k\pi,\ k\in\mathbb{Z} ;\\
\tan X = \tan \alpha &\iff X = \alpha + k\pi,\ k\in\mathbb{Z} \quad (\text{si } \cos X \neq 0 \text{ et } \cos\alpha \neq 0).
\end{align*}
\end{propriete}

\begin{popfigure}
\begin{center}
\begin{tikzpicture}[scale=2.6]
\draw[popline] (-1.4,0) -- (1.5,0) node[below right]{\small $x$};
\draw[popline] (0,-1.3) -- (0,1.4) node[above left]{\small $y$};
\draw[popblue,thick] (0,0) circle (1);
\coordinate (M) at (0.5,0.866);
\coordinate (M') at (0.5,-0.866);
\draw[poporange,very thick] (0,0) -- (M);
\draw[popgreen,very thick] (0,0) -- (M');
\draw[dashed,popmuted] (0.5,-1) -- (0.5,1);
\fill[poporange] (M) circle (0.035) node[above right]{\small $M$ : image de $\frac{\pi}{3}$};
\fill[popgreen] (M') circle (0.035) node[below right]{\small $M'$ : image de $-\frac{\pi}{3}$};
\draw[popink,thick] (0.5,0.06) -- (0.5,-0.06) node[below left=1pt and 2pt]{\small $\frac12$};
\draw[poporange,->] (0.45,0) arc (0:60:0.45) node[midway,right]{\small $\frac{\pi}{3}$};
\draw[popgreen,->] (0.55,0) arc (0:-60:0.55) node[midway,right]{\small $-\frac{\pi}{3}$};
\end{tikzpicture}
\end{center}
\legende{Équation $\cos X = \frac12$ : les deux points images ont la même abscisse $\frac12$ et sont symétriques par rapport à l'axe des abscisses. D'où $X = \pm\frac{\pi}{3} + 2k\pi$.}
\end{popfigure}

\begin{popfigure}
\begin{center}
\begin{tikzpicture}[scale=2.6]
\draw[popline] (-1.5,0) -- (1.5,0) node[below right]{\small $x$};
\draw[popline] (0,-1.3) -- (0,1.4) node[above left]{\small $y$};
\draw[popblue,thick] (0,0) circle (1);
\coordinate (M) at (0.707,0.707);
\coordinate (M') at (-0.707,0.707);
\draw[poporange,very thick] (0,0) -- (M);
\draw[poppurple,very thick] (0,0) -- (M');
\draw[dashed,popmuted] (-1,0.707) -- (1,0.707);
\fill[poporange] (M) circle (0.035) node[above right]{\small image de $\frac{\pi}{4}$};
\fill[poppurple] (M') circle (0.035) node[above left]{\small image de $\frac{3\pi}{4}=\pi-\frac{\pi}{4}$};
\draw[popink,thick] (-0.06,0.707) -- (0.06,0.707) node[above left=1pt and 1pt]{\small $\frac{\sqrt2}{2}$};
\draw[poporange,->] (0.5,0) arc (0:45:0.5) node[midway,right]{\small $\frac{\pi}{4}$};
\draw[poppurple,->] (0.62,0) arc (0:135:0.62) node[midway,above]{\small $\frac{3\pi}{4}$};
\end{tikzpicture}
\end{center}
\legende{Équation $\sin X = \frac{\sqrt2}{2}$ : les deux points images ont la même ordonnée et sont symétriques par rapport à l'axe des ordonnées. D'où $X = \frac{\pi}{4} + 2k\pi$ ou $X = \frac{3\pi}{4} + 2k\pi$.}
\end{popfigure}

\begin{attention}[Les deux erreurs classiques du Probatoire]
\begin{enumerate}
\item \textbf{Oublier la deuxième famille de solutions.} Pour $\cos X = \cos\alpha$, il y a $X = \alpha + 2k\pi$ \emph{et} $X = -\alpha + 2k\pi$. Pour $\sin X = \sin\alpha$, il y a $X = \alpha + 2k\pi$ \emph{et} $X = \pi - \alpha + 2k\pi$. Une seule famille = la moitié des points au mieux !
\item \textbf{Se tromper de période pour la tangente.} Pour $\tan X = \tan\alpha$, la période est $\cle{\pi}$ et non $2\pi$ : $X = \alpha + k\pi$, une seule famille.
\end{enumerate}
\end{attention}

\begin{aretenir}[Cas particuliers à reconnaître instantanément]
\begin{center}
\begin{tabular}{|c|c|}
\hline
\rowcolor{popblueL} Équation & Solutions ($k \in \mathbb{Z}$) \\ \hline
$\cos X = 0$ & $X = \dfrac{\pi}{2} + k\pi$ \\ \hline
$\cos X = 1$ & $X = 2k\pi$ \\ \hline
$\cos X = -1$ & $X = \pi + 2k\pi$ \\ \hline
$\sin X = 0$ & $X = k\pi$ \\ \hline
$\sin X = 1$ & $X = \dfrac{\pi}{2} + 2k\pi$ \\ \hline
$\sin X = -1$ & $X = -\dfrac{\pi}{2} + 2k\pi$ \\ \hline
$\tan X = 0$ & $X = k\pi$ \\ \hline
\end{tabular}
\end{center}
Dans ces cas, les deux familles de solutions se confondent en une seule.
\end{aretenir}

\courssub{Méthode de résolution et sélection des solutions dans un intervalle borné}

\begin{methode}[Résoudre $\cos X = a$, $\sin X = a$ ou $\tan X = a$]
\begin{enumerate}
\item \textbf{Vérifier l'existence} : si l'équation est $\cos X = a$ ou $\sin X = a$ avec $|a|>1$, conclure $S = \varnothing$. Pour la tangente, écrire la condition d'existence.
\item \textbf{Trouver une solution particulière $\alpha$} : une valeur remarquable ($\frac{\pi}{6}, \frac{\pi}{4}, \frac{\pi}{3}, \frac{2\pi}{3}\dots$) telle que $\cos\alpha = a$ (resp. $\sin\alpha = a$, $\tan\alpha = a$). Attention aux signes : $\cos\alpha = -\frac12$ donne $\alpha = \frac{2\pi}{3}$, pas $\frac{\pi}{3}$ !
\item \textbf{Écrire les familles de solutions} avec les théorèmes (deux familles pour cos et sin, une seule de période $\pi$ pour tan).
\item \textbf{Filtrer selon l'intervalle} : encadrer $k$ pour ne garder que les solutions appartenant à l'intervalle demandé, puis conclure par l'ensemble $S$.
\end{enumerate}
\end{methode}

\begin{exemplebox}[Résolution de $\cos X = -\dfrac{1}{2}$]
\textbf{Dans $\mathbb{R}$ :} on sait que $\cos\dfrac{2\pi}{3} = -\dfrac{1}{2}$. Donc
\[ \cos X = -\frac12 \iff \cos X = \cos\frac{2\pi}{3} \iff X = \frac{2\pi}{3} + 2k\pi \ \text{ou}\ X = -\frac{2\pi}{3} + 2k\pi,\ k\in\mathbb{Z}. \]
\textbf{Dans $[0\,;\,2\pi]$ :} pour la première famille, $0 \le \frac{2\pi}{3} + 2k\pi \le 2\pi$ impose $k = 0$, soit $X = \frac{2\pi}{3}$. Pour la seconde, $0 \le -\frac{2\pi}{3} + 2k\pi \le 2\pi$ impose $k = 1$, soit $X = \frac{4\pi}{3}$. Conclusion : $S = \left\lbrace \dfrac{2\pi}{3}\,;\,\dfrac{4\pi}{3} \right\rbrace$.
\end{exemplebox}

\begin{exoresolu}[{Résoudre $\sin\left(2X - \dfrac{\pi}{4}\right) = \dfrac{\sqrt{3}}{2}$ dans $[0\,;\,\pi]$}]
Résoudre dans $[0\,;\,\pi]$ l'équation $\sin\left(2X - \dfrac{\pi}{4}\right) = \dfrac{\sqrt{3}}{2}$.
\tcblower
\textbf{Étape 1 — Solution particulière.} On sait que $\sin\dfrac{\pi}{3} = \dfrac{\sqrt{3}}{2}$. Posons $\theta = 2X - \dfrac{\pi}{4}$. Alors :
\[ \sin\theta = \sin\frac{\pi}{3} \iff \theta = \frac{\pi}{3} + 2k\pi \ \text{ou}\ \theta = \pi - \frac{\pi}{3} + 2k\pi = \frac{2\pi}{3} + 2k\pi,\quad k\in\mathbb{Z}. \]
\textbf{Étape 2 — Retour à $X$.}
\begin{align*}
\text{Famille 1 : } 2X - \frac{\pi}{4} &= \frac{\pi}{3} + 2k\pi \iff X = \frac{7\pi}{24} + k\pi ;\\
\text{Famille 2 : } 2X - \frac{\pi}{4} &= \frac{2\pi}{3} + 2k\pi \iff X = \frac{11\pi}{24} + k\pi.
\end{align*}
\textbf{Étape 3 — Filtrage sur $[0\,;\,\pi]$.}
Famille 1 : $0 \le \frac{7\pi}{24} + k\pi \le \pi$. Pour $k=0$ : $X = \frac{7\pi}{24} \approx 0{,}29\pi$ convient ; pour $k=1$ : $X = \frac{31\pi}{24} > \pi$, rejeté ; $k=-1$ donne $X<0$, rejeté.
Famille 2 : $0 \le \frac{11\pi}{24} + k\pi \le \pi$. Pour $k=0$ : $X = \frac{11\pi}{24} \approx 0{,}46\pi$ convient ; $k=1$ donne $\frac{35\pi}{24} > \pi$, rejeté.
\textbf{Conclusion :} $S = \left\lbrace \dfrac{7\pi}{24}\,;\,\dfrac{11\pi}{24} \right\rbrace$.
\end{exoresolu}

\begin{attention}[Équations à argument composé : ne pas oublier de « redescendre »]
Dans $\sin(2X - \frac{\pi}{4}) = \frac{\sqrt3}{2}$, le $k\pi$ obtenu pour $X$ (et non $2k\pi$) vient de la division par $2$. De façon générale, si l'argument est $nX + b$, on divise toute la famille par $n$, et la période des solutions en $X$ devient $\frac{2\pi}{n}$ (ou $\frac{\pi}{n}$ pour la tangente). Vérifie toujours que le filtrage en $k$ a bien été fait \emph{après} cette division.
\end{attention}

\courssub{Équations se ramenant aux formes de base}

Beaucoup d'équations du Probatoire ne se présentent pas directement sous la forme $\cos X = a$, mais s'y ramènent par une factorisation ou un changement de variable.

\begin{exoresolu}[Équation factorisée et équation en $\cos^2 X$]
\textbf{1.} Résoudre dans $[0\,;\,2\pi]$ : $2\cos^2 X - 1 = 0$.
\textbf{2.} Résoudre dans $]-\pi\,;\,\pi]$ : $(2\sin X - 1)(\tan X + 1) = 0$.
\tcblower
\textbf{1.} $2\cos^2 X - 1 = 0 \iff \cos^2 X = \dfrac12 \iff \cos X = \dfrac{\sqrt2}{2}$ ou $\cos X = -\dfrac{\sqrt2}{2}$.
\begin{itemize}
\item $\cos X = \frac{\sqrt2}{2} = \cos\frac{\pi}{4}$ : $X = \pm\frac{\pi}{4} + 2k\pi$. Dans $[0\,;\,2\pi]$ : $\frac{\pi}{4}$ et $\frac{7\pi}{4}$.
\item $\cos X = -\frac{\sqrt2}{2} = \cos\frac{3\pi}{4}$ : $X = \pm\frac{3\pi}{4} + 2k\pi$. Dans $[0\,;\,2\pi]$ : $\frac{3\pi}{4}$ et $\frac{5\pi}{4}$.
\end{itemize}
$S = \left\lbrace \frac{\pi}{4}\,;\,\frac{3\pi}{4}\,;\,\frac{5\pi}{4}\,;\,\frac{7\pi}{4} \right\rbrace$.

\textbf{2.} Un produit est nul si et seulement si l'un des facteurs est nul.
\begin{itemize}
\item $2\sin X - 1 = 0 \iff \sin X = \frac12 = \sin\frac{\pi}{6}$ : $X = \frac{\pi}{6} + 2k\pi$ ou $X = \frac{5\pi}{6} + 2k\pi$. Dans $]-\pi\,;\,\pi]$ : $\frac{\pi}{6}$ et $\frac{5\pi}{6}$.
\item $\tan X = -1 = \tan\left(-\frac{\pi}{4}\right)$ : $X = -\frac{\pi}{4} + k\pi$. Dans $]-\pi\,;\,\pi]$ : $-\frac{\pi}{4}$ et $\frac{3\pi}{4}$ (la condition $\cos X \neq 0$ est respectée).
\end{itemize}
$S = \left\lbrace -\frac{\pi}{4}\,;\,\frac{\pi}{6}\,;\,\frac{3\pi}{4}\,;\,\frac{5\pi}{6} \right\rbrace$.
\end{exoresolu}

\begin{methode}[Se ramener aux équations de base]
\begin{enumerate}
\item \textbf{Changement de variable} : pour $\cos 2X = a$, poser $\theta = 2X$, résoudre en $\theta$, puis revenir à $X$ en divisant.
\item \textbf{Factorisation} : transformer l'équation en produit nul $(\cdots)(\cdots) = 0$ (facteur commun, identité $\cos^2 X + \sin^2 X = 1$, formules de duplication).
\item \textbf{Équation du second degré} : pour $2\cos^2 X + \cos X - 1 = 0$, poser $u = \cos X$, résoudre en $u$, vérifier $|u| \le 1$, puis résoudre $\cos X = u$.
\end{enumerate}
\end{methode}

\begin{exercice}[Pour t'entraîner]
\begin{enumerate}
\item Résoudre dans $\mathbb{R}$ puis dans $[0\,;\,2\pi]$ : $\sin X = -\dfrac{\sqrt3}{2}$.
\item Résoudre dans $[0\,;\,\pi]$ : $\tan\left(X - \dfrac{\pi}{6}\right) = 1$.
\item Résoudre dans $[-\pi\,;\,\pi]$ : $\cos 2X = \dfrac12$.
\item Résoudre dans $\mathbb{R}$ : $2\sin^2 X - 3\sin X + 1 = 0$ (poser $u = \sin X$).
\end{enumerate}
\end{exercice}

\begin{corrige}[Exercice]
\textbf{1.} $\sin X = -\frac{\sqrt3}{2} = \sin\left(-\frac{\pi}{3}\right)$, donc $X = -\frac{\pi}{3} + 2k\pi$ ou $X = \pi + \frac{\pi}{3} + 2k\pi = \frac{4\pi}{3} + 2k\pi$. Dans $[0\,;\,2\pi]$ : $S = \left\lbrace \frac{4\pi}{3}\,;\,\frac{5\pi}{3} \right\rbrace$ (pour la première famille, $k=1$ donne $\frac{5\pi}{3}$).

\textbf{2.} Condition : $X - \frac{\pi}{6} \neq \frac{\pi}{2} + k\pi$. $\tan\left(X - \frac{\pi}{6}\right) = 1 = \tan\frac{\pi}{4}$ donne $X - \frac{\pi}{6} = \frac{\pi}{4} + k\pi$, soit $X = \frac{5\pi}{12} + k\pi$. Dans $[0\,;\,\pi]$ : $k=0$ donne $\frac{5\pi}{12}$ ; $k=1$ donne $\frac{17\pi}{12} > \pi$. $S = \left\lbrace \frac{5\pi}{12} \right\rbrace$.

\textbf{3.} Posons $\theta = 2X$ : $\cos\theta = \frac12$ donne $\theta = \pm\frac{\pi}{3} + 2k\pi$, d'où $X = \pm\frac{\pi}{6} + k\pi$. Dans $[-\pi\,;\,\pi]$ : $k \in \{-1\,;\,0\,;\,1\}$ donne $X \in \left\lbrace -\frac{5\pi}{6}\,;\,-\frac{\pi}{6}\,;\,\frac{\pi}{6}\,;\,\frac{5\pi}{6} \right\rbrace$.

\textbf{4.} Avec $u = \sin X$ : $2u^2 - 3u + 1 = 0$, $\Delta = 9 - 8 = 1$, $u = 1$ ou $u = \frac12$ (toutes deux dans $[-1\,;\,1]$). $\sin X = 1$ donne $X = \frac{\pi}{2} + 2k\pi$ ; $\sin X = \frac12$ donne $X = \frac{\pi}{6} + 2k\pi$ ou $X = \frac{5\pi}{6} + 2k\pi$.
\end{corrige}

\begin{savaistu}[Pourquoi ces équations font tourner le monde]
Les équations $\cos X = a$ et $\sin X = a$ sont au cœur du traitement du signal : quand ton téléphone MTN ou Orange démodule une onde radio, il résout en permanence des équations trigonométriques pour retrouver l'information codée dans la phase du signal. Les ingénieurs de CAMTEL utilisent les mêmes outils pour synchroniser les horloges du réseau fibré. Les « familles de solutions en $2k\pi$ » que tu écris aujourd'hui correspondent, dans le monde réel, aux instants où une onde périodique repasse exactement par la même position.
\end{savaistu}

\begin{aretenir}[Bilan de la section]
\begin{itemize}
\item $\cos X = a$ ou $\sin X = a$ : pas de solution si $|a|>1$ ; sinon deux familles de solutions, de période $2\pi$.
\item $\tan X = a$ : toujours des solutions, une seule famille, de période $\pi$, sous la condition $\cos X \neq 0$.
\item Méthode : solution particulière remarquable $\rightarrow$ familles générales $\rightarrow$ filtrage par encadrement de $k$ sur l'intervalle.
\item Les équations complexes se ramènent aux formes de base par changement de variable, factorisation ou équation du second degré en $\cos X$ ou $\sin X$.
\end{itemize}
Ces outils seront indispensables dans la prochaine section, où nous résoudrons les équations du type $a\cos X + b\sin X = c$ grâce à la transformation étudiée précédemment.
\end{aretenir}

\coursec{Équations du type $a\cos X + b\sin X = c$}

\courssub{Pourquoi une nouvelle méthode ?}

Dans la section précédente, tu as appris à résoudre les équations \emph{de base} : $\cos X = a$, $\sin X = a$ et $\tan X = a$. Toute la stratégie consistait à lire les solutions sur le cercle trigonométrique, puis à les exprimer avec les familles $X = \pm\alpha + 2k\pi$ ou $X = \alpha + k\pi$.

Mais dans la réalité — et au Probatoire — les équations ne se présentent pas toujours sous cette forme idéale. Imagine un technicien de la SONAREL qui étudie la tension résultant de la superposition de deux signaux sinusoïdaux : il obtient naturellement une expression du type $\cos X + \sqrt{3}\sin X$, mélange de cosinus et de sinus. Comment trouver les instants $X$ où cette expression vaut une valeur donnée $c$ ?

\begin{experience}[L'obstacle et l'idée directrice]
Essaie de résoudre directement $\cos X + \sqrt{3}\sin X = 1$ : impossible d'isoler $X$, car il apparaît dans \emph{deux} fonctions différentes. L'idée géniale, étudiée dans la section 2, consiste à \cle{fusionner} les deux termes en un seul cosinus déphasé :
\[ a\cos X + b\sin X = A\cos(X+\varphi) \quad \text{avec} \quad A = \sqrt{a^2+b^2}. \]
L'équation devient alors $A\cos(X+\varphi) = c$, c'est-à-dire $\cos(X+\varphi) = \dfrac{c}{A}$ : une équation \emph{de base} en $\cos$, que tu sais résoudre ! Toute cette section repose sur ce détour.
\end{experience}

\courssub{Condition d'existence des solutions}

Avant de te lancer dans les calculs, pose-toi toujours la question : l'équation a-t-elle des solutions ? La réponse est remarquablement simple.

\begin{propriete}[Condition d'existence — à savoir rédiger]
Soient $a$, $b$, $c$ trois réels avec $(a,b) \neq (0,0)$. L'équation
\[ a\cos X + b\sin X = c \]
admet des solutions dans $\mathbb{R}$ \textbf{si et seulement si}
\[ c^2 \leq a^2 + b^2 \qquad \text{c'est-à-dire} \qquad |c| \leq \sqrt{a^2+b^2} = A. \]
\end{propriete}

\begin{experience}[Justification de la condition]
Posons $A = \sqrt{a^2+b^2}$ et écrivons la transformation $a\cos X + b\sin X = A\cos(X+\varphi)$. Or, pour tout réel $t$, on sait que $-1 \leq \cos t \leq 1$. En appliquant cela à $t = X + \varphi$ et en multipliant par $A > 0$ :
\[ -A \leq A\cos(X+\varphi) \leq A \quad \Longleftrightarrow \quad -A \leq a\cos X + b\sin X \leq A. \]
L'expression $a\cos X + b\sin X$ est donc \cle{encadrée} entre $-A$ et $A$, et elle atteint effectivement ces deux bornes (puisque $\cos$ prend toutes les valeurs de $[-1\,;\,1]$). Elle ne peut donc être égale à $c$ que si $c \in [-A\,;\,A]$, c'est-à-dire $|c| \leq A$, soit encore $c^2 \leq a^2+b^2$. \hfill$\blacksquare$
\end{experience}

\begin{aretenir}[À retenir absolument]
L'expression $a\cos X + b\sin X$ vit entièrement dans l'intervalle $\left[-\sqrt{a^2+b^2}\,;\,\sqrt{a^2+b^2}\right]$. Si $c$ est en dehors, l'équation $a\cos X + b\sin X = c$ n'a \textbf{aucune solution} : inutile de chercher !
\end{aretenir}

\begin{exemplebox}[Vérification rapide]
\begin{itemize}
\item L'équation $3\cos X + 4\sin X = 6$ a-t-elle des solutions ? On calcule $A = \sqrt{9+16} = 5$. Comme $|6| = 6 > 5$, il n'y a \textbf{aucune solution}.
\item L'équation $\cos X + \sqrt{3}\sin X = 1$ : $A = \sqrt{1+3} = 2$ et $|1| \leq 2$ : il y a des solutions, nous les trouverons plus bas.
\end{itemize}
\end{exemplebox}

\courssub{La méthode complète de résolution}

\begin{methode}[Résoudre $a\cos X + b\sin X = c$ en 4 étapes]
\begin{enumerate}
\item \textbf{Transformer.} Calculer $A = \sqrt{a^2+b^2}$, puis déterminer un angle $\varphi$ tel que $\cos\varphi = \dfrac{a}{A}$ et $\sin\varphi = -\dfrac{b}{A}$, de sorte que $a\cos X + b\sin X = A\cos(X+\varphi)$.
\emph{Remarque :} on peut aussi écrire $A\cos(X-\varphi')$ avec $\cos\varphi'=\frac{a}{A},\ \sin\varphi'=\frac{b}{A}$ ; les deux formes sont équivalentes ($\varphi'=-\varphi$).
\item \textbf{Vérifier l'existence.} Contrôler que $|c| \leq A$ (sinon : pas de solution, on s'arrête).
\item \textbf{Résoudre l'équation de base.} Écrire $\cos(X+\varphi) = \dfrac{c}{A}$, trouver $\alpha \in [0\,;\,\pi]$ tel que $\cos\alpha = \dfrac{c}{A}$, et conclure :
\[ X + \varphi = \alpha + 2k\pi \quad \text{ou} \quad X + \varphi = -\alpha + 2k\pi, \quad k \in \mathbb{Z}, \]
soit $X = -\varphi + \alpha + 2k\pi$ ou $X = -\varphi - \alpha + 2k\pi$.
\item \textbf{Sélectionner.} Si l'on travaille dans un intervalle donné, tester les valeurs de $k \in \mathbb{Z}$ pour ne garder que les solutions qui appartiennent à cet intervalle.
\end{enumerate}
\end{methode}

\begin{popfigure}
\begin{tikzpicture}[node distance=0.55cm, every node/.style={font=\small},
etape/.style={draw=popblue, fill=popblueL, rounded corners, text width=8.2cm, align=center, minimum height=1cm},
fleche/.style={-{Stealth[length=3mm]}, thick, popdark}]
\node[etape] (e1) {\textbf{Étape 1 :} $A=\sqrt{a^2+b^2}$ ; écrire $a\cos X+b\sin X = A\cos(X+\varphi)$};
\node[etape, below=of e1] (e2) {\textbf{Étape 2 :} vérifier $|c|\leq A$ \; (sinon : aucune solution)};
\node[etape, below=of e2] (e3) {\textbf{Étape 3 :} résoudre $\cos(X+\varphi)=\dfrac{c}{A}$ : $X=-\varphi\pm\alpha+2k\pi$};
\node[etape, below=of e3] (e4) {\textbf{Étape 4 :} sélectionner les solutions dans l'intervalle demandé};
\draw[fleche] (e1) -- (e2);
\draw[fleche] (e2) -- (e3);
\draw[fleche] (e3) -- (e4);
\end{tikzpicture}
\legende{Organigramme de la méthode de résolution de $a\cos X + b\sin X = c$.}
\end{popfigure}

\begin{popfigure}
\begin{tikzpicture}[scale=2.2]
\draw[popline] (-1.35,0) -- (1.35,0);
\draw[popline] (0,-1.35) -- (0,1.35);
\draw[thick, popdark] (0,0) circle (1);
\coordinate (M1) at (40:1);
\coordinate (M2) at (-40:1);
\draw[popline, dashed] (0.766,0) -- (M1);
\draw[popline, dashed] (0.766,0) -- (M2);
\draw[very thick, poporange] (0,0) -- (M1);
\draw[very thick, popgreen] (0,0) -- (M2);
\draw[popblue, very thick] (0.766,-0.12) -- (0.766,0.12);
\fill[popink] (M1) circle (0.025) node[above right=1pt, popink] {$M_1 : X+\varphi=\alpha$};
\fill[popink] (M2) circle (0.025) node[below right=1pt, popink] {$M_2 : X+\varphi=-\alpha$};
\node[popblue, left] at (0.72,0.06) {$\dfrac{c}{A}$};
\draw[popblue, -{Stealth}] (0.35,0) arc (0:40:0.35) node[midway, right, font=\scriptsize] {$\alpha$};
\draw[popblue, -{Stealth}] (0.3,0) arc (0:-40:0.3) node[midway, right, font=\scriptsize] {$-\alpha$};
\node[below left, popmuted, font=\scriptsize] at (0,0) {$O$};
\end{tikzpicture}
\legende{Les points images de $X+\varphi$ dont le cosinus vaut $\dfrac{c}{A}$ : deux points symétriques par rapport à l'axe des cosinus, d'où $X+\varphi = \pm\alpha + 2k\pi$.}
\end{popfigure}

\courssub{Exemple résolu de bout en bout}

\begin{exoresolu}[{Résoudre $\cos X + \sqrt{3}\sin X = 1$ dans $[0\,;\,2\pi]$}]
\textbf{Énoncé.} Résoudre dans $[0\,;\,2\pi]$ l'équation $(E)$ : $\cos X + \sqrt{3}\sin X = 1$.
\tcblower
\textbf{Solution détaillée.}

\textbf{Étape 1 : transformation.} Ici $a = 1$, $b = \sqrt{3}$, donc $A = \sqrt{1^2 + (\sqrt{3})^2} = \sqrt{4} = 2$. On factorise par $2$ :
\[ \cos X + \sqrt{3}\sin X = 2\left(\tfrac{1}{2}\cos X + \tfrac{\sqrt{3}}{2}\sin X\right). \]
On reconnaît $\cos\dfrac{\pi}{3} = \dfrac{1}{2}$ et $\sin\dfrac{\pi}{3} = \dfrac{\sqrt{3}}{2}$. Or $\cos X\cos\dfrac{\pi}{3} + \sin X\sin\dfrac{\pi}{3} = \cos\left(X - \dfrac{\pi}{3}\right)$. Donc :
\[ \cos X + \sqrt{3}\sin X = 2\cos\left(X - \tfrac{\pi}{3}\right). \]
\emph{(Ici $\varphi = -\frac{\pi}{3}$ car $\cos\varphi=\frac{1}{2},\ \sin\varphi=-\frac{\sqrt{3}}{2}$, ce qui donne bien $A\cos(X+\varphi)=2\cos(X-\frac{\pi}{3})$.)}

\textbf{Étape 2 : existence.} $|c| = 1 \leq A = 2$ : il y a des solutions.

\textbf{Étape 3 : équation de base.} $(E) \Longleftrightarrow 2\cos\left(X - \dfrac{\pi}{3}\right) = 1 \Longleftrightarrow \cos\left(X - \dfrac{\pi}{3}\right) = \dfrac{1}{2} = \cos\dfrac{\pi}{3}$. Donc :
\[ X - \tfrac{\pi}{3} = \tfrac{\pi}{3} + 2k\pi \quad \text{ou} \quad X - \tfrac{\pi}{3} = -\tfrac{\pi}{3} + 2k\pi, \quad k \in \mathbb{Z}, \]
c'est-à-dire $X = \dfrac{2\pi}{3} + 2k\pi$ \; ou \; $X = 2k\pi$.

\textbf{Étape 4 : sélection dans $[0\,;\,2\pi]$.}
\begin{itemize}
\item Famille $X = 2k\pi$ : $k=0$ donne $X = 0$ ; $k=1$ donne $X = 2\pi$. Les deux sont dans l'intervalle.
\item Famille $X = \dfrac{2\pi}{3} + 2k\pi$ : $k=0$ donne $X = \dfrac{2\pi}{3}$ ; les autres valeurs de $k$ sortent de l'intervalle.
\end{itemize}

\textbf{Vérification.} Pour $X = 0$ : $\cos 0 + \sqrt{3}\sin 0 = 1$ ✓. Pour $X = \dfrac{2\pi}{3}$ : $\cos\dfrac{2\pi}{3} + \sqrt{3}\sin\dfrac{2\pi}{3} = -\dfrac{1}{2} + \sqrt{3}\times\dfrac{\sqrt{3}}{2} = -\dfrac{1}{2} + \dfrac{3}{2} = 1$ ✓. Pour $X = 2\pi$ : même calcul qu'en $0$ ✓.

\textbf{Conclusion.} $\mathcal{S} = \left\{0\,;\,\dfrac{2\pi}{3}\,;\,2\pi\right\}$.
\end{exoresolu}

\begin{attention}[Précision sur l'énoncé du manuel]
Selon le choix de l'intervalle de résolution, les solutions s'écrivent différemment mais décrivent les mêmes points du cercle. Par exemple, dans $[-\pi\,;\,\pi]$, la famille $X = 2k\pi$ ne donne que $X = 0$, et l'on trouve $\mathcal{S} = \left\{0\,;\,\dfrac{2\pi}{3}\right\}$. Vérifie toujours soigneusement les bornes de l'intervalle demandé : c'est là que se perdent le plus de points au Probatoire.
\end{attention}

\courssub{Cas particuliers à connaître}

\textbf{Premier cas : $c = 0$ (équation homogène).}\par
L'équation $a\cos X + b\sin X = 0$ se résout plus vite sans passer par la transformation complète. En effet :
\[ a\cos X + b\sin X = 0 \Longleftrightarrow A\cos(X+\varphi) = 0 \Longleftrightarrow X + \varphi = \tfrac{\pi}{2} + k\pi, \]
ce qui donne une \emph{seule} famille de solutions, espacées de $\pi$ (et non de $2\pi$). Autre approche : lorsque $\cos X \neq 0$, on peut diviser par $\cos X$ pour obtenir $a + b\tan X = 0$, soit $\tan X = -\dfrac{a}{b}$.

\begin{attention}[Erreur fréquente : la division par $\cos X$]
Diviser $a\cos X + b\sin X = 0$ par $\cos X$ \textbf{sans précaution} est une faute classique. Avant de diviser, il faut vérifier séparément si les réels tels que $\cos X = 0$ (c'est-à-dire $X = \dfrac{\pi}{2} + k\pi$) sont solutions. Pour ces valeurs, $\sin X = \pm 1$, donc l'équation devient $\pm b = 0$ : si $b \neq 0$, elles ne sont \emph{pas} solutions et la division est alors légitime. Mais si tu oublies cette vérification, tu risques d'introduire ou de perdre des solutions. Rédige toujours : « Si $\cos X = 0$, alors ... donc ces valeurs ne conviennent pas. Supposons désormais $\cos X \neq 0$ et divisons... ».
\end{attention}

\textbf{Deuxième cas : $c = A$ ou $c = -A$ (solution unique modulo $2\pi$).}\par
Lorsque $c = A$, l'équation devient $\cos(X+\varphi) = 1$, d'où $X + \varphi = 2k\pi$ : une \cle{unique} famille de solutions $X = -\varphi + 2k\pi$. De même, $c = -A$ donne $\cos(X+\varphi) = -1$, soit $X = \pi - \varphi + 2k\pi$. Contrairement au cas général (deux familles), il n'y a ici qu'une seule solution par période : c'est l'instant où le signal atteint son amplitude maximale (ou minimale).

\begin{exemplebox}[Cas $c = A$]
Résoudre $\cos X + \sqrt{3}\sin X = 2$ dans $\mathbb{R}$. On a vu que le membre de gauche vaut $2\cos\left(X - \dfrac{\pi}{3}\right)$. L'équation s'écrit $\cos\left(X - \dfrac{\pi}{3}\right) = 1$, donc $X - \dfrac{\pi}{3} = 2k\pi$, soit $\mathcal{S} = \left\{\dfrac{\pi}{3} + 2k\pi,\ k \in \mathbb{Z}\right\}$. Une seule famille !
\end{exemplebox}

\courssub{Synthèse et ouverture}

\begin{savaistu}[Superposition de signaux]
Les équations du type $a\cos X + b\sin X = c$ ne sont pas un exercice artificiel : elles modélisent la recherche des instants où la \cle{superposition} de deux signaux sinusoïdaux (deux ondes radio, deux tensions alternatives, deux vibrations sonores) atteint une amplitude donnée. La transformation $A\cos(X+\varphi)$ montre que la somme de deux signaux de même fréquence est encore un signal de même fréquence, d'amplitude $A = \sqrt{a^2+b^2}$ et de déphasage $\varphi$. C'est ce principe qu'utilisent les ingénieurs de MTN ou d'Orange pour analyser les interférences entre antennes relais, et les techniciens d'ENEO pour l'étude des réseaux électriques alternatifs.
\end{savaistu}

\begin{exercice}[À toi de jouer]
\begin{enumerate}
\item Résoudre dans $\mathbb{R}$ puis dans $[0\,;\,2\pi]$ : $\sqrt{3}\cos X - \sin X = 1$.
\item Montrer que l'équation $2\cos X + 3\sin X = 4$ n'a aucune solution.
\item Résoudre dans $]-\pi\,;\,\pi]$ : $\cos X + \sin X = 0$ (deux méthodes possibles).
\end{enumerate}
\end{exercice}

\begin{corrige}[Exercice]
\begin{enumerate}
\item $A = \sqrt{3+1} = 2$ ; $\sqrt{3}\cos X - \sin X = 2\left(\tfrac{\sqrt{3}}{2}\cos X - \tfrac{1}{2}\sin X\right) = 2\cos\left(X + \tfrac{\pi}{6}\right)$. L'équation devient $\cos\left(X + \tfrac{\pi}{6}\right) = \tfrac{1}{2}$, d'où $X + \tfrac{\pi}{6} = \pm\tfrac{\pi}{3} + 2k\pi$, soit $X = \tfrac{\pi}{6} + 2k\pi$ ou $X = -\tfrac{\pi}{2} + 2k\pi$. Dans $[0\,;\,2\pi]$ : $\mathcal{S} = \left\{\tfrac{\pi}{6}\,;\,\tfrac{3\pi}{2}\right\}$.
\item $A = \sqrt{4+9} = \sqrt{13} \approx 3{,}6 < 4$ : la condition $|c| \leq A$ n'est pas vérifiée, donc aucune solution.
\item Pour $\cos X = 0$ ($X = \pm\tfrac{\pi}{2}$), l'équation donne $\sin X = 0$, impossible. On divise alors par $\cos X$ : $1 + \tan X = 0$, soit $\tan X = -1$, d'où $X = -\tfrac{\pi}{4} + k\pi$. Dans $]-\pi\,;\,\pi]$ : $\mathcal{S} = \left\{-\tfrac{\pi}{4}\,;\,\tfrac{3\pi}{4}\right\}$.
\end{enumerate}
\end{corrige}

Tu sais désormais transformer toute équation $a\cos X + b\sin X = c$ en une équation de base, vérifier son existence et en extraire les solutions dans n'importe quel intervalle. La section suivante reprend exactement la même démarche — lecture sur le cercle trigonométrique — mais pour les \emph{inéquations} : il s'agira cette fois de délimiter des arcs entiers du cercle.

\coursec{Inéquations trigonométriques de base}

Dans la section précédente, nous avons appris à résoudre des \emph{équations} : trouver les valeurs \emph{exactes} de $X$ qui vérifient $\cos X = a$, $\sin X = a$ ou $a\cos X + b\sin X = c$. Mais dans la vie courante, on cherche souvent non pas une valeur exacte, mais une \emph{plage} de valeurs. Un technicien de CAMTEL qui oriente une antenne veut que l'angle de visée reste \emph{compris} entre deux bornes ; un agronome de la région de l'Ouest suit une température qui oscille et veut savoir \emph{pendant quelle fraction de la journée} elle dépasse un seuil. Mathématiquement, cela revient à résoudre des \emph{inéquations} trigonométriques : $\cos X \leq a$, $\sin X > a$, $\tan X \geq a$, etc.

L'idée directrice de toute cette section est la suivante : une inéquation se résout \emph{géométriquement} sur le cercle trigonométrique, en deux temps. On résout d'abord l'\cle{équation associée} (qui donne les \cle{points frontières}), puis on détermine \emph{de quel côté} de ces frontières se trouvent les solutions.

\courssub{L'outil fondamental : la constance du signe}

Avant de manipuler les inéquations, il faut comprendre \emph{pourquoi} la méthode « équation puis arc » fonctionne. Tout repose sur la propriété suivante.

\begin{propriete}[Constance du signe (admise)]
Si une expression trigonométrique ne s'annule pas sur un intervalle $I$, alors elle garde un \cle{signe constant} sur $I$ : elle est soit strictement positive sur tout $I$, soit strictement négative sur tout $I$.
\end{propriete}

Cette propriété est \emph{admise} au niveau de la Première (sa démonstration rigoureuse utilise la continuité, notion qui sera précisée en Terminale). Mais son intuition est limpide : le cosinus, le sinus et la tangente varient \emph{sans saut brutal}. Pour passer d'une valeur positive à une valeur négative, une telle expression est \emph{obligée} de franchir la valeur $0$. Donc si elle ne s'annule jamais sur un intervalle, elle ne peut pas y changer de signe.

\begin{exemplebox}[Signe de $\cos X - \tfrac{1}{2}$]
L'expression $\cos X - \dfrac{1}{2}$ s'annule pour $X = \pm \dfrac{\pi}{3} + 2k\pi$. Sur l'intervalle $\left]-\dfrac{\pi}{3}\,;\,\dfrac{\pi}{3}\right[$, elle ne s'annule pas : elle y garde donc un signe constant. Pour le connaître, il suffit de tester \emph{une} valeur, par exemple $X = 0$ : $\cos 0 - \dfrac{1}{2} = \dfrac{1}{2} > 0$. Donc $\cos X - \dfrac{1}{2} > 0$ sur tout l'intervalle $\left]-\dfrac{\pi}{3}\,;\,\dfrac{\pi}{3}\right[$.
\end{exemplebox}

\begin{methode}[Test de signe sur un intervalle]
Pour déterminer le signe d'une expression trigonométrique $E(X)$ sur un intervalle où elle ne s'annule pas :
\begin{enumerate}
\item vérifier que $E$ ne s'annule pas sur cet intervalle (résoudre $E(X)=0$) ;
\item choisir une valeur \emph{simple} $X_0$ de l'intervalle (souvent $0$, $\dfrac{\pi}{6}$, $\dfrac{\pi}{4}$, $\dfrac{\pi}{2}$) ;
\item calculer $E(X_0)$ : son signe est le signe de $E$ sur \emph{tout} l'intervalle.
\end{enumerate}
\end{methode}

\courssub{Lecture géométrique sur le cercle trigonométrique}

Rappelons l'interprétation géométrique essentielle : sur le cercle trigonométrique, le point image $M$ du réel $X$ a pour coordonnées $(\cos X \,;\, \sin X)$.

\begin{aretenir}[Lecture des inéquations sur le cercle]
\begin{itemize}
\item $\cos X > a$ : les points images solutions sont ceux dont l'\emph{abscisse} dépasse $a$, c'est-à-dire l'arc situé \emph{à droite} de la droite verticale d'équation $x = a$.
\item $\sin X > a$ : les points images solutions sont ceux dont l'\emph{ordonnée} dépasse $a$, c'est-à-dire l'arc situé \emph{au-dessus} de la droite horizontale d'équation $y = a$.
\item $\tan X > a$ : on lit sur l'\emph{axe des tangentes} (droite verticale d'équation $x = 1$) : les solutions correspondent aux points $M$ dont la projection sur cet axe, le long de la droite $(OM)$, a une ordonnée supérieure à $a$.
\end{itemize}
Pour des inégalités larges ($\geq$, $\leq$), on \emph{inclut} les points frontières ; pour des inégalités strictes ($>$, $<$), on les \emph{exclut}.
\end{aretenir}

\begin{methode}[Résolution d'une inéquation trigonométrique de base]
Pour résoudre une inéquation du type $\cos X \leq a$ (ou $\sin X$, $\tan X$) sur un intervalle donné :
\begin{enumerate}
\item \textbf{Équation associée} : résoudre $\cos X = a$ pour obtenir les valeurs frontières.
\item \textbf{Points frontières} : placer sur le cercle trigonométrique les points images de ces valeurs.
\item \textbf{Arc solution} : délimiter l'arc (ou les arcs) solution par lecture géométrique (droite verticale pour le cosinus, horizontale pour le sinus, axe des tangentes pour la tangente). En cas de doute, tester une valeur simple grâce à la propriété de constance du signe.
\item \textbf{Traduction en inégalités} : écrire l'ensemble des $X$ correspondants dans l'intervalle demandé, en respectant les bornes strictes ou larges.
\item \textbf{Extension à $\mathbb{R}$} : si la résolution est demandée dans $\mathbb{R}$, ajouter la périodicité : $+\,2k\pi$ ($k \in \mathbb{Z}$) pour cosinus et sinus, $+\,k\pi$ pour la tangente.
\end{enumerate}
\end{methode}

\courssub{Inéquations avec le cosinus}

\begin{exoresolu}[{Résolution de $\cos X \leq \dfrac{1}{2}$ dans $[0\,;\,2\pi]$}]
Résoudre dans $[0\,;\,2\pi]$ l'inéquation $\cos X \leq \dfrac{1}{2}$, puis donner les solutions dans $\mathbb{R}$.
\tcblower
\textbf{Étape 1 : équation associée.} $\cos X = \dfrac{1}{2} = \cos \dfrac{\pi}{3}$, donc $X = \pm \dfrac{\pi}{3} + 2k\pi$. Dans $[0\,;\,2\pi]$, les valeurs frontières sont $X = \dfrac{\pi}{3}$ et $X = 2\pi - \dfrac{\pi}{3} = \dfrac{5\pi}{3}$.

\textbf{Étapes 2 et 3 : points frontières et arc solution.} On trace la droite verticale d'équation $x = \dfrac{1}{2}$. Elle coupe le cercle aux points images de $\dfrac{\pi}{3}$ et $\dfrac{5\pi}{3}$. L'inéquation $\cos X \leq \dfrac{1}{2}$ signifie « abscisse inférieure ou égale à $\dfrac{1}{2}$ » : les points solutions sont sur l'arc situé \emph{à gauche} de cette droite, c'est-à-dire le grand arc passant par le point image de $\pi$ (dont l'abscisse est $-1 \leq \dfrac{1}{2}$). Vérification par la constance du signe : sur $\left]\dfrac{\pi}{3}\,;\,\dfrac{5\pi}{3}\right[$, l'expression $\cos X - \dfrac{1}{2}$ ne s'annule pas ; en $X = \pi$, elle vaut $-1 - \dfrac{1}{2} < 0$ : elle est donc négative sur tout cet intervalle. C'est bien l'arc solution.

\textbf{Étape 4 : traduction.} L'inégalité est large, donc les bornes sont incluses :
\[ S_{[0\,;\,2\pi]} = \left[\dfrac{\pi}{3}\,;\,\dfrac{5\pi}{3}\right]. \]

\textbf{Étape 5 : extension à $\mathbb{R}$.} Par périodicité $2\pi$ :
\[ S_{\mathbb{R}} = \bigcup_{k \in \mathbb{Z}} \left[\dfrac{\pi}{3} + 2k\pi\,;\,\dfrac{5\pi}{3} + 2k\pi\right]. \]
\end{exoresolu}

\begin{popfigure}
\begin{center}
\begin{tikzpicture}[scale=2.6]
\draw[popmuted] (-1.35,0) -- (1.35,0);
\draw[popmuted] (0,-1.25) -- (0,1.25);
\draw[popdark,thick] (0,0) circle (1);
\draw[poporange,thick,dashed] (0.5,-1.2) -- (0.5,1.2);
\node[poporange,right] at (0.5,-1.15) {\footnotesize $x=\tfrac12$};
\draw[popblue,line width=2.2pt] (60:1) arc (60:300:1);
\fill[poppurple] (60:1) circle (0.035);
\fill[poppurple] (300:1) circle (0.035);
\node[above right] at (60:1.02) {\footnotesize $\frac{\pi}{3}$};
\node[below right] at (300:1.02) {\footnotesize $\frac{5\pi}{3}$};
\fill[popdark] (180:1) circle (0.03);
\node[left] at (180:1.02) {\footnotesize $\pi$};
\fill[popdark] (0,0) circle (0.025);
\node[below left] at (0,0) {\footnotesize $O$};
\node[below] at (1.3,0) {\footnotesize $\cos$};
\node[left] at (0,1.2) {\footnotesize $\sin$};
\end{tikzpicture}
\end{center}
\legende{Arc solution (en bleu) de $\cos X \leq \dfrac{1}{2}$ : à gauche de la droite $x = \dfrac{1}{2}$, bornes incluses (points pleins).}
\end{popfigure}

\courssub{Inéquations avec le sinus}

\begin{exoresolu}[{Résolution de $\sin X > -\dfrac{\sqrt{2}}{2}$ dans $[-\pi\,;\,\pi]$}]
Résoudre dans $[-\pi\,;\,\pi]$ l'inéquation $\sin X > -\dfrac{\sqrt{2}}{2}$.
\tcblower
\textbf{Étape 1 : équation associée.} $\sin X = -\dfrac{\sqrt{2}}{2} = \sin\left(-\dfrac{\pi}{4}\right)$, donc $X = -\dfrac{\pi}{4} + 2k\pi$ ou $X = \pi - \left(-\dfrac{\pi}{4}\right) + 2k\pi = \dfrac{5\pi}{4} + 2k\pi$. Dans $[-\pi\,;\,\pi]$, les valeurs frontières sont $X = -\dfrac{\pi}{4}$ et $X = \dfrac{5\pi}{4} - 2\pi = -\dfrac{3\pi}{4}$.

\textbf{Étapes 2 et 3 : arc solution.} On trace la droite horizontale $y = -\dfrac{\sqrt{2}}{2}$. Elle coupe le cercle aux points images de $-\dfrac{3\pi}{4}$ et $-\dfrac{\pi}{4}$. On cherche les points d'ordonnée \emph{strictement supérieure} à $-\dfrac{\sqrt{2}}{2}$ : c'est le grand arc passant par le point image de $\dfrac{\pi}{2}$ (ordonnée $1$). Vérification : sur $\left]-\dfrac{3\pi}{4}\,;\,-\dfrac{\pi}{4}\right[$, l'expression $\sin X + \dfrac{\sqrt{2}}{2}$ ne s'annule pas ; en $X = -\dfrac{\pi}{2}$, elle vaut $-1 + \dfrac{\sqrt{2}}{2} < 0$ : ce petit arc est donc \emph{hors} solution, ce qui confirme que la solution est le reste du cercle.

\textbf{Étape 4 : traduction.} Attention : sur $[-\pi\,;\,\pi]$, le grand arc solution se lit en \emph{deux morceaux} car il franchit la coupure en $\pm\pi$ : de $-\pi$ à $-\dfrac{3\pi}{4}$ (exclu), puis de $-\dfrac{\pi}{4}$ (exclu) à $\pi$. L'inégalité étant stricte, les frontières sont exclues :
\[ S = \left[-\pi\,;\,-\dfrac{3\pi}{4}\right[ \;\cup\; \left]-\dfrac{\pi}{4}\,;\,\pi\right]. \]
\end{exoresolu}

\begin{popfigure}
\begin{center}
\begin{tikzpicture}[scale=2.6]
\draw[popmuted] (-1.35,0) -- (1.35,0);
\draw[popmuted] (0,-1.25) -- (0,1.25);
\draw[popdark,thick] (0,0) circle (1);
\draw[poporange,thick,dashed] (-1.25,-0.7071) -- (1.25,-0.7071);
\node[poporange,right] at (1.05,-0.85) {\footnotesize $y=-\tfrac{\sqrt2}{2}$};
\draw[popblue,line width=2.2pt] (-45:1) arc (-45:225:1);
\draw[popblue,fill=popcream] (-45:1) circle (0.035);
\draw[popblue,fill=popcream] (225:1) circle (0.035);
\node[below right] at (-45:1.02) {\footnotesize $-\frac{\pi}{4}$};
\node[below left] at (225:1.02) {\footnotesize $-\frac{3\pi}{4}$};
\fill[popdark] (90:1) circle (0.03);
\node[above] at (90:1.02) {\footnotesize $\frac{\pi}{2}$};
\fill[popdark] (0,0) circle (0.025);
\node[below left] at (0,0) {\footnotesize $O$};
\end{tikzpicture}
\end{center}
\legende{Arc solution (en bleu) de $\sin X > -\dfrac{\sqrt{2}}{2}$ : au-dessus de la droite $y = -\dfrac{\sqrt{2}}{2}$, bornes exclues (points creux).}
\end{popfigure}

\begin{attention}[{L'arc « en deux morceaux » sur $[-\pi\,;\,\pi]$}]
C'est \emph{l'erreur classique} du Probatoire. Quand l'arc solution franchit le point image de $\pi$ (la « coupure » de l'intervalle $[-\pi\,;\,\pi]$), il se traduit par une \emph{réunion de deux intervalles}, pas par un seul. Écrire $\left]-\dfrac{\pi}{4}\,;\,-\dfrac{3\pi}{4}\right[$ serait doublement faux : les bornes sont dans le mauvais ordre et on a pris le \emph{mauvais} arc (le petit arc, qui est précisément l'ensemble des \emph{non}-solutions). De même, ne jamais orienter l'arc « à l'envers » : vérifie toujours ton choix en testant une valeur simple ($0$, $\dfrac{\pi}{2}$, $\pi$) grâce à la constance du signe.
\end{attention}

\courssub{Inéquations avec la tangente}

La tangente se lit sur l'\cle{axe des tangentes} : la droite verticale d'équation $x = 1$, tangente au cercle au point $I(1\,;\,0)$. Pour un point image $M$ de $X$ (avec $X \neq \dfrac{\pi}{2} + k\pi$), la droite $(OM)$ coupe cet axe en un point $T$ d'ordonnée $\tan X$. Deux différences majeures avec cosinus et sinus :
\begin{itemize}
\item la tangente a pour \emph{période} $\pi$ (et non $2\pi$) : l'extension à $\mathbb{R}$ se fait avec $+\,k\pi$ ;
\item la tangente n'est \emph{pas définie} pour $X = \dfrac{\pi}{2} + k\pi$ : ces valeurs doivent être \emph{exclues} de l'ensemble des solutions, même si l'inégalité est large.
\end{itemize}

\begin{exoresolu}[{Résolution de $\tan X \geq 1$ dans $[0\,;\,\pi]$}]
Résoudre dans $[0\,;\,\pi]$ l'inéquation $\tan X \geq 1$.
\tcblower
\textbf{Condition d'existence :} $X \neq \dfrac{\pi}{2}$ (dans $[0\,;\,\pi]$).

\textbf{Étape 1 : équation associée.} $\tan X = 1 = \tan \dfrac{\pi}{4}$, donc $X = \dfrac{\pi}{4} + k\pi$. Dans $[0\,;\,\pi]$, la seule valeur frontière est $X = \dfrac{\pi}{4}$.

\textbf{Étapes 2 et 3 : lecture sur l'axe des tangentes.} Sur l'axe des tangentes, on marque le point d'ordonnée $1$. Les solutions correspondent aux points $M$ tels que $T$ soit \emph{au-dessus} de cette marque. Sur $\left[0\,;\,\dfrac{\pi}{2}\right[$, la tangente est croissante de $0$ vers $+\infty$ : on obtient $\left[\dfrac{\pi}{4}\,;\,\dfrac{\pi}{2}\right[$. Sur $\left]\dfrac{\pi}{2}\,;\,\pi\right]$, la tangente est négative (elle croît de $-\infty$ vers $0$) : aucune solution. Vérification : $\tan \dfrac{\pi}{3} = \sqrt{3} \geq 1$ (solution) et $\tan \dfrac{2\pi}{3} = -\sqrt{3} < 1$ (non solution).

\textbf{Étape 4 : conclusion.} L'inégalité est large donc $\dfrac{\pi}{4}$ est incluse ; la valeur interdite $\dfrac{\pi}{2}$ est exclue :
\[ S = \left[\dfrac{\pi}{4}\,;\,\dfrac{\pi}{2}\right[. \]
Dans $\mathbb{R}$, par périodicité $\pi$ : $S_{\mathbb{R}} = \displaystyle\bigcup_{k\in\mathbb{Z}} \left[\dfrac{\pi}{4} + k\pi\,;\,\dfrac{\pi}{2} + k\pi\right[$.
\end{exoresolu}

\begin{popfigure}
\begin{center}
\begin{tikzpicture}[scale=2.6]
\draw[popmuted] (-1.35,0) -- (1.6,0);
\draw[popmuted] (0,-1.25) -- (0,1.4);
\draw[popdark,thick] (0,0) circle (1);
\draw[popmuted] (1,-1.25) -- (1,1.4);
\node[below] at (1,-1.25) {\footnotesize axe des tangentes};
\draw[poporange,thick,dashed] (0,0) -- (1,1);
\fill[poporange] (1,1) circle (0.035);
\node[right] at (1,1) {\footnotesize $1$};
\draw[popblue,line width=2.2pt] (45:1) arc (45:90:1);
\fill[poppurple] (45:1) circle (0.035);
\draw[popblue,fill=popcream] (90:1) circle (0.035);
\node[above right] at (45:1.0) {\footnotesize $\frac{\pi}{4}$};
\node[above left] at (90:1.02) {\footnotesize $\frac{\pi}{2}$ (exclu)};
\fill[popdark] (0,0) circle (0.025);
\node[below left] at (0,0) {\footnotesize $O$};
\end{tikzpicture}
\end{center}
\legende{Arc solution de $\tan X \geq 1$ dans $[0\,;\,\pi]$ : de $\dfrac{\pi}{4}$ (inclus) à $\dfrac{\pi}{2}$ (exclu, valeur interdite).}
\end{popfigure}

\begin{attention}[Tangente : les deux pièges mortels]
\begin{itemize}
\item \textbf{Valeurs interdites.} Même pour une inéquation large ($\geq$ ou $\leq$), les valeurs $X = \dfrac{\pi}{2} + k\pi$ ne peuvent \emph{jamais} être solutions, car $\tan X$ n'y existe pas. La borne correspondante est toujours \emph{ouverte}.
\item \textbf{Mauvaise période.} L'extension à $\mathbb{R}$ se fait avec $+\,k\pi$ pour la tangente, pas $+\,2k\pi$. Oublier cela, c'est perdre la moitié des solutions.
\end{itemize}
\end{attention}

\courssub{Inéquations composées : tableaux de signes}

Certaines inéquations font intervenir un \emph{produit} ou un \emph{quotient} d'expressions trigonométriques, par exemple $(2\cos X - 1)\sin X \geq 0$. Grâce à la propriété de constance du signe, on détermine le signe de chaque facteur séparément, puis on conclut par un \cle{tableau de signes}, exactement comme pour les polynômes en Seconde.

\begin{methode}[Tableau de signes trigonométrique]
\begin{enumerate}
\item Résoudre chaque équation associée ($2\cos X - 1 = 0$, $\sin X = 0$, \ldots) dans l'intervalle d'étude.
\item Ranger toutes les valeurs frontières dans l'ordre croissant : elles découpent l'intervalle en sous-intervalles.
\item Sur chaque sous-intervalle, chaque facteur garde un signe constant (propriété admise) : le déterminer par un test ou par lecture sur le cercle.
\item Appliquer la règle des signes pour le produit ou le quotient ; pour un quotient, exclure les valeurs qui annulent le dénominateur (double barre dans le tableau).
\end{enumerate}
\end{methode}

\begin{exoresolu}[{Résolution de $(2\cos X - 1)\sin X \geq 0$ dans $[0\,;\,2\pi]$}]
Résoudre dans $[0\,;\,2\pi]$ l'inéquation $(2\cos X - 1)\sin X \geq 0$.
\tcblower
\textbf{Équations associées.} $2\cos X - 1 = 0 \iff \cos X = \dfrac{1}{2} \iff X \in \left\{\dfrac{\pi}{3}\,;\,\dfrac{5\pi}{3}\right\}$ dans $[0\,;\,2\pi]$, et $\sin X = 0 \iff X \in \{0\,;\,\pi\,;\,2\pi\}$. Valeurs frontières rangées : $0 < \dfrac{\pi}{3} < \pi < \dfrac{5\pi}{3} < 2\pi$.

\textbf{Signes.} Le facteur $2\cos X - 1$ est positif là où $\cos X > \dfrac{1}{2}$, c'est-à-dire sur $\left[0\,;\,\dfrac{\pi}{3}\right[ \cup \left]\dfrac{5\pi}{3}\,;\,2\pi\right]$ (arc à droite de la droite $x = \dfrac12$). Le facteur $\sin X$ est positif sur $]0\,;\,\pi[$, négatif sur $]\pi\,;\,2\pi[$.

\begin{center}
\begin{tabular}{|c|ccccccccc|}\hline
$X$ & $0$ & & $\frac{\pi}{3}$ & & $\pi$ & & $\frac{5\pi}{3}$ & & $2\pi$ \\ \hline
$2\cos X - 1$ & $+$ & $+$ & $0$ & $-$ & $-$ & $-$ & $0$ & $+$ & $+$ \\ \hline
$\sin X$ & $0$ & $+$ & $+$ & $+$ & $0$ & $-$ & $-$ & $-$ & $0$ \\ \hline
Produit & $0$ & $+$ & $0$ & $-$ & $0$ & $+$ & $0$ & $-$ & $0$ \\ \hline
\end{tabular}
\end{center}

\textbf{Conclusion.} On retient les intervalles où le produit est positif ou nul (inégalité large, bornes incluses) :
\[ S = \left[0\,;\,\dfrac{\pi}{3}\right] \cup \left[\pi\,;\,\dfrac{5\pi}{3}\right] \cup \{2\pi\}. \]
Vérification : en $X = \dfrac{\pi}{2}$, $\left(2\cos\dfrac{\pi}{2}-1\right)\sin\dfrac{\pi}{2} = (-1)\times 1 = -1 < 0$ : $\dfrac{\pi}{2}$ n'est pas solution, conforme au tableau.
\end{exoresolu}

\begin{savaistu}[Pourquoi tant d'inéquations en physique ?]
Les grandeurs qui oscillent — tension alternative du réseau ENEO, marées, vibrations d'une corde de mvet — sont modélisées par des expressions du type $A\cos(\omega t + \varphi)$. Savoir \emph{pendant quelle fraction d'une période} la tension dépasse un seuil de sécurité, c'est exactement résoudre une inéquation trigonométrique : la réponse est un rapport de durées, lu comme un rapport d'arcs sur le cercle trigonométrique.
\end{savaistu}

\begin{exercice}[Pour s'entraîner]
\begin{enumerate}
\item Résoudre dans $[0\,;\,2\pi]$ puis dans $\mathbb{R}$ : $\cos X > -\dfrac{\sqrt{3}}{2}$.
\item Résoudre dans $[-\pi\,;\,\pi]$ : $\sin X \leq \dfrac{1}{2}$.
\item Résoudre dans $\left]-\dfrac{\pi}{2}\,;\,\dfrac{\pi}{2}\right[$ : $\tan X < \sqrt{3}$.
\item Résoudre dans $[0\,;\,2\pi]$ : $\sin X \,(2\sin X - 1) < 0$.
\end{enumerate}
\end{exercice}

\begin{corrige}[Exercice]
\begin{enumerate}
\item $\cos X = -\dfrac{\sqrt{3}}{2}$ pour $X = \dfrac{5\pi}{6}$ ou $X = \dfrac{7\pi}{6}$ dans $[0\,;\,2\pi]$. L'arc solution est à droite de la droite $x = -\dfrac{\sqrt{3}}{2}$ (grand arc passant par le point image de $0$), bornes exclues : $S = \left[0\,;\,\dfrac{5\pi}{6}\right[ \cup \left]\dfrac{7\pi}{6}\,;\,2\pi\right]$. Dans $\mathbb{R}$ : $S = \displaystyle\bigcup_{k\in\mathbb{Z}} \left]-\dfrac{5\pi}{6} + 2k\pi\,;\,\dfrac{5\pi}{6} + 2k\pi\right[$.
\item $\sin X = \dfrac{1}{2}$ pour $X = \dfrac{\pi}{6}$ ou $X = \dfrac{5\pi}{6}$ dans $[-\pi\,;\,\pi]$. On veut l'arc \emph{au-dessous} de la droite $y = \dfrac12$, bornes incluses ; cet arc contient le point image de $-\dfrac{\pi}{2}$ et franchit la coupure en $\pm\pi$, donc il se traduit par deux intervalles : $S = \left[-\pi\,;\,\dfrac{\pi}{6}\right] \cup \left[\dfrac{5\pi}{6}\,;\,\pi\right]$.
\item $\tan X = \sqrt{3}$ pour $X = \dfrac{\pi}{3}$. La tangente est strictement croissante sur $\left]-\dfrac{\pi}{2}\,;\,\dfrac{\pi}{2}\right[$, donc $\tan X < \sqrt{3}$ équivaut à $X < \dfrac{\pi}{3}$ : $S = \left]-\dfrac{\pi}{2}\,;\,\dfrac{\pi}{3}\right[$.
\item Le produit est strictement négatif quand les deux facteurs sont de signes contraires. $\sin X = 0$ en $0$, $\pi$, $2\pi$ ; $2\sin X - 1 = 0$ en $\dfrac{\pi}{6}$ et $\dfrac{5\pi}{6}$. Sur $]0\,;\,\pi[$, $\sin X > 0$ et $2\sin X - 1 < 0$ exactement sur $\left]0\,;\,\dfrac{\pi}{6}\right[ \cup \left]\dfrac{5\pi}{6}\,;\,\pi\right[$ ; sur $]\pi\,;\,2\pi[$, $\sin X < 0$ et $2\sin X - 1 < 0$ aussi, donc le produit y est positif. Le tableau de signes donne : $S = \left]0\,;\,\dfrac{\pi}{6}\right[ \cup \left]\dfrac{5\pi}{6}\,;\,\pi\right[$.
\end{enumerate}
\end{corrige}

\begin{aretenir}[Bilan de la section]
\begin{itemize}
\item Une expression trigonométrique qui ne s'annule pas sur un intervalle y garde un \emph{signe constant} (propriété admise) : c'est le fondement de toutes les méthodes de cette section.
\item Méthode en cinq étapes : équation associée $\rightarrow$ points frontières sur le cercle $\rightarrow$ arc solution (lecture géométrique ou test de signe) $\rightarrow$ inégalités sur $X$ $\rightarrow$ extension par périodicité ($2k\pi$ pour cos et sin, $k\pi$ pour tan).
\item Inégalités larges : bornes incluses ; strictes : bornes exclues ; pour la tangente, les valeurs $\dfrac{\pi}{2} + k\pi$ sont toujours exclues.
\item Sur $[-\pi\,;\,\pi]$, un arc solution peut se traduire par une \emph{réunion de deux intervalles}.
\item Pour un produit ou un quotient, on utilise un tableau de signes, comme en Seconde.
\end{itemize}
\end{aretenir}

Dans la section suivante, nous combinerons ces techniques avec la transformation de $a\cos X + b\sin X$ étudiée précédemment pour résoudre des inéquations plus élaborées, notamment graphiquement.

\coursec{Inéquations avec $a\cos X + b\sin X$ et synthèse}

Dans la section précédente, tu as appris à résoudre les inéquations trigonométriques de base : $\cos X \leq a$, $\sin X > a$, $\tan X \geq a$, en t'appuyant sur le cercle trigonométrique. Mais dans la pratique — en physique des vibrations, en électricité, ou dans les problèmes de Probatoire — l'expression qui apparaît est souvent un mélange de cosinus et de sinus : $a\cos X + b\sin X$. Comment comparer une telle expression à un réel $c$ ? Toute la stratégie repose sur une idée que tu maîtrises déjà depuis la section 2 : \textbf{transformer} $a\cos X + b\sin X$ en une seule fonction trigonométrique $A\cos(X+\varphi)$, puis se ramener à une inéquation de base par un changement de variable. C'est exactement le même raisonnement que pour les équations, avec une vigilance supplémentaire : l'\emph{intervalle} d'étude se déplace lui aussi lors du changement de variable. Cette dernière section te donne la méthode complète, algébrique puis graphique, avant de dresser le bilan de toute la leçon.

\courssub{Résolution algébrique : transformer, changer de variable, revenir}

\textbf{Intuition.} Résoudre $\cos X + \sin X \geq 1$ directement semble difficile : cosinus et sinus varient « en décalé ». Mais on sait (section 2) que $\cos X + \sin X = \sqrt{2}\cos\left(X - \dfrac{\pi}{4}\right)$. L'inéquation devient alors $\cos\left(X - \dfrac{\pi}{4}\right) \geq \dfrac{1}{\sqrt{2}}$, c'est-à-dire une inéquation \emph{de base} en la variable $t = X - \dfrac{\pi}{4}$. On résout en $t$, puis on revient à $X$.

\begin{methode}[Inéquation en $a\cos X + b\sin X$ : la méthode en cinq étapes]
Pour résoudre une inéquation du type $a\cos X + b\sin X \;\square\; c$ (où $\square$ désigne $<$, $\leq$, $>$ ou $\geq$) sur un intervalle $I$ :
\begin{enumerate}
\item \textbf{Transformer} : écrire $a\cos X + b\sin X = A\cos(X + \varphi)$ avec $A = \sqrt{a^2+b^2}$ et $\varphi$ tel que $\cos\varphi = \dfrac{a}{A}$ et $\sin\varphi = -\dfrac{b}{A}$ (ou la forme $A\sin(X+\psi)$ si elle est plus commode).
\item \textbf{Changer de variable} : poser $t = X + \varphi$ et \textbf{transporter l'intervalle} : si $X \in I = [\alpha\,;\,\beta]$, alors $t \in [\alpha + \varphi\,;\,\beta + \varphi]$.
\item \textbf{Résoudre en $t$} : l'inéquation devient $\cos t \;\square\; \dfrac{c}{A}$ (attention au sens de l'inégalité : $A > 0$, donc le sens est conservé). On la résout sur le cercle trigonométrique, \emph{dans l'intervalle des $t$}.
\item \textbf{Revenir à $X$} : remplacer $t$ par $X + \varphi$ dans chaque solution, isoler $X$, puis \textbf{intersecter avec l'intervalle $I$ de départ}.
\item \textbf{Vérifier} : tester au moins une valeur de l'ensemble trouvé (et une valeur hors ensemble) dans l'inéquation de départ.
\end{enumerate}
\end{methode}

\begin{attention}[Le changement de variable décale l'intervalle d'étude]
C'est \textbf{l'erreur la plus fréquente} au Probatoire. Quand tu poses $t = X - \dfrac{\pi}{4}$, les bornes changent : si $X \in [0\,;\,2\pi]$, alors $t \in \left[-\dfrac{\pi}{4}\,;\,\dfrac{7\pi}{4}\right]$, et \emph{non pas} $t \in [0\,;\,2\pi]$. Résoudre en $t$ sur le mauvais intervalle fait perdre ou ajouter des solutions. Écris systématiquement le nouvel intervalle \emph{avant} de résoudre, et n'oublie pas à la fin d'intersecter avec l'intervalle imposé à $X$.
\end{attention}

\courssub{Exemple complet : $\cos X + \sin X \geq 1$ dans $[0\,;\,2\pi]$}

Traitons intégralement un exemple, en suivant la méthode pas à pas.

\textbf{Étape 1 — Transformation.} Avec $a = b = 1$ : $A = \sqrt{1^2+1^2} = \sqrt{2}$. On factorise :
\[
\cos X + \sin X = \sqrt{2}\left(\frac{1}{\sqrt{2}}\cos X + \frac{1}{\sqrt{2}}\sin X\right) = \sqrt{2}\cos\left(X - \frac{\pi}{4}\right),
\]
car $\cos\dfrac{\pi}{4} = \sin\dfrac{\pi}{4} = \dfrac{\sqrt{2}}{2} = \dfrac{1}{\sqrt{2}}$.

\textbf{Étape 2 — Changement de variable.} Posons $t = X - \dfrac{\pi}{4}$. Comme $X \in [0\,;\,2\pi]$, on a :
\[
t \in \left[-\frac{\pi}{4}\,;\,2\pi - \frac{\pi}{4}\right] = \left[-\frac{\pi}{4}\,;\,\frac{7\pi}{4}\right].
\]
L'inéquation devient : $\sqrt{2}\cos t \geq 1$, c'est-à-dire $\cos t \geq \dfrac{1}{\sqrt{2}} = \dfrac{\sqrt{2}}{2}$.

\textbf{Étape 3 — Résolution en $t$.} Sur le cercle trigonométrique, $\cos t \geq \dfrac{\sqrt{2}}{2}$ signifie que le point image est à droite de la droite verticale d'abscisse $\dfrac{\sqrt{2}}{2}$. Les valeurs de référence sont $t = -\dfrac{\pi}{4}$ et $t = \dfrac{\pi}{4}$ (angles dont le cosinus vaut $\dfrac{\sqrt{2}}{2}$). Sur l'intervalle $\left[-\dfrac{\pi}{4}\,;\,\dfrac{7\pi}{4}\right]$, on obtient :
\[
t \in \left[-\frac{\pi}{4}\,;\,\frac{\pi}{4}\right].
\]
(La borne $-\dfrac{\pi}{4}$ est incluse car l'inégalité est large, et elle appartient à l'intervalle des $t$.)
Mais attention : l'intervalle des $t$ se termine en $\dfrac{7\pi}{4}$, ce qui correspond à $X = 2\pi$. Le point $X = 2\pi$ donne $t = \dfrac{7\pi}{4}$, et $\cos\dfrac{7\pi}{4} = \dfrac{\sqrt{2}}{2}$ : l'inégalité large est vérifiée ! En fait, sur $\left[-\dfrac{\pi}{4}\,;\,\dfrac{7\pi}{4}\right]$, la solution complète en $t$ est $\left[-\dfrac{\pi}{4}\,;\,\dfrac{\pi}{4}\right] \cup \left\{\dfrac{7\pi}{4}\right\}$ (le cosinus redescend sous $\frac{\sqrt2}{2}$ sur $\left]\frac{\pi}{4}\,;\,\frac{7\pi}{4}\right[$ et ne remonte qu'au point final).

\textbf{Étape 4 — Retour à $X$.} Comme $t = X - \dfrac{\pi}{4}$, on a $X = t + \dfrac{\pi}{4}$, donc :
\[
-\frac{\pi}{4} \leq t \leq \frac{\pi}{4} \iff 0 \leq X \leq \frac{\pi}{2},
\]
et $t = \dfrac{7\pi}{4} \iff X = 2\pi$.
On intersecte avec $[0\,;\,2\pi]$ : cela donne $X \in \left[0\,;\,\dfrac{\pi}{2}\right] \cup \{2\pi\}$.

\textbf{Étape 5 — Vérification.} Pour $X = \dfrac{\pi}{4}$ (dans $S$) : $\cos\dfrac{\pi}{4} + \sin\dfrac{\pi}{4} = \sqrt{2} \geq 1$ : vrai. Pour $X = \pi$ (hors $S$) : $\cos\pi + \sin\pi = -1 \geq 1$ : faux. Pour $X = 2\pi$ : $1 + 0 = 1 \geq 1$ : vrai (inégalité large). La solution est cohérente.

\begin{exemplebox}[Résultat à retenir]
L'ensemble des solutions de $\cos X + \sin X \geq 1$ dans $[0\,;\,2\pi]$ est :
\[
S = \left[0\,;\,\frac{\pi}{2}\right] \cup \{2\pi\}.
\]
Le point isolé $2\pi$ s'explique par la périodicité : c'est le « même point » que $0$ sur le cercle, mais il appartient bien à l'intervalle fermé $[0\,;\,2\pi]$.
\end{exemplebox}

\begin{attention}[Toujours vérifier une solution par substitution]
Après une résolution, prends l'habitude de substituer dans l'inéquation \emph{de départ} (pas dans l'inéquation transformée) une valeur de ton ensemble solution et une valeur extérieure. Cette vérification de dix secondes détecte immédiatement les erreurs de signe, les oublis de points isolés (comme $2\pi$ ci-dessus) et les mauvais transports d'intervalle. Au Probatoire, c'est un réflexe qui sauve des points.
\end{attention}

\courssub{Résolution graphique et cohérence des deux approches}

\textbf{Idée.} On trace sur le même repère la courbe de $f(X) = a\cos X + b\sin X$ et la droite horizontale $y = c$. L'ensemble solution de $f(X) \geq c$ est la réunion des intervalles où la courbe est \emph{au-dessus} (ou sur) la droite. Cette lecture graphique ne remplace pas le calcul exact, mais elle permet de \textbf{contrôler} le résultat algébrique : nombre d'intervalles, positions approximatives des bornes.

\begin{popfigure}
\begin{center}
\begin{tikzpicture}
\begin{axis}[
  width=0.95\linewidth, height=6.2cm,
  axis lines=middle,
  xlabel={$X$}, ylabel={$y$},
  xmin=-0.2, xmax=6.6, ymin=-1.8, ymax=1.8,
  xtick={0,1.5708,3.1416,4.7124,6.2832},
  xticklabels={$0$,$\frac{\pi}{2}$,$\pi$,$\frac{3\pi}{2}$,$2\pi$},
  ytick={-1,1},
  domain=0:6.2832, samples=200, clip=false
]
\addplot[popcurve, thick]{cos(deg(x)) + sin(deg(x))};
\addplot[poporange, thick, dashed, domain=0:6.2832]{1};
\node[popink] at (axis cs:3.6,1.35) {\small $y = \cos X + \sin X$};
\node[poporange] at (axis cs:5.6,0.72) {\small $y = 1$};
\addplot[only marks, mark=*, popdark, mark size=1.6pt] coordinates {(0,1) (1.5708,1) (6.2832,1)};
\end{axis}
\end{tikzpicture}
\end{center}
\legende{Courbe de $f(X) = \cos X + \sin X$ sur $[0\,;\,2\pi]$ et droite $y = 1$. La courbe est au-dessus de la droite exactement sur $\left[0\,;\,\frac{\pi}{2}\right]$ et au point $X = 2\pi$ : on retrouve $S = \left[0\,;\,\frac{\pi}{2}\right] \cup \{2\pi\}$.}
\end{popfigure}

\begin{exoresolu}[Vérification croisée algébrique / graphique]
\textbf{Énoncé.} On a résolu algébriquement $\cos X + \sin X \geq 1$ dans $[0\,;\,2\pi]$ et trouvé $S = \left[0\,;\,\dfrac{\pi}{2}\right] \cup \{2\pi\}$. Confirmer graphiquement et par le cercle trigonométrique.
\tcblower
\textbf{Solution.} Sur la courbe ci-dessus, les intersections avec $y = 1$ se lisent en $X = 0$, $X = \dfrac{\pi}{2} \approx 1{,}57$ et $X = 2\pi \approx 6{,}28$. Entre $0$ et $\dfrac{\pi}{2}$, la courbe culmine à $\sqrt{2} \approx 1{,}41$ (en $X = \dfrac{\pi}{4}$) : elle est bien au-dessus de $1$. Entre $\dfrac{\pi}{2}$ et $2\pi$, elle passe sous la droite (minimum $-\sqrt{2}$ en $X = \dfrac{5\pi}{4}$) et ne la rejoint qu'en $2\pi$. Le cercle trigonométrique confirme : avec $t = X - \dfrac{\pi}{4}$, la condition $\cos t \geq \dfrac{\sqrt{2}}{2}$ délimite l'arc $\left[-\dfrac{\pi}{4}\,;\,\dfrac{\pi}{4}\right]$, qui correspond aux points du cercle situés à droite de la verticale d'abscisse $\dfrac{\sqrt{2}}{2}$. Les deux approches donnent le même ensemble : la solution est validée.
\end{exoresolu}

\begin{popfigure}
\begin{center}
\begin{tikzpicture}[scale=2.6]
\draw[popline] (-1.35,0) -- (1.35,0) node[below right] {\small cos};
\draw[popline] (0,-1.35) -- (0,1.35) node[above left] {\small sin};
\draw[popink, thick] (0,0) circle (1);
% droite verticale cos = sqrt2/2
\draw[poporange, thick, dashed] (0.7071,-1.25) -- (0.7071,1.25);
\node[poporange, below] at (0.7071,-1.25) {\small $\frac{\sqrt{2}}{2}$};
% arc solution en t
\draw[popgreen, very thick] (0.7071,-0.7071) arc[start angle=-45, end angle=45, radius=1];
% points remarquables
\fill[popdark] (0.7071,0.7071) circle (0.028) node[above right, popdark] {\small $t=\frac{\pi}{4}$};
\fill[popdark] (0.7071,-0.7071) circle (0.028) node[below right, popdark] {\small $t=-\frac{\pi}{4}$};
\fill[popblue] (1,0) circle (0.028) node[right, popblue] {\small $t=0$};
\draw[->, popmuted] (0.35,0) arc[start angle=0, end angle=-45, radius=0.35];
\node[popmuted] at (0.5,-0.22) {\small $-\frac{\pi}{4}$};
\end{tikzpicture}
\end{center}
\legende{Cercle trigonométrique pour la variable $t = X - \dfrac{\pi}{4}$ : la condition $\cos t \geq \dfrac{\sqrt{2}}{2}$ sélectionne l'arc épais $\left[-\dfrac{\pi}{4}\,;\,\dfrac{\pi}{4}\right]$, à droite de la droite verticale. En revenant à $X = t + \dfrac{\pi}{4}$, cet arc devient $\left[0\,;\,\dfrac{\pi}{2}\right]$.}
\end{popfigure}

\courssub{Tableau de synthèse de la leçon}

\begin{center}
\renewcommand{\arraystretch}{1.5}
\begin{tabularx}{\linewidth}{|l|X|X|}
\hline
\rowcolor{popblueL} \textbf{Notion} & \textbf{Formule / résultat} & \textbf{Condition, remarque} \\
\hline
Addition & $\cos(a\pm b) = \cos a\cos b \mp \sin a\sin b$ ; $\sin(a\pm b) = \sin a\cos b \pm \cos a\sin b$ & Valables pour tous réels $a, b$ \\
\hline
Tangente & $\tan(a+b) = \dfrac{\tan a + \tan b}{1 - \tan a\tan b}$ & Si toutes les tangentes existent et $1 - \tan a\tan b \neq 0$ \\
\hline
Duplication & $\cos 2a = \cos^2 a - \sin^2 a = 2\cos^2 a - 1$ ; $\sin 2a = 2\sin a\cos a$ & Utile pour $\cos^2 a = \dfrac{1+\cos 2a}{2}$ \\
\hline
Transformation & $a\cos X + b\sin X = A\cos(X+\varphi)$, $A = \sqrt{a^2+b^2}$ & $\cos\varphi = \dfrac{a}{A}$, $\sin\varphi = -\dfrac{b}{A}$ \\
\hline
$\cos X = a$ & $X = \pm\alpha + 2k\pi$ si $|a| \leq 1$, avec $\cos\alpha = a$ & Aucune solution si $|a| > 1$ \\
\hline
$\sin X = a$ & $X = \alpha + 2k\pi$ ou $X = \pi - \alpha + 2k\pi$ si $|a| \leq 1$ & Aucune solution si $|a| > 1$ \\
\hline
$\tan X = a$ & $X = \alpha + k\pi$, avec $\tan\alpha = a$ & Solution pour tout $a \in \mathbb{R}$ \\
\hline
$a\cos X + b\sin X = c$ & Se ramène à $\cos t = \dfrac{c}{A}$, $t = X + \varphi$ & Solutions si $|c| \leq A$ \\
\hline
Inéquations & On lit les arcs sur le cercle trigonométrique, puis on ajoute $2k\pi$ & Signe constant sur un intervalle sans annulation (propriété admise) \\
\hline
\end{tabularx}
\end{center}

\courssub{Bilan des méthodes, erreurs à éviter et ouverture}

\textbf{Bilan des méthodes.} Cette leçon t'a donné une boîte à outils complète : les formules d'addition et de duplication pour \emph{calculer} et \emph{transformer} ; la réduction de $a\cos X + b\sin X$ pour \emph{unifier} ; le cercle trigonométrique pour \emph{résoudre} équations et inéquations ; et la vérification graphique pour \emph{contrôler}. La démarche générale est toujours la même : transformer pour se ramener à une situation de base, résoudre dans la bonne variable et le bon intervalle, revenir à la variable initiale, vérifier.

\begin{attention}[Les cinq erreurs classiques de la leçon]
\begin{enumerate}
\item Oublier de \textbf{transporter l'intervalle} lors du changement de variable $t = X + \varphi$.
\item Confondre les solutions de $\sin X = a$ (deux familles : $\alpha$ \emph{et} $\pi - \alpha$) avec celles de $\cos X = a$ ($\pm\alpha$).
\item Perdre les \textbf{points isolés} aux bornes d'un intervalle fermé (comme $X = 2\pi$ dans notre exemple).
\item Diviser par une expression trigonométrique sans discuter son signe ou son annulation (par exemple diviser par $\cos X$ en oubliant le cas $\cos X = 0$).
\item Ne pas vérifier par substitution : une erreur de signe dans $\varphi$ fausse tout l'ensemble final.
\end{enumerate}
\end{attention}

\begin{savaistu}[Et en Terminale ?]
Tout ce que tu viens d'apprendre deviendra un outil quotidien en Terminale C, D, E et TI. Tu y étudieras les \emph{fonctions trigonométriques} $X \mapsto a\cos X + b\sin X$ comme de véritables fonctions : dérivation ($(\cos X)' = -\sin X$, $(\sin X)' = \cos X$), tableaux de variations, limites, et résolution d'équations différentielles des oscillateurs (ressort, circuit électrique). La transformation $a\cos X + b\sin X = A\cos(X+\varphi)$ y apparaîtra sous le nom de \emph{mise sous forme canonique d'un signal sinusoïdal} : $A$ est l'amplitude, $\varphi$ la phase. C'est exactement le langage utilisé par les ingénieurs de CAMTEL ou d'ENEO pour décrire les signaux alternatifs qui parcourent les réseaux du Cameroun. Ce complément sera traité en Terminale — tu es déjà prêt.
\end{savaistu}

\begin{exercice}[Pour t'entraîner]
\begin{enumerate}
\item Résoudre dans $[0\,;\,2\pi]$ l'inéquation $\cos X - \sqrt{3}\sin X \leq 1$ (indication : écrire le membre de gauche sous la forme $2\cos\left(X + \dfrac{\pi}{3}\right)$).
\item Résoudre dans $[-\pi\,;\,\pi]$ : $\sin X + \cos X > 0$.
\item Déterminer graphiquement, à l'aide d'une courbe bien tracée, le nombre de solutions de $\sqrt{3}\cos X + \sin X = 1$ dans $[0\,;\,2\pi]$, puis confirmer par le calcul.
\end{enumerate}
\end{exercice}

\begin{corrige}[Exercice — indications]
\begin{enumerate}
\item $\cos X - \sqrt{3}\sin X = 2\cos\left(X + \dfrac{\pi}{3}\right)$. On pose $t = X + \dfrac{\pi}{3} \in \left[\dfrac{\pi}{3}\,;\,\dfrac{7\pi}{3}\right]$ ; on résout $\cos t \leq \dfrac{1}{2}$, soit $t \in \left[\dfrac{\pi}{3}\,;\,\dfrac{5\pi}{3}\right] \cup \left\{\dfrac{7\pi}{3}\right\}$ (car $\cos\frac{7\pi}{3}=\frac12$ et l'inégalité est large). Retour : $X \in \left[0\,;\,\dfrac{4\pi}{3}\right] \cup \{2\pi\}$.
\item $\sin X + \cos X = \sqrt{2}\cos\left(X - \dfrac{\pi}{4}\right) > 0$ avec $t = X - \dfrac{\pi}{4} \in \left[-\dfrac{5\pi}{4}\,;\,\dfrac{3\pi}{4}\right]$ ; $\cos t > 0$ pour $t \in \left(-\dfrac{\pi}{2}\,;\,\dfrac{\pi}{2}\right)$, d'où $X \in \left(-\dfrac{\pi}{4}\,;\,\dfrac{3\pi}{4}\right)$.
\item $\sqrt{3}\cos X + \sin X = 2\cos\left(X - \dfrac{\pi}{6}\right) = 1$ donne $\cos t = \dfrac{1}{2}$ avec $t \in \left[-\dfrac{\pi}{6}\,;\,\dfrac{11\pi}{6}\right]$ : deux solutions, $X = \dfrac{\pi}{2}$ et $X = \dfrac{11\pi}{6}$, confirmées par les deux intersections de la courbe avec la droite $y = 1$.
\end{enumerate}
\end{corrige}

\newpage

\coursec{Activité d'intégration}

\coursec{Leçon 3 : Trigonométrie — Formules, équations et inéquations}

\courssub{Activité d'introduction : Le panneau solaire de Maroua}

\courssub{Objectifs de l'activité}

À l'issue de cette activité, tu seras capable de :
\begin{itemize}
\item transformer une expression du type $a\cos X+b\sin X$ en $A\cos(X-\varphi)$ dans un contexte concret ;
\item exploiter cette écriture pour déterminer le maximum d'une grandeur périodique et l'instant où il est atteint ;
\item résoudre une équation du type $a\cos X+b\sin X=c$ sur un intervalle borné ;
\item résoudre une inéquation du type $a\cos X+b\sin X\geq c$ et interpréter l'ensemble des solutions comme une plage horaire ;
\item confronter les résultats algébriques à une vérification graphique, et développer ainsi l'esprit critique vis-à-vis d'un modèle mathématique.
\end{itemize}

\courssub{Situation}

Dans la région de l'Extrême-Nord, l'électrification rurale repose de plus en plus sur l'énergie solaire. À Maroua, une coopérative agricole a installé un panneau photovoltaïque orienté plein sud pour alimenter une pompe à eau destinée à l'irrigation d'un champ de niébé. L'ensoleillement y est remarquable, mais la puissance délivrée varie au cours de la journée : nulle la nuit, elle croît le matin, atteint un sommet autour de midi, puis décroît le soir.

Un technicien de la coopérative a mesuré la puissance reçue et propose le modèle suivant, valable pour une journée type :
\[ P(t)=120\cos\left(\frac{\pi t}{12}\right)+120\sqrt{3}\,\sin\left(\frac{\pi t}{12}\right) \]
où :
\begin{itemize}
\item $t$ désigne l'heure de la journée, en heures, avec $0\leq t\leq 24$ ;
\item $t=0$ correspond à minuit et $t=12$ à midi ;
\item $P(t)$ est la puissance reçue par le panneau, exprimée en watts (W).
\end{itemize}

La pompe à eau de la coopérative exige une puissance d'au moins $120$ W pour fonctionner. Le technicien se pose alors quatre questions :
\begin{itemize}
\item À quelle heure de la journée la puissance reçue est-elle maximale, et quelle est cette puissance maximale ?
\item À quelles heures la puissance vaut-elle exactement $120$ W ?
\item Pendant quelle(s) plage(s) horaire(s) la pompe peut-elle fonctionner ?
\item Le modèle est-il cohérent avec l'observation (nuit sans production) ?
\end{itemize}

\courssub{Consignes}

Travail en groupes de trois ou quatre élèves. Chaque groupe rédige un compte rendu.

\begin{enumerate}
\item \textbf{Mise en forme de l'expression.} On pose $X=\dfrac{\pi t}{12}$. En utilisant la méthode de transformation de $a\cos X+b\sin X$, montrer que l'on peut écrire :
\[ P(t)=240\cos\left(\frac{\pi t}{12}-\frac{\pi}{3}\right). \]
On précisera les valeurs de $A$ et de $\varphi$ et on justifiera chaque étape du calcul.
\item \textbf{Maximum de puissance.} En utilisant le fait que, pour tout réel $u$, $-1\leq\cos u\leq 1$, déterminer la puissance maximale $P_{\max}$ et l'heure $t_0$ à laquelle elle est atteinte. Le résultat est-il conforme à l'intuition ?
\item \textbf{Équation $P(t)=120$.} Résoudre dans $[0\,;\,24]$ l'équation $P(t)=120$. On pourra d'abord résoudre $\cos u=\dfrac{1}{2}$, puis revenir à la variable $t$. Donner les solutions en heures et minutes.
\item \textbf{Plage de fonctionnement de la pompe.} Résoudre dans $[0\,;\,24]$ l'inéquation $P(t)\geq 120$. On s'aidera du cercle trigonométrique pour déterminer les valeurs de $u$ telles que $\cos u\geq \dfrac{1}{2}$. En déduire la ou les plages horaires pendant lesquelles la pompe peut fonctionner, ainsi que la durée totale de fonctionnement.
\item \textbf{Vérification graphique.} À l'aide d'une calculatrice graphique ou d'un tracé soigné sur papier millimétré (quelques points calculés : $t=0$, $4$, $8$, $12$, $16$, $20$, $24$), représenter la courbe de $P$ sur $[0\,;\,24]$. Vérifier graphiquement les résultats des questions 2, 3 et 4. Commenter le signe de $P(t)$ la nuit : le modèle est-il réaliste ? Que faudrait-il préciser ?
\end{enumerate}

\courssub{Ressources à mobiliser}

\begin{aretenir}[Boîte à outils de la leçon]
\begin{itemize}
\item \textbf{Formules d'addition (rappel) :} $\cos(a-b)=\cos a\cos b+\sin a\sin b$ ; $\sin(a-b)=\sin a\cos b-\cos a\sin b$.
\item \textbf{Transformation :} pour $(a,b)\neq(0,0)$, il existe un réel $\varphi$ tel que
\[ a\cos X+b\sin X=\sqrt{a^{2}+b^{2}}\,\cos(X-\varphi),\quad \text{avec } \cos\varphi=\frac{a}{\sqrt{a^{2}+b^{2}}},\ \sin\varphi=\frac{b}{\sqrt{a^{2}+b^{2}}}. \]
\item \textbf{Équation $\cos u=\cos\alpha$ :} $u=\alpha+2k\pi$ ou $u=-\alpha+2k\pi$, $k\in\mathbb{Z}$.
\item \textbf{Inéquation $\cos u\geq c$ :} on lit les solutions sur le cercle trigonométrique (points dont l'abscisse est supérieure ou égale à $c$).
\item \textbf{Propriété admise :} une expression trigonométrique qui ne s'annule pas sur un intervalle y garde un signe constant.
\item \textbf{Valeurs remarquables :} $\cos\dfrac{\pi}{3}=\dfrac{1}{2}$, $\sin\dfrac{\pi}{3}=\dfrac{\sqrt{3}}{2}$, $\cos\dfrac{\pi}{6}=\dfrac{\sqrt{3}}{2}$, $\sin\dfrac{\pi}{6}=\dfrac{1}{2}$.
\end{itemize}
\end{aretenir}

\begin{methode}[Transformer $a\cos X+b\sin X$ en $A\cos(X-\varphi)$]
\begin{enumerate}
\item Calculer $A=\sqrt{a^{2}+b^{2}}$ (amplitude).
\item Factoriser par $A$ : $a\cos X+b\sin X = A\left(\frac{a}{A}\cos X+\frac{b}{A}\sin X\right)$.
\item Identifier $\varphi$ tel que $\cos\varphi=\frac{a}{A}$ et $\sin\varphi=\frac{b}{A}$ (déterminer $\varphi$ modulo $2\pi$ à l'aide du cercle trigonométrique ou des valeurs remarquables).
\item Écrire $a\cos X+b\sin X = A\cos(X-\varphi)$.
\end{enumerate}
\end{methode}

\begin{methode}[Résoudre $a\cos X+b\sin X = c$]
\begin{enumerate}
\item Transformer le membre de gauche en $A\cos(X-\varphi)$.
\item L'équation devient $\cos(X-\varphi) = \frac{c}{A}$.
\item Si $\left|\frac{c}{A}\right| > 1$, pas de solution. Sinon, poser $u=X-\varphi$ et résoudre $\cos u = \frac{c}{A}$ : $u = \pm\arccos(c/A) + 2k\pi$.
\item Revenir à $X$ : $X = \varphi \pm \arccos(c/A) + 2k\pi$.
\item Sélectionner les solutions dans l'intervalle demandé.
\end{enumerate}
\end{methode}

\begin{methode}[Résoudre $a\cos X+b\sin X \geq c$ sur un intervalle]
\begin{enumerate}
\item Transformer en $A\cos(X-\varphi) \geq c$ $\iff$ $\cos(X-\varphi) \geq \frac{c}{A}$.
\item Poser $u = X-\varphi$. Déterminer l'intervalle $J$ parcouru par $u$ quand $X$ varie dans l'intervalle donné.
\item Sur le cercle trigonométrique, repérer l'ensemble des points de $J$ dont l'abscisse est $\geq c/A$.
\item Traduire cet ensemble en inéquation(s) sur $u$, puis revenir à $X$.
\end{enumerate}
\end{methode}

\courssub{Démarche et corrigé commenté}

\begin{exoresolu}[Le panneau solaire de Maroua — corrigé complet]
On rappelle que $P(t)=120\cos\left(\dfrac{\pi t}{12}\right)+120\sqrt{3}\sin\left(\dfrac{\pi t}{12}\right)$ pour $t\in[0\,;\,24]$, et que la pompe exige $P(t)\geq 120$ W.
\tcblower
\textbf{Question 1 : écriture sous la forme $A\cos(X-\varphi)$.}

On pose $X=\dfrac{\pi t}{12}$, $a=120$ et $b=120\sqrt{3}$. On calcule l'amplitude :
\[ A=\sqrt{a^{2}+b^{2}}=\sqrt{120^{2}+\left(120\sqrt{3}\right)^{2}}=\sqrt{14400+43200}=\sqrt{57600}=240. \]
On factorise par $A$ :
\[ P(t)=240\left(\frac{120}{240}\cos X+\frac{120\sqrt{3}}{240}\sin X\right)=240\left(\frac{1}{2}\cos X+\frac{\sqrt{3}}{2}\sin X\right). \]
On cherche $\varphi$ tel que $\cos\varphi=\dfrac{1}{2}$ et $\sin\varphi=\dfrac{\sqrt{3}}{2}$. Sur le cercle trigonométrique, cela correspond à $\varphi=\dfrac{\pi}{3}$ (premier quadrant). D'après la formule $\cos(X-\varphi)=\cos X\cos\varphi+\sin X\sin\varphi$, on obtient :
\[ P(t)=240\cos\left(X-\frac{\pi}{3}\right)=240\cos\left(\frac{\pi t}{12}-\frac{\pi}{3}\right). \]
On a donc $A=240$ et $\varphi=\dfrac{\pi}{3}$.

\textbf{Question 2 : puissance maximale et heure correspondante.}

Pour tout réel $u$, $-1\leq\cos u\leq 1$, donc $P(t)\leq 240$. La puissance maximale est $P_{\max}=240$ W. Elle est atteinte lorsque $\cos\left(\frac{\pi t}{12}-\frac{\pi}{3}\right)=1$, c'est-à-dire :
\[ \frac{\pi t}{12}-\frac{\pi}{3}=2k\pi \iff \frac{\pi t}{12}=\frac{\pi}{3}+2k\pi \iff t=4+24k,\quad k\in\mathbb{Z}. \]
Dans $[0\,;\,24]$, seule $t_0=4$ convient. Cela correspond à \textbf{4 h du matin}. Ce résultat contredit l'intuition physique (maximum attendu vers midi) : le modèle mathématique est correct, mais son déphasage le rend physiquement inadapté sans correction. \emph{Pour la suite, on poursuit l'étude du modèle tel qu'il est proposé.}

\textbf{Question 3 : résolution de $P(t)=120$.}
\begin{align*}
P(t)=120 &\iff 240\cos\left(\frac{\pi t}{12}-\frac{\pi}{3}\right)=120 \iff \cos\left(\frac{\pi t}{12}-\frac{\pi}{3}\right)=\frac{1}{2}.
\end{align*}
On pose $u=\dfrac{\pi t}{12}-\dfrac{\pi}{3}$. L'équation $\cos u=\dfrac{1}{2}=\cos\dfrac{\pi}{3}$ admet pour solutions :
\[ u=\frac{\pi}{3}+2k\pi \quad\text{ou}\quad u=-\frac{\pi}{3}+2k\pi,\quad k\in\mathbb{Z}. \]

\emph{Premier cas :} $\dfrac{\pi t}{12}-\dfrac{\pi}{3}=\dfrac{\pi}{3}+2k\pi \iff \dfrac{\pi t}{12}=\dfrac{2\pi}{3}+2k\pi \iff t=8+24k$.

\emph{Deuxième cas :} $\dfrac{\pi t}{12}-\dfrac{\pi}{3}=-\dfrac{\pi}{3}+2k\pi \iff \dfrac{\pi t}{12}=2k\pi \iff t=24k$.

Dans l'intervalle $[0\,;\,24]$, les solutions sont $t=0$, $t=8$ et $t=24$. La puissance vaut exactement $120$ W à \textbf{minuit} ($t=0$ et $t=24$, même instant) et à \textbf{8 h du matin}.

\textbf{Question 4 : résolution de $P(t)\geq 120$.}

On pose $u=\dfrac{\pi t}{12}-\dfrac{\pi}{3}$. Quand $t$ décrit $[0\,;\,24]$, $u$ décrit l'intervalle $J=\left[-\dfrac{\pi}{3}\,;\,2\pi-\dfrac{\pi}{3}\right]=\left[-\dfrac{\pi}{3}\,;\,\dfrac{5\pi}{3}\right]$.
On cherche $u\in J$ tels que $\cos u\geq \dfrac{1}{2}$.
Sur le cercle trigonométrique, $\cos u = \frac{1}{2}$ pour $u=-\frac{\pi}{3}$, $u=\frac{\pi}{3}$ et $u=\frac{5\pi}{3}$ (qui est $-\frac{\pi}{3}+2\pi$). L'abscisse est $\geq \frac{1}{2}$ sur l'arc allant de $-\frac{\pi}{3}$ à $\frac{\pi}{3}$, et en le point $\frac{5\pi}{3}$.
Donc $u\in\left[-\dfrac{\pi}{3}\,;\,\dfrac{\pi}{3}\right] \cup \left\{\dfrac{5\pi}{3}\right\}$.

On revient à $t$ :
\begin{itemize}
\item $-\dfrac{\pi}{3}\leq\dfrac{\pi t}{12}-\dfrac{\pi}{3}\leq\dfrac{\pi}{3} \iff 0\leq\dfrac{\pi t}{12}\leq\dfrac{2\pi}{3} \iff 0\leq t\leq 8$.
\item $\dfrac{\pi t}{12}-\dfrac{\pi}{3}=\dfrac{5\pi}{3} \iff \dfrac{\pi t}{12}=2\pi \iff t=24$.
\end{itemize}
L'ensemble des solutions dans $[0\,;\,24]$ est $\mathcal{S}=[0\,;\,8] \cup \{24\}$. Comme $t=0$ et $t=24$ représentent le même instant (minuit), la pompe peut fonctionner de \textbf{minuit à 8 h du matin}, soit une durée totale de \textbf{8 heures}. (Le décalage du modèle persiste : physiquement, on attendrait une plage centrée sur midi.)

\textbf{Question 5 : vérification graphique.}

On calcule quelques valeurs pour tracer la courbe représentative de $P$ sur $[0\,;\,24]$ :
\begin{center}
\begin{tabular}{|c|c|c|c|c|c|c|c|}\hline
$t$ & 0 & 4 & 8 & 12 & 16 & 20 & 24 \\ \hline
$P(t)$ (W) & 120 & 240 & 120 & -120 & -240 & -120 & 120 \\ \hline
\end{tabular}
\end{center}

\begin{popfigure}
\begin{tikzpicture}
\begin{axis}[
    width=\linewidth, height=8cm,
    axis lines=middle,
    xlabel=$t$ (h), ylabel=$P(t)$ (W),
    xmin=-0.5, xmax=24.5, ymin=-260, ymax=260,
    xtick={0,4,8,12,16,20,24},
    ytick={-240,-120,0,120,240},
    grid=major, grid style={popmuted},
    tick label style={font=\small},
    label style={font=\small},
    domain=0:24, samples=200,
    clip=false
]
\addplot[popcurve, thick, name path=P] {240*cos(15*x - 60)};
\addplot[poporange, dashed, thick, name path=S] coordinates {(0,120) (24,120)} node[pos=0.95, above right, font=\small, poporange] {$y=120$};
\addplot[popgreen, dashed] coordinates {(0,0) (24,0)};
% Points clés
\addplot[only marks, mark=*, popblue, mark size=2.5] coordinates {(0,120) (4,240) (8,120) (12,-120) (16,-240) (20,-120) (24,120)};
% Annotations
\node[popblue, anchor=south west, font=\tiny] at (axis cs:4,240) {Max 240 W};
\node[poporange, anchor=north west, font=\tiny] at (axis cs:0,120) {Seuil 120 W};
\node[poporange, anchor=north west, font=\tiny] at (axis cs:8,120) {Seuil 120 W};
\node[poporange, anchor=north east, font=\tiny] at (axis cs:24,120) {Seuil 120 W};
% Zone de fonctionnement
\addplot[popgreenL!50, opacity=0.3] fill between[of=P and S, soft clip={domain=0:8}];
\end{axis}
\end{tikzpicture}
\legende{Courbe de la puissance $P(t)$ et zone de fonctionnement de la pompe ($P(t)\geq 120$ W en vert).}
\end{popfigure}

La courbe confirme :
\begin{itemize}
\item Maximum de $240$ W en $t=4$ (4 h).
\item Passage par $120$ W en $t=0$, $t=8$ et $t=24$.
\item $P(t)\geq 120$ exactement sur $[0\,;\,8]$ (zone verte).
\item $P(t)<0$ pour $t\in]10\,;\,22[$ environ : une puissance négative n'a pas de sens physique. Le modèle n'est qu'une approximation mathématique ; il faudrait le recalibrer (par exemple en remplaçant $\frac{\pi t}{12}$ par $\frac{\pi(t-12)}{12}$ pour centrer le maximum sur midi).
\end{itemize}
\end{exoresolu}

\begin{attention}[Pièges rencontrés dans l'activité]
\begin{itemize}
\item Ne pas oublier de factoriser par $A=\sqrt{a^{2}+b^{2}}$ \emph{avant} d'identifier $\cos\varphi$ et $\sin\varphi$ : les nombres $\dfrac{a}{A}$ et $\dfrac{b}{A}$ doivent être compris entre $-1$ et $1$.
\item Dans le retour à la variable $t$, diviser par le coefficient $\dfrac{\pi}{12}$ revient à \emph{multiplier par $\dfrac{12}{\pi}$} : les erreurs de calcul sont fréquentes à cette étape.
\item Pour l'inéquation, les solutions se lisent sur le cercle trigonométrique \emph{dans l'intervalle parcouru par $u$} (ici $[-\pi/3\,;\,5\pi/3]$), qui n'est pas forcément $[0\,;\,2\pi]$. Ne pas oublier les solutions aux bornes de cet intervalle (ici $u=5\pi/3 \leftrightarrow t=24$).
\item Vérifier toujours la condition $\left|\frac{c}{A}\right|\leq 1$ avant de résoudre $\cos u = \frac{c}{A}$.
\end{itemize}
\end{attention}

\courssub{Bilan de l'activité}

Cette activité a permis de mettre en évidence que :
\begin{itemize}
\item l'écriture $a\cos X+b\sin X=A\cos(X-\varphi)$ est un outil puissant : elle ramène l'étude d'une somme de deux fonctions trigonométriques à celle d'un \emph{seul} cosinus, dont on connaît le maximum ($A$), le minimum ($-A$) et les variations ;
\item les équations et inéquations du type $a\cos X+b\sin X=c$ se résolvent en deux temps : transformation, puis résolution élémentaire $\cos u=\alpha$ ou $\cos u\geq\alpha$ à l'aide du cercle trigonométrique, avant le retour à la variable initiale ;
\item les solutions d'une inéquation trigonométrique sur un intervalle borné s'interprètent concrètement (ici, une plage horaire et une durée de fonctionnement) ;
\item un modèle mathématique doit toujours être confronté à la réalité : les calculs peuvent être exacts alors que le modèle est imparfait, et la vérification graphique est un réflexe indispensable.
\end{itemize}

\begin{savaistu}[Le soleil, premier fournisseur d'énergie du Sahel]
À Maroua, l'ensoleillement dépasse $3000$ heures par an, parmi les plus élevés du Cameroun. Les ingénieurs qui dimensionnent les installations solaires utilisent quotidiennement des fonctions trigonométriques : inclinaison optimale des panneaux, hauteur du soleil selon la saison, puissance instantanée... Les formules que tu viens d'utiliser sont exactement celles de leurs cahiers de calcul !
\end{savaistu}

\courssub{Exercices d'entraînement}

\begin{exercice}[Transformation et valeurs exactes]
\begin{enumerate}
\item Transformer les expressions suivantes sous la forme $A\cos(x-\varphi)$ avec $A>0$ et $\varphi\in[0\,;\,2\pi[$ :
\begin{itemize}
\item $E_1(x)=\cos x-\sqrt{3}\sin x$
\item $E_2(x)=2\sin x-2\cos x$
\item $E_3(x)=-\sqrt{3}\cos x-\sin x$
\end{itemize}
\item En déduire les valeurs exactes de $\cos\dfrac{\pi}{12}$ et $\sin\dfrac{\pi}{12}$ en utilisant $E_1\left(\dfrac{\pi}{12}\right)$ et la valeur connue de $\cos\dfrac{\pi}{4}$.
\end{enumerate}
\end{exercice}

\begin{exercice}[Équation $a\cos x+b\sin x=c$]
Résoudre dans $[0\,;\,2\pi]$ les équations suivantes :
\begin{enumerate}
\item $\sqrt{3}\cos x+\sin x = \sqrt{2}$
\item $2\cos x-2\sin x = 1$
\item $\cos x+\sin x = -1$
\end{enumerate}
\end{exercice}

\begin{exercice}[Inéquation trigonométrique]
Résoudre dans $[0\,;\,2\pi]$ les inéquations :
\begin{enumerate}
\item $\cos x-\sqrt{3}\sin x \geq 0$
\item $2\sin x+2\cos x < \sqrt{2}$
\end{enumerate}
\end{exercice}

\begin{exercice}[Modélisation : Marée à Douala]
Le niveau de l'eau (en mètres) dans le port de Douala est modélisé par la fonction $h(t)=2\cos\left(\frac{\pi t}{6}\right)+2\sqrt{3}\sin\left(\frac{\pi t}{6}\right)$ où $t$ est le temps en heures après minuit.
\begin{enumerate}
\item Écrire $h(t)$ sous la forme $A\cos(\omega t-\varphi)$.
\item Déterminer les heures des marées hautes (maxima) et marées basses (minima) dans la journée $t\in[0\,;\,24]$.
\item Un bateau a besoin de $3,5$ m d'eau pour entrer au port. Pendant quelles plages horaires peut-il entrer ?
\end{enumerate}
\end{exercice}

\begin{corrige}[Exercices d'entraînement]
\textbf{Exercice 1.}
\begin{enumerate}
\item $E_1(x)=2\cos\left(x+\frac{\pi}{3}\right)$ ; $E_2(x)=2\sqrt{2}\cos\left(x-\frac{3\pi}{4}\right)$ ; $E_3(x)=2\cos\left(x-\frac{5\pi}{6}\right)$.
\item $E_1(\pi/12)=2\cos(\pi/12+\pi/3)=2\cos(5\pi/12)=\cos(\pi/4)-\sqrt{3}\sin(\pi/4)=\sqrt{2}/2-\sqrt{6}/2$. D'où $\cos(5\pi/12)=(\sqrt{2}-\sqrt{6})/4$. En utilisant $\cos(\pi/12)=\sin(5\pi/12)$, on trouve $\cos(\pi/12)=(\sqrt{6}+\sqrt{2})/4$, $\sin(\pi/12)=(\sqrt{6}-\sqrt{2})/4$.
\end{enumerate}
\textbf{Exercice 2.}
1. $2\cos(x-\pi/6)=\sqrt{2} \iff \cos(x-\pi/6)=\sqrt{2}/2 \iff x-\pi/6=\pm\pi/4+2k\pi \iff x=\pi/6\pm\pi/4+2k\pi$. Dans $[0,2\pi]$ : $x=\pi/12,\ 5\pi/12,\ 13\pi/12,\ 17\pi/12$ (vérifier l'appartenance).
2. $2\sqrt{2}\cos(x+\pi/4)=1 \iff \cos(x+\pi/4)=\sqrt{2}/4$. Solutions : $x+\pi/4=\pm\arccos(\sqrt{2}/4)+2k\pi$.
3. $\sqrt{2}\cos(x-\pi/4)=-1 \iff \cos(x-\pi/4)=-\sqrt{2}/2 \iff x-\pi/4=\pm 3\pi/4+2k\pi \iff x=\pi/4\pm 3\pi/4+2k\pi$. Dans $[0,2\pi]$ : $x=\pi,\ 3\pi/2$.
\textbf{Exercice 3.}
1. $2\cos(x+\pi/3)\geq 0 \iff \cos(x+\pi/3)\geq 0$. $u=x+\pi/3\in[\pi/3,7\pi/3]$. $\cos u\geq 0 \iff u\in[\pi/3,\pi/2]\cup[3\pi/2,5\pi/2]\cup\{7\pi/3\}$. Donc $x\in[0,\pi/6]\cup[7\pi/6,13\pi/6]\cap[0,2\pi]=[0,\pi/6]\cup[7\pi/6,2\pi]$.
2. $2\sqrt{2}\sin(x+\pi/4)<\sqrt{2} \iff \sin(x+\pi/4)<1/2$. Résolution similaire sur le cercle.
\textbf{Exercice 4.}
1. $h(t)=4\cos(\pi t/6-\pi/3)$. Amplitude 4 m.
2. Max ($h=4$) : $\pi t/6-\pi/3=2k\pi \Rightarrow t=2+12k \Rightarrow$ 2 h et 14 h. Min ($h=-4$) : $\pi t/6-\pi/3=\pi+2k\pi \Rightarrow t=8+12k \Rightarrow$ 8 h et 20 h.
3. $h(t)\geq 3.5 \iff 4\cos(\pi t/6-\pi/3)\geq 3.5 \iff \cos u\geq 0.875$. $u=\pi t/6-\pi/3$. $\arccos(0.875)\approx 0.505$ rad. $u\in[-0.505+2k\pi, 0.505+2k\pi]$. Pour $k=0$ : $t\in[0.03, 1.97]$ (env. 0 h 02 à 1 h 58). Pour $k=1$ : $t\in[12.03, 13.97]$ (env. 12 h 02 à 13 h 58).
\end{corrige}

\newpage

\coursec{Méthodes et exercices résolus}

\courssub{Fiche 1 : Utiliser les formules d'addition et de duplication}

\begin{aretenir}[Les formules à connaître par c\oe ur]
\textbf{Formules d'addition :}
\[ \cos(a+b)=\cos a\cos b-\sin a\sin b \qquad \cos(a-b)=\cos a\cos b+\sin a\sin b \]
\[ \sin(a+b)=\sin a\cos b+\cos a\sin b \qquad \sin(a-b)=\sin a\cos b-\cos a\sin b \]
\[ \tan(a+b)=\dfrac{\tan a+\tan b}{1-\tan a\tan b} \qquad \tan(a-b)=\dfrac{\tan a-\tan b}{1+\tan a\tan b} \]
(sous réserve que toutes les tangentes soient définies et que les dénominateurs soient non nuls).

\textbf{Formules de duplication} (obtenues en posant $b=a$) :
\[ \cos 2a=\cos^2 a-\sin^2 a=2\cos^2 a-1=1-2\sin^2 a \qquad \sin 2a=2\sin a\cos a \qquad \tan 2a=\dfrac{2\tan a}{1-\tan^2 a}. \]

\textbf{Formules de linéarisation} (on isole $\cos^2 a$ et $\sin^2 a$ dans la duplication) :
\[ \cos^2 a=\dfrac{1+\cos 2a}{2} \qquad \sin^2 a=\dfrac{1-\cos 2a}{2}. \]
\end{aretenir}

\begin{methode}[Calculer une valeur exacte grâce aux formules d'addition]
Pour calculer $\cos\theta$, $\sin\theta$ ou $\tan\theta$ lorsque $\theta$ n'est pas une valeur remarquable :
\begin{enumerate}
\item \textbf{Décomposer} $\theta$ en somme ou différence d'angles remarquables ($\dfrac{\pi}{6}$, $\dfrac{\pi}{4}$, $\dfrac{\pi}{3}$, $\dfrac{\pi}{2}$). Exemples utiles : $\dfrac{\pi}{12}=\dfrac{\pi}{3}-\dfrac{\pi}{4}$ ; $\dfrac{7\pi}{12}=\dfrac{\pi}{3}+\dfrac{\pi}{4}$ ; $\dfrac{5\pi}{12}=\dfrac{\pi}{4}+\dfrac{\pi}{6}$.
\item \textbf{Choisir la bonne formule} d'addition ou de duplication.
\item \textbf{Substituer} les valeurs remarquables connues, puis simplifier (factoriser, réduire au même dénominateur, rendre rationnel le dénominateur).
\item Pour la tangente, utiliser $\tan(a+b)=\dfrac{\tan a+\tan b}{1-\tan a\tan b}$, \textbf{à condition} que les tangentes existent et que $1-\tan a\tan b\neq 0$.
\item \textbf{Contrôler le signe} du résultat à l'aide du quadrant auquel appartient $\theta$.
\end{enumerate}
\end{methode}

\begin{exoresolu}[Valeurs exactes de $\cos\dfrac{\pi}{12}$ et $\tan\dfrac{7\pi}{12}$]
\begin{enumerate}
\item Calculer la valeur exacte de $\cos\dfrac{\pi}{12}$.
\item En remarquant que $\dfrac{7\pi}{12}=\dfrac{\pi}{3}+\dfrac{\pi}{4}$, calculer la valeur exacte de $\tan\dfrac{7\pi}{12}$.
\end{enumerate}
\tcblower
\textbf{1.} On écrit $\dfrac{\pi}{12}=\dfrac{\pi}{3}-\dfrac{\pi}{4}$ (vérification : $\dfrac{\pi}{3}-\dfrac{\pi}{4}=\dfrac{4\pi-3\pi}{12}=\dfrac{\pi}{12}$). On applique la formule d'addition du cosinus :
\begin{align*}
\cos\dfrac{\pi}{12} &= \cos\left(\dfrac{\pi}{3}-\dfrac{\pi}{4}\right)=\cos\dfrac{\pi}{3}\cos\dfrac{\pi}{4}+\sin\dfrac{\pi}{3}\sin\dfrac{\pi}{4}\\
&= \dfrac{1}{2}\times\dfrac{\sqrt{2}}{2}+\dfrac{\sqrt{3}}{2}\times\dfrac{\sqrt{2}}{2}=\dfrac{\sqrt{2}+\sqrt{6}}{4}.
\end{align*}
\textbf{Conclusion :} $\cos\dfrac{\pi}{12}=\dfrac{\sqrt{6}+\sqrt{2}}{4}$. (Cohérent : $\dfrac{\pi}{12}\in\left]0;\dfrac{\pi}{2}\right[$, donc son cosinus est positif.)

\textbf{2.} On a $\tan\dfrac{\pi}{3}=\sqrt{3}$ et $\tan\dfrac{\pi}{4}=1$. Comme $1-\tan\dfrac{\pi}{3}\tan\dfrac{\pi}{4}=1-\sqrt{3}\neq 0$, la formule s'applique :
\[ \tan\dfrac{7\pi}{12}=\tan\left(\dfrac{\pi}{3}+\dfrac{\pi}{4}\right)=\dfrac{\sqrt{3}+1}{1-\sqrt{3}}=\dfrac{(\sqrt{3}+1)(1+\sqrt{3})}{(1-\sqrt{3})(1+\sqrt{3})}=\dfrac{4+2\sqrt{3}}{1-3}=\dfrac{4+2\sqrt{3}}{-2}=-2-\sqrt{3}. \]
\textbf{Conclusion :} $\tan\dfrac{7\pi}{12}=-2-\sqrt{3}$. (Cohérent : $\dfrac{7\pi}{12}\in\left]\dfrac{\pi}{2};\pi\right[$, donc sa tangente est négative.)
\end{exoresolu}

\begin{exemplebox}[Linéariser pour simplifier un calcul]
Calculer $\sin^2\dfrac{\pi}{8}$. D'après la formule de linéarisation avec $a=\dfrac{\pi}{8}$ (donc $2a=\dfrac{\pi}{4}$) :
\[ \sin^2\dfrac{\pi}{8}=\dfrac{1-\cos\dfrac{\pi}{4}}{2}=\dfrac{1-\dfrac{\sqrt{2}}{2}}{2}=\dfrac{2-\sqrt{2}}{4}. \]
Comme $\dfrac{\pi}{8}\in\left]0;\dfrac{\pi}{2}\right[$, $\sin\dfrac{\pi}{8}>0$, donc $\sin\dfrac{\pi}{8}=\dfrac{\sqrt{2-\sqrt{2}}}{2}$.
\end{exemplebox}

\courssub{Fiche 2 : Écrire $a\cos X+b\sin X$ sous la forme $A\cos(X+\varphi)$}

\begin{propriete}[Transformation de $a\cos X+b\sin X$ — admise]
Soit $a$ et $b$ deux réels non tous les deux nuls. En posant $A=\sqrt{a^2+b^2}>0$, il existe un réel $\varphi$ tel que, pour tout réel $X$ :
\[ a\cos X+b\sin X=A\cos(X+\varphi). \]
On peut aussi écrire $a\cos X+b\sin X=A\sin\left(\dfrac{\pi}{2}-X-\varphi\right)$, puisque $\cos t=\sin\left(\dfrac{\pi}{2}-t\right)$ pour tout réel $t$.
\end{propriete}

\begin{experience}[Pourquoi cette transformation est-elle possible ?]
Factorisons par $A=\sqrt{a^2+b^2}$ :
\[ a\cos X+b\sin X=A\left(\dfrac{a}{A}\cos X+\dfrac{b}{A}\sin X\right). \]
Le point $M$ de coordonnées $\left(\dfrac{a}{A};\dfrac{b}{A}\right)$ vérifie $\left(\dfrac{a}{A}\right)^2+\left(\dfrac{b}{A}\right)^2=\dfrac{a^2+b^2}{A^2}=1$ : il est donc \textbf{sur le cercle trigonométrique}. Il existe alors un réel $\varphi_0$ tel que $\cos\varphi_0=\dfrac{a}{A}$ et $\sin\varphi_0=\dfrac{b}{A}$. On reconnaît la formule de $\cos(X-\varphi_0)$ :
\[ a\cos X+b\sin X=A\left(\cos X\cos\varphi_0+\sin X\sin\varphi_0\right)=A\cos(X-\varphi_0). \]
Toute la difficulté pratique consiste donc à \textbf{identifier $\varphi_0$} à l'aide des valeurs remarquables.
\end{experience}

\begin{methode}[Transformation de $a\cos X+b\sin X$ en pratique]
Soit $a$ et $b$ deux réels non tous les deux nuls.
\begin{enumerate}
\item Calculer $A=\sqrt{a^2+b^2}>0$.
\item Factoriser : $a\cos X+b\sin X=A\left(\dfrac{a}{A}\cos X+\dfrac{b}{A}\sin X\right)$.
\item Déterminer $\varphi_0$ tel que $\cos\varphi_0=\dfrac{a}{A}$ et $\sin\varphi_0=\dfrac{b}{A}$, en plaçant le point $\left(\dfrac{a}{A};\dfrac{b}{A}\right)$ sur le cercle trigonométrique.
\item Conclure : $a\cos X+b\sin X=A\cos(X-\varphi_0)$.
\item \textbf{Vérifier} en développant $A\cos(X-\varphi_0)$ avec la formule d'addition : on doit retrouver $a\cos X+b\sin X$.
\end{enumerate}
\textbf{Remarque :} pour obtenir la forme $A\cos(X+\varphi)$, il suffit de prendre $\varphi=-\varphi_0$, c'est-à-dire $\cos\varphi=\dfrac{a}{A}$ et $\sin\varphi=-\dfrac{b}{A}$. Les deux écritures sont acceptées au Probatoire ; l'essentiel est la vérification finale.
\end{methode}

\begin{exoresolu}[Transformer $\cos X+\sqrt{3}\sin X$]
Soit $f(X)=\cos X+\sqrt{3}\sin X$. Écrire $f(X)$ sous la forme $A\cos(X-\varphi)$ avec $A>0$ et $\varphi\in\left]-\pi;\pi\right]$.
\tcblower
\textbf{Étape 1 :} $a=1$, $b=\sqrt{3}$, donc $A=\sqrt{1^2+(\sqrt{3})^2}=\sqrt{4}=2$.

\textbf{Étape 2 :} $f(X)=2\left(\dfrac{1}{2}\cos X+\dfrac{\sqrt{3}}{2}\sin X\right)$.

\textbf{Étape 3 :} On cherche $\varphi$ tel que $\cos\varphi=\dfrac{1}{2}$ et $\sin\varphi=\dfrac{\sqrt{3}}{2}$ : c'est $\varphi=\dfrac{\pi}{3}$.

\textbf{Étape 4 :} $\cos X\cos\dfrac{\pi}{3}+\sin X\sin\dfrac{\pi}{3}=\cos\left(X-\dfrac{\pi}{3}\right)$.

\textbf{Vérification :} $2\cos\left(X-\dfrac{\pi}{3}\right)=2\left(\cos X\times\dfrac{1}{2}+\sin X\times\dfrac{\sqrt{3}}{2}\right)=\cos X+\sqrt{3}\sin X$. \checkmark

\textbf{Conclusion :} $\cos X+\sqrt{3}\sin X=2\cos\left(X-\dfrac{\pi}{3}\right)$.
\end{exoresolu}

\begin{exemplebox}[Un deuxième exemple pour s'entraîner l'œil]
Écrivons $\sin X-\cos X$ sous la forme $A\cos(X-\varphi)$. Ici $a=-1$, $b=1$, donc $A=\sqrt{2}$ et
\[ \sin X-\cos X=\sqrt{2}\left(-\dfrac{1}{\sqrt{2}}\cos X+\dfrac{1}{\sqrt{2}}\sin X\right)=\sqrt{2}\left(\cos X\cos\dfrac{3\pi}{4}+\sin X\sin\dfrac{3\pi}{4}\right)=\sqrt{2}\cos\left(X-\dfrac{3\pi}{4}\right). \]
Vérification : $\sqrt{2}\left(\cos X\times\left(-\dfrac{\sqrt{2}}{2}\right)+\sin X\times\dfrac{\sqrt{2}}{2}\right)=-\cos X+\sin X$. \checkmark
\end{exemplebox}

\courssub{Fiche 3 : Résoudre $\cos X=a$, $\sin X=a$, $\tan X=a$}

\begin{propriete}[Solutions des équations élémentaires — à savoir rédiger]
Pour tout réel $\alpha$ et tout entier $k\in\mathbb{Z}$ :
\begin{itemize}
\item $\cos X=\cos\alpha \iff X=\alpha+2k\pi$ \textbf{ou} $X=-\alpha+2k\pi$ ;
\item $\sin X=\sin\alpha \iff X=\alpha+2k\pi$ \textbf{ou} $X=\pi-\alpha+2k\pi$ ;
\item $\tan X=\tan\alpha \iff X=\alpha+k\pi$ (avec $X\neq\dfrac{\pi}{2}+k\pi$).
\end{itemize}
\textbf{Existence :} $\cos X=a$ et $\sin X=a$ n'ont de solutions que si $a\in[-1;1]$ ; $\tan X=a$ a des solutions pour tout réel $a$.
\end{propriete}

\begin{methode}[Résoudre une équation trigonométrique élémentaire]
\begin{enumerate}
\item \textbf{Vérifier la condition d'existence} : $a\in[-1;1]$ pour le cosinus et le sinus.
\item \textbf{Trouver une solution particulière} $\alpha$ (valeur remarquable ou valeur donnée par l'énoncé).
\item \textbf{Écrire la solution générale} avec les deux familles (une seule pour la tangente), en précisant $k\in\mathbb{Z}$.
\item \textbf{Sélectionner} les solutions appartenant à l'intervalle demandé en testant des valeurs entières de $k$ (souvent $k=-1,0,1$).
\end{enumerate}
\textbf{Réflexe visuel :} sur le cercle trigonométrique, $\cos X=a$ se lit par une droite \emph{verticale} d'abscisse $a$, $\sin X=a$ par une droite \emph{horizontale} d'ordonnée $a$ : chacune coupe le cercle en deux points (deux solutions par période), tandis que $\tan X=a$ donne un seul point par intervalle de longueur $\pi$.
\end{methode}

\begin{exoresolu}[Résoudre $\cos X=-\dfrac{1}{2}$ puis $\sin X=\dfrac{\sqrt{2}}{2}$]
\begin{enumerate}
\item Résoudre dans $\mathbb{R}$ l'équation $\cos X=-\dfrac{1}{2}$, puis donner les solutions dans $[0;2\pi]$.
\item Résoudre dans $]-\pi;\pi]$ l'équation $\sin X=\dfrac{\sqrt{2}}{2}$.
\end{enumerate}
\tcblower
\textbf{1.} $-\dfrac{1}{2}\in[-1;1]$ : des solutions existent. On sait que $\cos\dfrac{2\pi}{3}=-\dfrac{1}{2}$. Donc :
\[ \cos X=\cos\dfrac{2\pi}{3}\iff X=\dfrac{2\pi}{3}+2k\pi \ \text{ou}\ X=-\dfrac{2\pi}{3}+2k\pi,\quad k\in\mathbb{Z}. \]
Dans $[0;2\pi]$ : pour $k=0$, $X=\dfrac{2\pi}{3}$ ; pour $k=1$, $X=-\dfrac{2\pi}{3}+2\pi=\dfrac{4\pi}{3}$.
\textbf{Conclusion :} $\mathcal{S}_{[0;2\pi]}=\left\{\dfrac{2\pi}{3};\dfrac{4\pi}{3}\right\}$.

\textbf{2.} $\sin\dfrac{\pi}{4}=\dfrac{\sqrt{2}}{2}$, donc $X=\dfrac{\pi}{4}+2k\pi$ ou $X=\pi-\dfrac{\pi}{4}+2k\pi=\dfrac{3\pi}{4}+2k\pi$, $k\in\mathbb{Z}$.
Dans $]-\pi;\pi]$, seules les valeurs avec $k=0$ conviennent : $\dfrac{\pi}{4}$ et $\dfrac{3\pi}{4}$.
\textbf{Conclusion :} $\mathcal{S}=\left\{\dfrac{\pi}{4};\dfrac{3\pi}{4}\right\}$.
\end{exoresolu}

\begin{exoresolu}[{Résoudre $\tan X=\sqrt{3}$ dans $[0;2\pi]$}]
Résoudre dans $[0;2\pi]$ l'équation $\tan X=\sqrt{3}$.
\tcblower
La tangente est définie pour $X\neq\dfrac{\pi}{2}+k\pi$. On sait que $\tan\dfrac{\pi}{3}=\sqrt{3}$. Donc :
\[ \tan X=\tan\dfrac{\pi}{3}\iff X=\dfrac{\pi}{3}+k\pi,\quad k\in\mathbb{Z}. \]
Dans $[0;2\pi]$ : pour $k=0$, $X=\dfrac{\pi}{3}$ ; pour $k=1$, $X=\dfrac{\pi}{3}+\pi=\dfrac{4\pi}{3}$. (Pour $k=2$, $X=\dfrac{7\pi}{3}>2\pi$ : rejeté.)

\textbf{Vérification :} $\tan\dfrac{4\pi}{3}=\dfrac{\sin\frac{4\pi}{3}}{\cos\frac{4\pi}{3}}=\dfrac{-\frac{\sqrt{3}}{2}}{-\frac{1}{2}}=\sqrt{3}$. \checkmark

\textbf{Conclusion :} $\mathcal{S}=\left\{\dfrac{\pi}{3};\dfrac{4\pi}{3}\right\}$.
\end{exoresolu}

\courssub{Fiche 4 : Résoudre $a\cos X+b\sin X=c$}

\begin{methode}[Équation $a\cos X+b\sin X=c$]
\begin{enumerate}
\item \textbf{Transformer} le premier membre : écrire $a\cos X+b\sin X=A\cos(X-\varphi)$ avec $A=\sqrt{a^2+b^2}$ (Fiche 2).
\item L'équation devient $\cos(X-\varphi)=\dfrac{c}{A}$. \textbf{Condition d'existence :} $\left|\dfrac{c}{A}\right|\leq 1$, c'est-à-dire $|c|\leq\sqrt{a^2+b^2}$. Si elle n'est pas vérifiée, conclure $\mathcal{S}=\varnothing$ en le justifiant.
\item Résoudre l'équation élémentaire en $X-\varphi$ : $X-\varphi=\pm\alpha+2k\pi$ où $\alpha$ vérifie $\cos\alpha=\dfrac{c}{A}$.
\item \textbf{Isoler $X$} (ne pas oublier d'ajouter $\varphi$ !) puis sélectionner les solutions dans l'intervalle demandé.
\end{enumerate}
\end{methode}

\begin{exoresolu}[{Résoudre $\cos X+\sqrt{3}\sin X=1$ dans $[0;2\pi]$}]
Résoudre dans $[0;2\pi]$ l'équation $(E)$ : $\cos X+\sqrt{3}\sin X=1$.
\tcblower
\textbf{Étape 1 :} D'après la Fiche 2, $\cos X+\sqrt{3}\sin X=2\cos\left(X-\dfrac{\pi}{3}\right)$.

\textbf{Étape 2 :} $(E)\iff 2\cos\left(X-\dfrac{\pi}{3}\right)=1\iff\cos\left(X-\dfrac{\pi}{3}\right)=\dfrac{1}{2}$. Comme $\dfrac{1}{2}\in[-1;1]$ (et $|1|\leq 2=\sqrt{1^2+(\sqrt{3})^2}$), il y a des solutions.

\textbf{Étape 3 :} $\cos\left(X-\dfrac{\pi}{3}\right)=\cos\dfrac{\pi}{3}$ donne :
\[ X-\dfrac{\pi}{3}=\dfrac{\pi}{3}+2k\pi \quad\text{ou}\quad X-\dfrac{\pi}{3}=-\dfrac{\pi}{3}+2k\pi,\quad k\in\mathbb{Z}. \]

\textbf{Étape 4 :} $X=\dfrac{2\pi}{3}+2k\pi$ ou $X=2k\pi$.

Dans $[0;2\pi]$ : pour la première famille, $k=0$ donne $X=\dfrac{2\pi}{3}$ ; pour la seconde, $k=0$ donne $X=0$ et $k=1$ donne $X=2\pi$.

\textbf{Vérification :} pour $X=\dfrac{2\pi}{3}$ : $\cos\dfrac{2\pi}{3}+\sqrt{3}\sin\dfrac{2\pi}{3}=-\dfrac{1}{2}+\sqrt{3}\times\dfrac{\sqrt{3}}{2}=-\dfrac{1}{2}+\dfrac{3}{2}=1$. \checkmark

\textbf{Conclusion :} $\mathcal{S}=\left\{0;\dfrac{2\pi}{3};2\pi\right\}$.
\end{exoresolu}

\courssub{Fiche 5 : Résoudre une inéquation trigonométrique de base}

\begin{propriete}[Signe constant — propriété admise au programme]
Si une expression trigonométrique ne s'annule pas sur un intervalle, alors elle garde un \textbf{signe constant} sur cet intervalle. Conséquence pratique : une fois connues les solutions de l'équation associée, il suffit de \textbf{tester une valeur} dans chacun des intervalles qu'elles délimitent pour connaître le signe de l'expression partout.
\end{propriete}

\begin{methode}[Inéquations $\cos X\geq a$, $\sin X<a$, $\tan X\leq a$]
\begin{enumerate}
\item \textbf{Résoudre d'abord l'équation} associée ($\cos X=a$, etc.) pour obtenir les valeurs frontières.
\item \textbf{Placer ces valeurs sur le cercle trigonométrique} et hachurer l'arc solution :
\begin{itemize}
\item $\cos X\geq a$ : points dont l'\emph{abscisse} est $\geq a$ (arc à droite de la droite verticale d'abscisse $a$) ;
\item $\sin X\geq a$ : points dont l'\emph{ordonnée} est $\geq a$ (arc au-dessus de la droite horizontale d'ordonnée $a$) ;
\item $\tan X$ se lit sur l'axe des tangentes (droite verticale passant par le point $(1;0)$).
\end{itemize}
\item \textbf{Lire l'intervalle solution} sur une période, en respectant l'intervalle imposé par l'énoncé (attention à l'ordre des bornes et aux inégalités strictes : bornes exclues).
\item \textbf{Étendre par périodicité} si l'intervalle est quelconque : ajouter $2k\pi$ (ou $k\pi$ pour la tangente), $k\in\mathbb{Z}$.
\end{enumerate}
\end{methode}

\begin{popfigure}
\begin{center}
\begin{tikzpicture}[scale=2.2]
\draw[popline] (-1.3,0) -- (1.3,0) node[below right] {\footnotesize cos};
\draw[popline] (0,-1.3) -- (0,1.3) node[above left] {\footnotesize sin};
\draw[popblue,thick] (0,0) circle (1);
\draw[poporange,dashed] (-1.15,0.5) -- (1.15,0.5) node[right,popdark] {\footnotesize $y=\tfrac12$};
\draw[popgreen,very thick] (30:1) arc (30:150:1);
\fill[popink] (30:1) circle (0.03) node[below right,popdark] {\footnotesize $\frac{\pi}{6}$};
\fill[popink] (150:1) circle (0.03) node[below left,popdark] {\footnotesize $\frac{5\pi}{6}$};
\fill[popink] (0,0) circle (0.025);
\end{tikzpicture}
\end{center}
\legende{L'ensemble des points images des réels $X$ tels que $\sin X\geq\dfrac{1}{2}$ est l'arc épais compris entre $\dfrac{\pi}{6}$ et $\dfrac{5\pi}{6}$ (points au-dessus de la droite d'ordonnée $\dfrac{1}{2}$).}
\end{popfigure}

\begin{exoresolu}[{Résoudre $2\sin X-1\geq 0$ dans $[0;2\pi]$}]
Résoudre dans $[0;2\pi]$ l'inéquation $2\sin X-1\geq 0$.
\tcblower
\textbf{Étape 1 :} L'inéquation équivaut à $\sin X\geq\dfrac{1}{2}$. Équation associée : $\sin X=\dfrac{1}{2}$, dont les solutions dans $[0;2\pi]$ sont $X=\dfrac{\pi}{6}$ et $X=\pi-\dfrac{\pi}{6}=\dfrac{5\pi}{6}$.

\textbf{Étape 2 :} Sur le cercle trigonométrique (figure ci-dessus), $\sin X\geq\dfrac{1}{2}$ correspond aux points situés \emph{au-dessus} de la droite horizontale d'ordonnée $\dfrac{1}{2}$ : c'est l'arc compris entre $\dfrac{\pi}{6}$ et $\dfrac{5\pi}{6}$.

\textbf{Étape 3 :} Test de cohérence (propriété du signe constant) : pour $X=\dfrac{\pi}{2}\in\left[\dfrac{\pi}{6};\dfrac{5\pi}{6}\right]$, $\sin\dfrac{\pi}{2}=1\geq\dfrac{1}{2}$. \checkmark\ Pour $X=\pi$, $\sin\pi=0<\dfrac{1}{2}$ : hors de l'arc. \checkmark

L'inégalité est large, donc les bornes sont incluses.

\textbf{Conclusion :} $\mathcal{S}=\left[\dfrac{\pi}{6};\dfrac{5\pi}{6}\right]$.
\end{exoresolu}

\begin{exoresolu}[{Résoudre $\cos X<\dfrac{\sqrt{2}}{2}$ dans $]-\pi;\pi]$}]
Résoudre dans $]-\pi;\pi]$ l'inéquation $\cos X<\dfrac{\sqrt{2}}{2}$.
\tcblower
\textbf{Étape 1 :} Équation associée : $\cos X=\dfrac{\sqrt{2}}{2}=\cos\dfrac{\pi}{4}$, de solutions $X=\dfrac{\pi}{4}$ et $X=-\dfrac{\pi}{4}$ dans $]-\pi;\pi]$.

\textbf{Étape 2 :} $\cos X<\dfrac{\sqrt{2}}{2}$ correspond aux points dont l'\emph{abscisse} est strictement inférieure à $\dfrac{\sqrt{2}}{2}$, c'est-à-dire situés à \emph{gauche} de la droite verticale d'abscisse $\dfrac{\sqrt{2}}{2}$ : c'est le grand arc passant par $\pi$, entre $\dfrac{\pi}{4}$ et $-\dfrac{\pi}{4}$ en tournant par les angles positifs au-delà de $\dfrac{\pi}{4}$ et les angles négatifs au-delà de $-\dfrac{\pi}{4}$.

\textbf{Étape 3 :} Test : $\cos\pi=-1<\dfrac{\sqrt{2}}{2}$ \checkmark\ (la zone autour de $\pi$ convient) ; $\cos 0=1>\dfrac{\sqrt{2}}{2}$ (la zone autour de $0$ ne convient pas).

L'inégalité est stricte : les bornes $\dfrac{\pi}{4}$ et $-\dfrac{\pi}{4}$ sont exclues. En parcourant $]-\pi;\pi]$, l'arc solution se découpe en deux morceaux à cause de la coupure en $\pi$ :

\textbf{Conclusion :} $\mathcal{S}=\left]-\pi;-\dfrac{\pi}{4}\right[\cup\left]\dfrac{\pi}{4};\pi\right]$.
\end{exoresolu}

\courssub{Fiche 6 : Inéquation avec $a\cos X+b\sin X$ et synthèse}

\begin{methode}[Inéquation du type $a\cos X+b\sin X\leq c$]
\begin{enumerate}
\item \textbf{Transformer} : $a\cos X+b\sin X=A\cos(X-\varphi)$ (Fiche 2).
\item Se ramener à $\cos(X-\varphi)\leq\dfrac{c}{A}$ (attention : $A>0$, le sens de l'inégalité est \textbf{conservé}).
\item \textbf{Poser $u=X-\varphi$} et résoudre $\cos u\leq\dfrac{c}{A}$ sur un intervalle adapté pour $u$ : si $X\in I$, alors $u\in I-\varphi$ (\textbf{translater les bornes} de l'intervalle !).
\item \textbf{Revenir à $X$} en ajoutant $\varphi$ à chaque borne, puis conclure par une phrase.
\end{enumerate}
\textbf{Réflexe graphique :} on peut contrôler le résultat en traçant la courbe de $X\mapsto A\cos(X-\varphi)$ et la droite d'équation $y=c$ : les solutions sont les abscisses des points où la courbe est au-dessous de la droite.
\end{methode}

\begin{exoresolu}[{Résoudre $\cos X+\sqrt{3}\sin X\leq 1$ dans $[0;2\pi]$}]
Résoudre dans $[0;2\pi]$ l'inéquation $\cos X+\sqrt{3}\sin X\leq 1$.
\tcblower
\textbf{Étape 1 :} $\cos X+\sqrt{3}\sin X=2\cos\left(X-\dfrac{\pi}{3}\right)$ (Fiche 2).

\textbf{Étape 2 :} L'inéquation équivaut à $\cos\left(X-\dfrac{\pi}{3}\right)\leq\dfrac{1}{2}$ ($2>0$, sens conservé).

\textbf{Étape 3 :} Posons $u=X-\dfrac{\pi}{3}$. Comme $X\in[0;2\pi]$, on a $u\in\left[-\dfrac{\pi}{3};\dfrac{5\pi}{3}\right]$.

Sur cet intervalle, résolvons $\cos u\leq\dfrac{1}{2}$ : les frontières sont $u=\dfrac{\pi}{3}$ et $u=\dfrac{5\pi}{3}$ (car $\cos\dfrac{\pi}{3}=\cos\left(-\dfrac{\pi}{3}\right)=\dfrac{1}{2}$, et $-\dfrac{\pi}{3}+2\pi=\dfrac{5\pi}{3}$). Le cosinus est inférieur à $\dfrac{1}{2}$ lorsque le point image est à \emph{gauche} de la droite verticale d'abscisse $\dfrac{1}{2}$, donc :
\[ u\in\left[\dfrac{\pi}{3};\dfrac{5\pi}{3}\right]. \]
(Test : $u=\pi$ donne $\cos\pi=-1\leq\dfrac{1}{2}$. \checkmark)

\textbf{Étape 4 :} On revient à $X=u+\dfrac{\pi}{3}$ : $X\in\left[\dfrac{\pi}{3}+\dfrac{\pi}{3}\;;\;\dfrac{5\pi}{3}+\dfrac{\pi}{3}\right]=\left[\dfrac{2\pi}{3};2\pi\right]$.

Cohérence avec l'exercice résolu de la Fiche 4 : l'égalité a lieu exactement en $X=\dfrac{2\pi}{3}$, $X=0$ et $X=2\pi$ ; sur $\left]0;\dfrac{2\pi}{3}\right[$ l'expression est strictement supérieure à $1$ (test en $X=\dfrac{\pi}{2}$ : $0+\sqrt{3}\approx 1,73>1$ \checkmark), donc l'inégalité $\leq 1$ est bien réalisée sur $\left[\dfrac{2\pi}{3};2\pi\right]$ (en $X=0$ : $1+0=1\leq 1$ \checkmark).

\textbf{Conclusion :} $\mathcal{S}=\left[\dfrac{2\pi}{3};2\pi\right]$.
\end{exoresolu}

\begin{exercice}[Pour s'entraîner avant le Probatoire]
\begin{enumerate}
\item Calculer la valeur exacte de $\sin\dfrac{\pi}{12}$ et de $\cos\dfrac{7\pi}{12}$.
\item Écrire $\sqrt{3}\cos X-\sin X$ sous la forme $A\cos(X-\varphi)$ avec $A>0$.
\item Résoudre dans $[0;2\pi]$ l'équation $\sqrt{3}\cos X-\sin X=1$.
\item Résoudre dans $[0;2\pi]$ l'inéquation $2\cos X-\sqrt{3}\leq 0$.
\end{enumerate}
\end{exercice}

\begin{corrige}[Exercice : Pour s'entraîner avant le Probatoire]
\textbf{1.} $\sin\dfrac{\pi}{12}=\sin\left(\dfrac{\pi}{3}-\dfrac{\pi}{4}\right)=\sin\dfrac{\pi}{3}\cos\dfrac{\pi}{4}-\cos\dfrac{\pi}{3}\sin\dfrac{\pi}{4}=\dfrac{\sqrt{3}}{2}\times\dfrac{\sqrt{2}}{2}-\dfrac{1}{2}\times\dfrac{\sqrt{2}}{2}=\dfrac{\sqrt{6}-\sqrt{2}}{4}$.

$\cos\dfrac{7\pi}{12}=\cos\left(\dfrac{\pi}{3}+\dfrac{\pi}{4}\right)=\cos\dfrac{\pi}{3}\cos\dfrac{\pi}{4}-\sin\dfrac{\pi}{3}\sin\dfrac{\pi}{4}=\dfrac{\sqrt{2}}{4}-\dfrac{\sqrt{6}}{4}=\dfrac{\sqrt{2}-\sqrt{6}}{4}$ (négatif, cohérent car $\dfrac{7\pi}{12}\in\left]\dfrac{\pi}{2};\pi\right[$).

\textbf{2.} $A=\sqrt{(\sqrt{3})^2+(-1)^2}=2$. Alors $\sqrt{3}\cos X-\sin X=2\left(\dfrac{\sqrt{3}}{2}\cos X-\dfrac{1}{2}\sin X\right)$. On cherche $\varphi$ avec $\cos\varphi=\dfrac{\sqrt{3}}{2}$ et $\sin\varphi=-\dfrac{1}{2}$ : $\varphi=-\dfrac{\pi}{6}$. Donc $\sqrt{3}\cos X-\sin X=2\cos\left(X+\dfrac{\pi}{6}\right)$. Vérification : $2\left(\cos X\times\dfrac{\sqrt{3}}{2}-\sin X\times\dfrac{1}{2}\right)=\sqrt{3}\cos X-\sin X$. \checkmark

\textbf{3.} L'équation s'écrit $\cos\left(X+\dfrac{\pi}{6}\right)=\dfrac{1}{2}=\cos\dfrac{\pi}{3}$, d'où $X+\dfrac{\pi}{6}=\pm\dfrac{\pi}{3}+2k\pi$, soit $X=\dfrac{\pi}{6}+2k\pi$ ou $X=-\dfrac{\pi}{2}+2k\pi$. Dans $[0;2\pi]$ : $X=\dfrac{\pi}{6}$ et $X=-\dfrac{\pi}{2}+2\pi=\dfrac{3\pi}{2}$. Donc $\mathcal{S}=\left\{\dfrac{\pi}{6};\dfrac{3\pi}{2}\right\}$.

\textbf{4.} L'inéquation équivaut à $\cos X\leq\dfrac{\sqrt{3}}{2}$. Équation associée : $\cos X=\dfrac{\sqrt{3}}{2}$, de solutions $\dfrac{\pi}{6}$ et $\dfrac{11\pi}{6}$ dans $[0;2\pi]$. Le cosinus est inférieur à $\dfrac{\sqrt{3}}{2}$ à gauche de la droite d'abscisse $\dfrac{\sqrt{3}}{2}$ (test : $\cos\pi=-1$ \checkmark). Inégalité large : bornes incluses. $\mathcal{S}=\left[\dfrac{\pi}{6};\dfrac{11\pi}{6}\right]$.
\end{corrige}

\begin{attention}[Les 5 erreurs qui coûtent des points au Probatoire]
\begin{enumerate}
\item \textbf{Oublier la condition d'existence.} Avant de résoudre $\cos X=a$ ou $\sin X=a$, vérifier que $a\in[-1;1]$. Pour $a\cos X+b\sin X=c$, vérifier $|c|\leq\sqrt{a^2+b^2}$. Une équation sans solution doit être signalée et justifiée.
\item \textbf{N'écrire qu'une seule famille de solutions.} $\cos X=\cos\alpha$ donne \emph{deux} familles ($X=\pm\alpha+2k\pi$) ; $\sin X=\sin\alpha$ aussi ($X=\alpha+2k\pi$ ou $X=\pi-\alpha+2k\pi$). Seule la tangente donne une seule famille ($X=\alpha+k\pi$). Préciser $k\in\mathbb{Z}$ est obligatoire.
\item \textbf{Oublier d'isoler $X$ après le changement de variable.} Dans $2\cos\left(X-\dfrac{\pi}{3}\right)=1$, les solutions en $X-\dfrac{\pi}{3}$ ne sont pas les solutions en $X$ : il faut ajouter $\dfrac{\pi}{3}$. De même, pour une inéquation, translater les bornes de l'intervalle pour la variable auxiliaire.
\item \textbf{Confondre les arcs sur le cercle.} $\cos X\geq a$ se lit sur l'axe des \emph{abscisses} (droite verticale), $\sin X\geq a$ sur l'axe des \emph{ordonnées} (droite horizontale). Un dessin rapide du cercle sur la copie évite l'inversion et vaut souvent des points de présentation.
\item \textbf{Mal gérer les bornes et le sens des inégalités.} Inégalité stricte $\Rightarrow$ bornes exclues (intervalles ouverts) ; inégalité large $\Rightarrow$ bornes incluses. Diviser par $A=\sqrt{a^2+b^2}>0$ ne change pas le sens, mais diviser par une quantité négative l'inverserait : toujours justifier le signe du diviseur.
\end{enumerate}
\end{attention}

\begin{savaistu}[La trigonométrie, outil des ondes et des réseaux]
L'écriture $a\cos X+b\sin X=A\cos(X-\varphi)$ porte un nom en physique : c'est la \textbf{composition de deux vibrations sinusoïdales} de même période. Le réel $A$ est l'\emph{amplitude} et $\varphi$ le \emph{déphasage}. C'est exactement cette formule qui explique comment deux signaux radio émis par deux antennes (par exemple celles d'un réseau mobile) se combinent en un seul signal d'amplitude $\sqrt{a^2+b^2}$. Les ingénieurs des télécommunications utilisent quotidiennement ces formules d'addition, héritées des travaux des mathématiciens arabes du Moyen Âge (Al-Battani, Abu al-Wafa) qui avaient établi les premières tables trigonométriques précises pour l'astronomie.
\end{savaistu}

\newpage

\coursec{Exercices}

\courssub{Je vérifie mes connaissances}

\begin{exercice}[QCM : une seule réponse exacte]
Pour chaque question, choisir la bonne réponse.
\begin{enumerate}
\item $\cos(a+b)$ est égal à :
\begin{enumerate}
\item[a)] $\cos a+\cos b$
\item[b)] $\cos a\cos b-\sin a\sin b$
\item[c)] $\cos a\cos b+\sin a\sin b$
\end{enumerate}
\item $\sin(2a)$ est égal à :
\begin{enumerate}
\item[a)] $2\sin a$
\item[b)] $\sin^2 a-\cos^2 a$
\item[c)] $2\sin a\cos a$
\end{enumerate}
\item L'équation $\cos X=\dfrac{3}{2}$ dans $\mathbb{R}$ :
\begin{enumerate}
\item[a)] admet une infinité de solutions
\item[b)] admet exactement deux solutions
\item[c)] n'admet aucune solution
\end{enumerate}
\item L'ensemble des solutions dans $\mathbb{R}$ de l'équation $\tan X=1$ est :
\begin{enumerate}
\item[a)] $\left\{\dfrac{\pi}{4}+2k\pi,\ k\in\mathbb{Z}\right\}$
\item[b)] $\left\{\dfrac{\pi}{4}+k\pi,\ k\in\mathbb{Z}\right\}$
\item[c)] $\left\{\dfrac{\pi}{4}+k\dfrac{\pi}{2},\ k\in\mathbb{Z}\right\}$
\end{enumerate}
\end{enumerate}
\end{exercice}

\begin{exercice}[Vrai ou faux, en justifiant]
Dire si chacune des affirmations suivantes est vraie ou fausse, en justifiant la réponse.
\begin{enumerate}
\item Pour tout réel $a$, $\cos(a-b)=\cos a\cos b+\sin a\sin b$.
\item Pour tout réel $X$, $\cos^2 X+\sin^2 X=1$.
\item L'équation $\sin X=a$ admet des solutions dans $\mathbb{R}$ si et seulement si $-1\leq a\leq 1$.
\item Si une expression trigonométrique ne s'annule pas sur un intervalle, alors elle garde un signe constant sur cet intervalle.
\item Les réels $\dfrac{\pi}{6}$ et $\dfrac{5\pi}{6}$ ont le même cosinus.
\end{enumerate}
\end{exercice}

\begin{exercice}[Calculs directs de valeurs exactes]
En utilisant les formules d'addition et les valeurs remarquables, calculer les valeurs exactes de :
\begin{enumerate}
\item $\cos\dfrac{7\pi}{12}$, $\sin\dfrac{7\pi}{12}$ et $\tan\dfrac{7\pi}{12}$, en remarquant que $\dfrac{7\pi}{12}=\dfrac{\pi}{3}+\dfrac{\pi}{4}$ ;
\item $\cos\dfrac{\pi}{12}$ et $\sin\dfrac{\pi}{12}$, en remarquant que $\dfrac{\pi}{12}=\dfrac{\pi}{3}-\dfrac{\pi}{4}$ ;
\item $\cos\dfrac{5\pi}{12}$ et $\tan\dfrac{5\pi}{12}$.
\end{enumerate}
\end{exercice}

\begin{exercice}[Lecture sur le cercle trigonométrique]
Le plan est muni d'un repère orthonormé $(O\,;\vec{i},\vec{j})$ et du cercle trigonométrique de centre $O$.
\begin{enumerate}
\item Placer sur le cercle les points images des réels $\dfrac{\pi}{6}$, $\dfrac{5\pi}{6}$, $-\dfrac{\pi}{3}$ et $\dfrac{7\pi}{6}$.
\item Délimiter sur le cercle l'ensemble des points images des réels $X$ de $[0\,;2\pi]$ tels que $\cos X\geq \dfrac{1}{2}$.
\item Délimiter sur le cercle l'ensemble des points images des réels $X$ de $[0\,;2\pi]$ tels que $\sin X<-\dfrac{\sqrt{2}}{2}$.
\end{enumerate}
\end{exercice}

\courssub{Je m'entraîne}

\begin{exercice}[Formules d'addition en situation]
On donne $\cos a=\dfrac{3}{5}$ avec $a\in\left[0\,;\dfrac{\pi}{2}\right]$ et $\sin b=\dfrac{5}{13}$ avec $b\in\left[0\,;\dfrac{\pi}{2}\right]$.
\begin{enumerate}
\item Calculer $\sin a$ et $\cos b$.
\item En déduire les valeurs exactes de $\cos(a+b)$, $\sin(a-b)$ et $\tan(a+b)$.
\end{enumerate}
\end{exercice}

\begin{exercice}[Équations trigonométriques de base]
Résoudre chacune des équations suivantes dans $\mathbb{R}$, puis donner les solutions appartenant à l'intervalle $[-\pi\,;\pi]$ :
\begin{enumerate}
\item $\cos X=-\dfrac{\sqrt{3}}{2}$ ;
\item $\sin(2X)=\dfrac{1}{2}$ ;
\item $\tan\left(X-\dfrac{\pi}{6}\right)=1$.
\end{enumerate}
\end{exercice}

\begin{exercice}[Transformation de $a\cos X+b\sin X$]
On considère l'expression $E(X)=\sqrt{3}\cos X-\sin X$.
\begin{enumerate}
\item Écrire $E(X)$ sous la forme $A\cos(X+\varphi)$, où $A$ est un réel strictement positif et $\varphi$ un réel de $[-\pi\,;\pi]$ à déterminer.
\item En déduire la résolution dans $[0\,;2\pi]$ de l'équation $\sqrt{3}\cos X-\sin X=\sqrt{2}$.
\end{enumerate}
\end{exercice}

\begin{exercice}[Inéquations trigonométriques de base]
Résoudre chacune des inéquations suivantes dans $[0\,;2\pi]$, puis dans $\mathbb{R}$ (on s'aidera du cercle trigonométrique) :
\begin{enumerate}
\item $\cos X>\dfrac{\sqrt{2}}{2}$ ;
\item $\sin X\leq -\dfrac{1}{2}$ ;
\item $\tan X<\sqrt{3}$ (pour cette dernière, on se limitera à l'intervalle $\left]-\dfrac{\pi}{2}\,;\dfrac{\pi}{2}\right[$ avant de généraliser).
\end{enumerate}
\end{exercice}

\courssub{Je me prépare au Probatoire}

\begin{exercice}[Type Probatoire : étude complète d'une équation et d'une inéquation]
\textbf{Partie A.} On considère l'équation $(E)$ : $2\cos^2 X-3\cos X+1=0$, d'inconnue réelle $X$.
\begin{enumerate}
\item En posant $u=\cos X$, montrer que résoudre $(E)$ revient à résoudre l'équation $2u^2-3u+1=0$ avec la condition $-1\leq u\leq 1$.
\item Résoudre l'équation en $u$, puis en déduire toutes les solutions de $(E)$ dans $[0\,;2\pi]$.
\item Donner l'ensemble des solutions de $(E)$ dans $\mathbb{R}$.
\end{enumerate}
\textbf{Partie B.} On considère la fonction $f$ définie sur $[0\,;2\pi]$ par $f(X)=(2\cos X-1)(2\sin X+1)$.
\begin{enumerate}
\item Résoudre dans $[0\,;2\pi]$ les équations $2\cos X-1=0$ et $2\sin X+1=0$.
\item À l'aide du cercle trigonométrique, déterminer le signe de chacun des facteurs $2\cos X-1$ et $2\sin X+1$ sur $[0\,;2\pi]$.
\item Dresser le tableau de signes de $f(X)$ sur $[0\,;2\pi]$ et en déduire l'ensemble des solutions de l'inéquation $f(X)\geq 0$.
\end{enumerate}
\end{exercice}

\begin{exercice}[Type Probatoire : hauteur du soleil à Yaoundé]
Un technicien de la Cameroon Civil Aviation Authority modélise, un jour donné, la hauteur angulaire (en radians) du soleil au-dessus de l'horizon à Yaoundé par la fonction
\[h(t)=\sqrt{3}\cos\left(\dfrac{\pi t}{12}\right)+\sin\left(\dfrac{\pi t}{12}\right),\]
où $t$ désigne l'heure de la journée, $t\in[0\,;12]$ (de 6 h à 18 h, l'origine étant prise à 6 h).
\begin{enumerate}
\item Montrer que pour tout réel $X$, $\sqrt{3}\cos X+\sin X=2\cos\left(X-\dfrac{\pi}{6}\right)$.
\item En déduire que $h(t)=2\cos\left(\dfrac{\pi t}{12}-\dfrac{\pi}{6}\right)$ pour tout $t\in[0\,;12]$.
\item Résoudre dans $[0\,;12]$ l'équation $h(t)=\sqrt{2}$. À quelles heures de la journée la hauteur angulaire du soleil est-elle égale à $\sqrt{2}$ radian ?
\item Résoudre dans $[0\,;12]$ l'inéquation $h(t)\geq 1$. Interpréter le résultat : sur quelle(s) plage(s) horaire(s) la hauteur angulaire du soleil est-elle supérieure ou égale à $1$ radian ?
\end{enumerate}
\end{exercice}

\begin{exercice}[Type Probatoire : fonction trigonométrique et lecture graphique]
\textbf{Partie A.} On considère la fonction $f$ définie sur $[0\,;2\pi]$ par
\[f(X)=2\cos^2 X+2\sin X-1.\]
\begin{enumerate}
\item Montrer que pour tout réel $X$, $f(X)=-2\sin^2 X+2\sin X+1$.
\item En posant $u=\sin X$, résoudre dans $[0\,;2\pi]$ l'équation $f(X)=0$ (on donnera les valeurs exactes des solutions).
\item Résoudre dans $[0\,;2\pi]$ l'inéquation $f(X)>0$.
\end{enumerate}
\textbf{Partie B.} On a tracé ci-dessous la courbe représentative de la fonction cosinus sur l'intervalle $[-\pi\,;3\pi]$.
\begin{popfigure}
\begin{center}
\begin{tikzpicture}
\begin{axis}[width=0.9\linewidth,height=5cm,axis lines=middle,xmin=-3.6,xmax=9.8,ymin=-1.6,ymax=1.6,xtick={-3.1416,0,3.1416,6.2832,9.4248},xticklabels={$-\pi$,$0$,$\pi$,$2\pi$,$3\pi$},ytick={-1,0,1},samples=200]
\addplot[popcurve,domain=-3.1416:9.4248]{cos(deg(x))};
\addplot[poporange,dashed,domain=-3.1416:9.4248]{0.5};
\end{axis}
\end{tikzpicture}
\end{center}
\legende{Courbe de la fonction cosinus sur $[-\pi\,;3\pi]$ et droite d'équation $y=\dfrac{1}{2}$.}
\end{popfigure}
En justifiant graphiquement à l'aide de cette courbe, déterminer le nombre de solutions de l'équation $\cos X=a$ dans l'intervalle $[-\pi\,;3\pi]$ dans chacun des cas suivants :
\begin{enumerate}
\item $a=2$ ;
\item $a=1$ ;
\item $a=\dfrac{1}{2}$ ;
\item $a=-1$.
\end{enumerate}
\end{exercice}

\newpage

\coursec{Corrigés des exercices}

\begin{corrige}[Exercice 1 — QCM]
\begin{enumerate}
\item \textbf{Réponse b)} : $\cos(a+b)=\cos a\cos b-\sin a\sin b$. C'est la formule d'addition pour le cosinus ; le piège classique est le signe : pour le cosinus, le signe \emph{change} ($+$ devient $-$), alors que pour le sinus il est conservé.
\item \textbf{Réponse c)} : $\sin(2a)=2\sin a\cos a$. On obtient cette formule de duplication en posant $b=a$ dans $\sin(a+b)=\sin a\cos b+\cos a\sin b$. La réponse a) est fausse : par exemple $\sin\left(2\times\dfrac{\pi}{2}\right)=\sin\pi=0\neq 2\sin\dfrac{\pi}{2}=2$.
\item \textbf{Réponse c)} : l'équation $\cos X=a$ n'admet de solutions que si $a\in[-1\,;1]$. Or $\dfrac{3}{2}>1$, donc l'équation n'admet \textbf{aucune solution} dans $\mathbb{R}$.
\item \textbf{Réponse b)} : la fonction tangente est périodique de période $\pi$ (et non $2\pi$). Comme $\tan\dfrac{\pi}{4}=1$, l'ensemble des solutions est $\left\{\dfrac{\pi}{4}+k\pi,\ k\in\mathbb{Z}\right\}$.
\end{enumerate}
\end{corrige}

\begin{corrige}[Exercice 2 — Vrai ou faux]
\begin{enumerate}
\item \textbf{Vrai.} C'est la formule d'addition $\cos(a-b)=\cos a\cos b+\sin a\sin b$, obtenue en remplaçant $b$ par $-b$ dans la formule de $\cos(a+b)$ et en utilisant la parité ($\cos(-b)=\cos b$) et l'imparité ($\sin(-b)=-\sin b$).
\item \textbf{Vrai.} Pour tout réel $X$, le point image de $X$ sur le cercle trigonométrique a pour coordonnées $(\cos X\,;\sin X)$ et appartient au cercle de centre $O$ et de rayon $1$ ; d'où $\cos^2 X+\sin^2 X=1$ (relation fondamentale).
\item \textbf{Vrai.} Pour tout réel $X$, $-1\leq\sin X\leq 1$, donc une solution n'existe que si $-1\leq a\leq 1$. Réciproquement, si $a\in[-1\,;1]$, la droite horizontale d'équation $y=a$ coupe le cercle trigonométrique, donc l'équation admet des solutions (une infinité, par périodicité).
\item \textbf{Vrai.} C'est la propriété admise du cours : une expression trigonométrique qui ne s'annule pas sur un intervalle y garde un signe constant. C'est cette propriété qui justifie les tableaux de signes trigonométriques.
\item \textbf{Faux.} $\cos\dfrac{\pi}{6}=\dfrac{\sqrt{3}}{2}$ alors que $\cos\dfrac{5\pi}{6}=\cos\left(\pi-\dfrac{\pi}{6}\right)=-\cos\dfrac{\pi}{6}=-\dfrac{\sqrt{3}}{2}$. Ce sont les \emph{sinus} qui sont égaux : $\sin\dfrac{\pi}{6}=\sin\dfrac{5\pi}{6}=\dfrac{1}{2}$ (angles supplémentaires).
\end{enumerate}
\end{corrige}

\begin{corrige}[Exercice 3 — Calculs directs de valeurs exactes]
\begin{enumerate}
\item On écrit $\dfrac{7\pi}{12}=\dfrac{\pi}{3}+\dfrac{\pi}{4}$ et on applique les formules d'addition avec $\cos\dfrac{\pi}{3}=\dfrac{1}{2}$, $\sin\dfrac{\pi}{3}=\dfrac{\sqrt{3}}{2}$, $\cos\dfrac{\pi}{4}=\sin\dfrac{\pi}{4}=\dfrac{\sqrt{2}}{2}$.
\begin{align*}
\cos\dfrac{7\pi}{12}&=\cos\dfrac{\pi}{3}\cos\dfrac{\pi}{4}-\sin\dfrac{\pi}{3}\sin\dfrac{\pi}{4}=\dfrac{1}{2}\times\dfrac{\sqrt{2}}{2}-\dfrac{\sqrt{3}}{2}\times\dfrac{\sqrt{2}}{2}=\dfrac{\sqrt{2}-\sqrt{6}}{4}.\\[2pt]
\sin\dfrac{7\pi}{12}&=\sin\dfrac{\pi}{3}\cos\dfrac{\pi}{4}+\cos\dfrac{\pi}{3}\sin\dfrac{\pi}{4}=\dfrac{\sqrt{3}}{2}\times\dfrac{\sqrt{2}}{2}+\dfrac{1}{2}\times\dfrac{\sqrt{2}}{2}=\dfrac{\sqrt{6}+\sqrt{2}}{4}.
\end{align*}
Pour la tangente :
\[\tan\dfrac{7\pi}{12}=\dfrac{\sin\frac{7\pi}{12}}{\cos\frac{7\pi}{12}}=\dfrac{\sqrt{6}+\sqrt{2}}{\sqrt{2}-\sqrt{6}}=\dfrac{(\sqrt{6}+\sqrt{2})(\sqrt{2}+\sqrt{6})}{(\sqrt{2}-\sqrt{6})(\sqrt{2}+\sqrt{6})}=\dfrac{(\sqrt{6}+\sqrt{2})^2}{2-6}=\dfrac{8+4\sqrt{3}}{-4}=-2-\sqrt{3}.\]
\item On écrit $\dfrac{\pi}{12}=\dfrac{\pi}{3}-\dfrac{\pi}{4}$ :
\begin{align*}
\cos\dfrac{\pi}{12}&=\cos\dfrac{\pi}{3}\cos\dfrac{\pi}{4}+\sin\dfrac{\pi}{3}\sin\dfrac{\pi}{4}=\dfrac{\sqrt{2}}{4}+\dfrac{\sqrt{6}}{4}=\dfrac{\sqrt{2}+\sqrt{6}}{4}.\\[2pt]
\sin\dfrac{\pi}{12}&=\sin\dfrac{\pi}{3}\cos\dfrac{\pi}{4}-\cos\dfrac{\pi}{3}\sin\dfrac{\pi}{4}=\dfrac{\sqrt{6}}{4}-\dfrac{\sqrt{2}}{4}=\dfrac{\sqrt{6}-\sqrt{2}}{4}.
\end{align*}
\item On remarque que $\dfrac{5\pi}{12}=\dfrac{\pi}{2}-\dfrac{\pi}{12}$. Donc, par les formules des angles complémentaires :
\[\cos\dfrac{5\pi}{12}=\sin\dfrac{\pi}{12}=\dfrac{\sqrt{6}-\sqrt{2}}{4}\qquad\text{et}\qquad \sin\dfrac{5\pi}{12}=\cos\dfrac{\pi}{12}=\dfrac{\sqrt{6}+\sqrt{2}}{4}.\]
\[\tan\dfrac{5\pi}{12}=\dfrac{\sqrt{6}+\sqrt{2}}{\sqrt{6}-\sqrt{2}}=\dfrac{(\sqrt{6}+\sqrt{2})^2}{(\sqrt{6}-\sqrt{2})(\sqrt{6}+\sqrt{2})}=\dfrac{8+4\sqrt{3}}{6-2}=\dfrac{8+4\sqrt{3}}{4}=2+\sqrt{3}.\]
Vérification de cohérence : $\dfrac{7\pi}{12}=\dfrac{\pi}{2}+\dfrac{\pi}{12}$, donc $\tan\dfrac{7\pi}{12}=-\dfrac{1}{\tan\frac{\pi}{12}}$. Or $\tan\dfrac{\pi}{12}=\dfrac{\sqrt{6}-\sqrt{2}}{\sqrt{6}+\sqrt{2}}=2-\sqrt{3}$, d'où $\tan\dfrac{7\pi}{12}=-\dfrac{1}{2-\sqrt{3}}=-\dfrac{2+\sqrt{3}}{(2-\sqrt{3})(2+\sqrt{3})}=-(2+\sqrt{3})$. Le résultat de la question 1 est confirmé.
\end{enumerate}
\end{corrige}

\begin{corrige}[Exercice 4 — Lecture sur le cercle trigonométrique]
\begin{enumerate}
\item On place les points images par enroulement de la droite des réels sur le cercle, dans le sens trigonométrique (sens inverse des aiguilles d'une montre) :
\begin{itemize}
\item $\dfrac{\pi}{6}$ : premier quadrant, point de coordonnées $\left(\dfrac{\sqrt{3}}{2}\,;\dfrac{1}{2}\right)$ ;
\item $\dfrac{5\pi}{6}=\pi-\dfrac{\pi}{6}$ : deuxième quadrant, symétrique du précédent par rapport à l'axe des ordonnées, coordonnées $\left(-\dfrac{\sqrt{3}}{2}\,;\dfrac{1}{2}\right)$ ;
\item $-\dfrac{\pi}{3}$ : quatrième quadrant (on tourne dans le sens horaire), coordonnées $\left(\dfrac{1}{2}\,;-\dfrac{\sqrt{3}}{2}\right)$ ;
\item $\dfrac{7\pi}{6}=\pi+\dfrac{\pi}{6}$ : troisième quadrant, coordonnées $\left(-\dfrac{\sqrt{3}}{2}\,;-\dfrac{1}{2}\right)$.
\end{itemize}
\item $\cos X\geq \dfrac{1}{2}$ signifie que l'\emph{abscisse} du point image est supérieure ou égale à $\dfrac{1}{2}$. On trace la droite verticale d'équation $x=\dfrac{1}{2}$ : elle coupe le cercle aux points images de $\dfrac{\pi}{3}$ et de $-\dfrac{\pi}{3}$ (c'est-à-dire $\dfrac{5\pi}{3}$). L'ensemble cherché est l'arc de cercle situé \textbf{à droite} de cette droite, extrémités comprises : cela correspond à $X\in\left[0\,;\dfrac{\pi}{3}\right]\cup\left[\dfrac{5\pi}{3}\,;2\pi\right]$.
\item $\sin X<-\dfrac{\sqrt{2}}{2}$ signifie que l'\emph{ordonnée} du point image est strictement inférieure à $-\dfrac{\sqrt{2}}{2}$. On trace la droite horizontale d'équation $y=-\dfrac{\sqrt{2}}{2}$ : elle coupe le cercle aux points images de $\dfrac{5\pi}{4}$ et $\dfrac{7\pi}{4}$. L'ensemble cherché est l'arc situé \textbf{au-dessous} de cette droite, extrémités \emph{exclues} (inégalité stricte) : cela correspond à $X\in\left]\dfrac{5\pi}{4}\,;\dfrac{7\pi}{4}\right[$.
\end{enumerate}
\end{corrige}

\begin{corrige}[Exercice 5 — Formules d'addition en situation]
\begin{enumerate}
\item On utilise la relation fondamentale $\cos^2 x+\sin^2 x=1$.
\[\sin^2 a=1-\cos^2 a=1-\dfrac{9}{25}=\dfrac{16}{25}.\]
Comme $a\in\left[0\,;\dfrac{\pi}{2}\right]$, $\sin a\geq 0$, donc $\sin a=\dfrac{4}{5}$.
\[\cos^2 b=1-\sin^2 b=1-\dfrac{25}{169}=\dfrac{144}{169}.\]
Comme $b\in\left[0\,;\dfrac{\pi}{2}\right]$, $\cos b\geq 0$, donc $\cos b=\dfrac{12}{13}$.
\item On applique les formules d'addition :
\[\cos(a+b)=\cos a\cos b-\sin a\sin b=\dfrac{3}{5}\times\dfrac{12}{13}-\dfrac{4}{5}\times\dfrac{5}{13}=\dfrac{36}{65}-\dfrac{20}{65}=\dfrac{16}{65}.\]
\[\sin(a-b)=\sin a\cos b-\cos a\sin b=\dfrac{4}{5}\times\dfrac{12}{13}-\dfrac{3}{5}\times\dfrac{5}{13}=\dfrac{48}{65}-\dfrac{15}{65}=\dfrac{33}{65}.\]
Pour la tangente, on calcule d'abord :
\[\sin(a+b)=\sin a\cos b+\cos a\sin b=\dfrac{48}{65}+\dfrac{15}{65}=\dfrac{63}{65},\]
puis
\[\tan(a+b)=\dfrac{\sin(a+b)}{\cos(a+b)}=\dfrac{63/65}{16/65}=\dfrac{63}{16}.\]
Vérification : $\left(\dfrac{16}{65}\right)^2+\left(\dfrac{63}{65}\right)^2=\dfrac{256+3969}{4225}=\dfrac{4225}{4225}=1$. Le résultat est cohérent avec la relation fondamentale.
\end{enumerate}
\end{corrige}

\begin{corrige}[Exercice 6 — Équations trigonométriques de base]
\begin{enumerate}
\item $\cos X=-\dfrac{\sqrt{3}}{2}$. Or $\cos\dfrac{5\pi}{6}=-\dfrac{\sqrt{3}}{2}$. Donc :
\[X=\dfrac{5\pi}{6}+2k\pi\quad\text{ou}\quad X=-\dfrac{5\pi}{6}+2k\pi,\quad k\in\mathbb{Z}.\]
Dans $[-\pi\,;\pi]$ : $\mathcal{S}=\left\{-\dfrac{5\pi}{6}\,;\dfrac{5\pi}{6}\right\}$.
\item $\sin(2X)=\dfrac{1}{2}$. On pose $\theta=2X$. Comme $\sin\dfrac{\pi}{6}=\dfrac{1}{2}$ :
\[2X=\dfrac{\pi}{6}+2k\pi\quad\text{ou}\quad 2X=\pi-\dfrac{\pi}{6}+2k\pi=\dfrac{5\pi}{6}+2k\pi,\]
donc, en divisant chaque terme par $2$ (y compris $2k\pi$) :
\[X=\dfrac{\pi}{12}+k\pi\quad\text{ou}\quad X=\dfrac{5\pi}{12}+k\pi,\quad k\in\mathbb{Z}.\]
Dans $[-\pi\,;\pi]$ : pour la première famille, $k=0$ donne $\dfrac{\pi}{12}$ et $k=-1$ donne $-\dfrac{11\pi}{12}$ ; pour la seconde, $k=0$ donne $\dfrac{5\pi}{12}$ et $k=-1$ donne $-\dfrac{7\pi}{12}$.
\[\mathcal{S}_{[-\pi;\pi]}=\left\{-\dfrac{11\pi}{12}\,;-\dfrac{7\pi}{12}\,;\dfrac{\pi}{12}\,;\dfrac{5\pi}{12}\right\}.\]
\item $\tan\left(X-\dfrac{\pi}{6}\right)=1$. Comme $\tan\dfrac{\pi}{4}=1$ et que la tangente est $\pi$-périodique :
\[X-\dfrac{\pi}{6}=\dfrac{\pi}{4}+k\pi\quad\Longleftrightarrow\quad X=\dfrac{\pi}{4}+\dfrac{\pi}{6}+k\pi=\dfrac{5\pi}{12}+k\pi,\quad k\in\mathbb{Z}.\]
Dans $[-\pi\,;\pi]$ : $k=0$ donne $\dfrac{5\pi}{12}$ ; $k=-1$ donne $-\dfrac{7\pi}{12}$.
\[\mathcal{S}_{[-\pi;\pi]}=\left\{-\dfrac{7\pi}{12}\,;\dfrac{5\pi}{12}\right\}.\]
\end{enumerate}
\textbf{Point de vigilance :} pour $\sin(2X)=\dfrac{1}{2}$, n'oublie pas de diviser \emph{aussi} le terme $2k\pi$ par $2$, ce qui donne du $k\pi$ : il y a alors deux fois plus de solutions par période.
\end{corrige}

\begin{corrige}[Exercice 7 — Transformation de $a\cos X+b\sin X$]
\begin{enumerate}
\item On calcule $A=\sqrt{a^2+b^2}=\sqrt{(\sqrt{3})^2+(-1)^2}=\sqrt{4}=2$. On factorise :
\[E(X)=2\left(\dfrac{\sqrt{3}}{2}\cos X-\dfrac{1}{2}\sin X\right).\]
On cherche $\varphi$ tel que $\cos\varphi=\dfrac{\sqrt{3}}{2}$ et $\sin\varphi=\dfrac{1}{2}$ : $\varphi=\dfrac{\pi}{6}$ convient. Ainsi, en reconnaissant la formule $\cos(a+b)=\cos a\cos b-\sin a\sin b$ :
\[E(X)=2\left(\cos X\cos\dfrac{\pi}{6}-\sin X\sin\dfrac{\pi}{6}\right)=2\cos\left(X+\dfrac{\pi}{6}\right).\]
Vérification par développement : $2\cos\left(X+\dfrac{\pi}{6}\right)=2\left(\dfrac{\sqrt{3}}{2}\cos X-\dfrac{1}{2}\sin X\right)=\sqrt{3}\cos X-\sin X$. C'est bien $E(X)$.
\item L'équation s'écrit $2\cos\left(X+\dfrac{\pi}{6}\right)=\sqrt{2}$, soit $\cos\left(X+\dfrac{\pi}{6}\right)=\dfrac{\sqrt{2}}{2}$. Donc :
\[X+\dfrac{\pi}{6}=\dfrac{\pi}{4}+2k\pi\quad\text{ou}\quad X+\dfrac{\pi}{6}=-\dfrac{\pi}{4}+2k\pi,\]
\[X=\dfrac{\pi}{12}+2k\pi\quad\text{ou}\quad X=-\dfrac{5\pi}{12}+2k\pi,\quad k\in\mathbb{Z}.\]
Dans $[0\,;2\pi]$ : $k=0$ donne $X=\dfrac{\pi}{12}$ ; pour la seconde famille, $k=1$ donne $X=-\dfrac{5\pi}{12}+2\pi=\dfrac{19\pi}{12}$.
\[\mathcal{S}_{[0;2\pi]}=\left\{\dfrac{\pi}{12}\,;\dfrac{19\pi}{12}\right\}.\]
Vérification pour $X=\dfrac{\pi}{12}$ : $\sqrt{3}\cos\dfrac{\pi}{12}-\sin\dfrac{\pi}{12}=\sqrt{3}\times\dfrac{\sqrt{6}+\sqrt{2}}{4}-\dfrac{\sqrt{6}-\sqrt{2}}{4}=\dfrac{3\sqrt{2}+\sqrt{6}-\sqrt{6}+\sqrt{2}}{4}=\dfrac{4\sqrt{2}}{4}=\sqrt{2}$. Correct.
\end{enumerate}
\end{corrige}

\begin{corrige}[Exercice 8 — Inéquations trigonométriques de base]
\begin{enumerate}
\item $\cos X>\dfrac{\sqrt{2}}{2}$. Sur le cercle trigonométrique, la droite $x=\dfrac{\sqrt{2}}{2}$ coupe le cercle aux points images de $\dfrac{\pi}{4}$ et $-\dfrac{\pi}{4}$ (soit $\dfrac{7\pi}{4}$). Les points d'abscisse strictement supérieure à $\dfrac{\sqrt{2}}{2}$ forment l'arc à droite de cette droite, extrémités exclues.
\[\mathcal{S}_{[0;2\pi]}=\left[0\,;\dfrac{\pi}{4}\right[\cup\left]\dfrac{7\pi}{4}\,;2\pi\right].\]
Dans $\mathbb{R}$, par $2\pi$-périodicité :
\[\mathcal{S}_{\mathbb{R}}=\bigcup_{k\in\mathbb{Z}}\left]-\dfrac{\pi}{4}+2k\pi\,;\dfrac{\pi}{4}+2k\pi\right[.\]
\item $\sin X\leq -\dfrac{1}{2}$. La droite $y=-\dfrac{1}{2}$ coupe le cercle aux points images de $\pi+\dfrac{\pi}{6}=\dfrac{7\pi}{6}$ et $2\pi-\dfrac{\pi}{6}=\dfrac{11\pi}{6}$. Les points d'ordonnée inférieure ou égale à $-\dfrac{1}{2}$ forment l'arc au-dessous de cette droite, extrémités comprises.
\[\mathcal{S}_{[0;2\pi]}=\left[\dfrac{7\pi}{6}\,;\dfrac{11\pi}{6}\right]\quad;\quad \mathcal{S}_{\mathbb{R}}=\bigcup_{k\in\mathbb{Z}}\left[\dfrac{7\pi}{6}+2k\pi\,;\dfrac{11\pi}{6}+2k\pi\right].\]
\item $\tan X<\sqrt{3}$ sur $\left]-\dfrac{\pi}{2}\,;\dfrac{\pi}{2}\right[$. La fonction tangente est strictement croissante sur cet intervalle et $\tan\dfrac{\pi}{3}=\sqrt{3}$, donc :
\[\tan X<\sqrt{3}\Longleftrightarrow X<\dfrac{\pi}{3},\qquad \mathcal{S}=\left]-\dfrac{\pi}{2}\,;\dfrac{\pi}{3}\right[.\]
Par $\pi$-périodicité de la tangente, dans $\mathbb{R}$ (avec $X\neq \dfrac{\pi}{2}+k\pi$) :
\[\mathcal{S}_{\mathbb{R}}=\bigcup_{k\in\mathbb{Z}}\left]-\dfrac{\pi}{2}+k\pi\,;\dfrac{\pi}{3}+k\pi\right[.\]
\end{enumerate}
\textbf{Point de vigilance :} pour les inéquations avec la tangente, il faut exclure les valeurs interdites $X=\dfrac{\pi}{2}+k\pi$ et raisonner avec la période $\pi$, pas $2\pi$.
\end{corrige}

\begin{corrige}[Exercice 9 — Type Probatoire : étude complète d'une équation et d'une inéquation]
\textbf{Partie A.}
\begin{enumerate}
\item En posant $u=\cos X$, l'équation $(E)$ : $2\cos^2 X-3\cos X+1=0$ s'écrit $2u^2-3u+1=0$. Comme pour tout réel $X$ on a $-1\leq\cos X\leq 1$, seules les solutions $u$ vérifiant $-1\leq u\leq 1$ donnent des solutions en $X$. Réciproquement, toute solution $u\in[-1\,;1]$ donne des réels $X$ tels que $\cos X=u$.
\item $\Delta=(-3)^2-4\times 2\times 1=9-8=1>0$. Deux solutions :
\[u_1=\dfrac{3-1}{4}=\dfrac{1}{2}\qquad\text{et}\qquad u_2=\dfrac{3+1}{4}=1.\]
Les deux appartiennent à $[-1\,;1]$, donc sont toutes deux recevables.
\begin{itemize}
\item $\cos X=1$ : dans $[0\,;2\pi]$, $X=0$ ou $X=2\pi$.
\item $\cos X=\dfrac{1}{2}$ : dans $[0\,;2\pi]$, $X=\dfrac{\pi}{3}$ ou $X=2\pi-\dfrac{\pi}{3}=\dfrac{5\pi}{3}$.
\end{itemize}
\[\mathcal{S}_{[0;2\pi]}=\left\{0\,;\dfrac{\pi}{3}\,;\dfrac{5\pi}{3}\,;2\pi\right\}.\]
\item Dans $\mathbb{R}$, par $2\pi$-périodicité :
\[\mathcal{S}_{\mathbb{R}}=\left\{2k\pi\,;\ \dfrac{\pi}{3}+2k\pi\,;\ -\dfrac{\pi}{3}+2k\pi,\ k\in\mathbb{Z}\right\}.\]
\end{enumerate}
\textbf{Partie B.}
\begin{enumerate}
\item $2\cos X-1=0\Longleftrightarrow\cos X=\dfrac{1}{2}\Longleftrightarrow X=\dfrac{\pi}{3}$ ou $X=\dfrac{5\pi}{3}$ dans $[0\,;2\pi]$.
$2\sin X+1=0\Longleftrightarrow\sin X=-\dfrac{1}{2}\Longleftrightarrow X=\dfrac{7\pi}{6}$ ou $X=\dfrac{11\pi}{6}$ dans $[0\,;2\pi]$.
\item Signe de $2\cos X-1$ : $\cos X>\dfrac{1}{2}$ sur $\left[0\,;\dfrac{\pi}{3}\right[\cup\left]\dfrac{5\pi}{3}\,;2\pi\right]$ (arc à droite de la droite $x=\dfrac12$), donc $2\cos X-1>0$ sur cet ensemble et $2\cos X-1<0$ sur $\left]\dfrac{\pi}{3}\,;\dfrac{5\pi}{3}\right[$.

Signe de $2\sin X+1$ : $\sin X>-\dfrac{1}{2}$ sur $\left[0\,;\dfrac{7\pi}{6}\right[\cup\left]\dfrac{11\pi}{6}\,;2\pi\right]$ (arc au-dessus de la droite $y=-\dfrac12$), donc $2\sin X+1>0$ sur cet ensemble et $2\sin X+1<0$ sur $\left]\dfrac{7\pi}{6}\,;\dfrac{11\pi}{6}\right[$.

Ces conclusions sont licites grâce à la propriété admise : chaque facteur ne s'annulant qu'aux points trouvés, il garde un signe constant sur chaque intervalle intermédiaire.
\item Tableau de signes sur $[0\,;2\pi]$ :
\begin{center}
\begin{tabular}{|c|ccccccccc|}\hline
$X$ & $0$ & & $\dfrac{\pi}{3}$ & & $\dfrac{7\pi}{6}$ & & $\dfrac{5\pi}{3}$ & & $\dfrac{11\pi}{6}$ \\ \hline
$2\cos X-1$ & & $+$ & $0$ & $-$ & & $-$ & $0$ & $+$ & \\ \hline
$2\sin X+1$ & & $+$ & & $+$ & $0$ & $-$ & & $-$ & $0$ \\ \hline
$f(X)$ & & $+$ & $0$ & $-$ & $0$ & $+$ & $0$ & $-$ & $0$ \\ \hline
\end{tabular}
\end{center}
(Le dernier tronçon, de $\dfrac{11\pi}{6}$ à $2\pi$, donne $(+)\times(+)=+$.)

L'inéquation $f(X)\geq 0$ a pour ensemble de solutions :
\[\mathcal{S}=\left[0\,;\dfrac{\pi}{3}\right]\cup\left[\dfrac{7\pi}{6}\,;\dfrac{5\pi}{3}\right]\cup\left[\dfrac{11\pi}{6}\,;2\pi\right].\]
\end{enumerate}
\textbf{Barème indicatif (sur 10) :} Partie A : changement de variable et condition (1 pt), résolution en $u$ (1 pt), solutions dans $[0\,;2\pi]$ (1,5 pt), solutions dans $\mathbb{R}$ (1 pt). Partie B : équations (1,5 pt), signes justifiés (2 pts), tableau de signes (1 pt), conclusion (1 pt).

\textbf{Points de vigilance :} citer la condition $-1\leq u\leq 1$ et vérifier la recevabilité de chaque racine ; ne pas oublier $X=2\pi$ pour $\cos X=1$ sur l'intervalle \emph{fermé} $[0\,;2\pi]$ ; dans le tableau, respecter l'ordre croissant des zéros : $\dfrac{\pi}{3}<\dfrac{7\pi}{6}<\dfrac{5\pi}{3}<\dfrac{11\pi}{6}$.
\end{corrige}

\begin{corrige}[Exercice 10 — Type Probatoire : hauteur du soleil à Yaoundé]
\begin{enumerate}
\item Pour tout réel $X$, développons avec la formule $\cos(a-b)=\cos a\cos b+\sin a\sin b$ :
\[2\cos\left(X-\dfrac{\pi}{6}\right)=2\left(\cos X\cos\dfrac{\pi}{6}+\sin X\sin\dfrac{\pi}{6}\right)=2\left(\dfrac{\sqrt{3}}{2}\cos X+\dfrac{1}{2}\sin X\right)=\sqrt{3}\cos X+\sin X.\]
\item En appliquant le résultat précédent avec $X=\dfrac{\pi t}{12}$ :
\[h(t)=\sqrt{3}\cos\left(\dfrac{\pi t}{12}\right)+\sin\left(\dfrac{\pi t}{12}\right)=2\cos\left(\dfrac{\pi t}{12}-\dfrac{\pi}{6}\right).\]
\item $h(t)=\sqrt{2}\Longleftrightarrow\cos\left(\dfrac{\pi t}{12}-\dfrac{\pi}{6}\right)=\dfrac{\sqrt{2}}{2}$.
Posons $\theta=\dfrac{\pi t}{12}-\dfrac{\pi}{6}$. Quand $t$ décrit $[0\,;12]$, $\theta$ décrit $\left[-\dfrac{\pi}{6}\,;\dfrac{5\pi}{6}\right]$.
\[\theta=\dfrac{\pi}{4}+2k\pi\quad\text{ou}\quad \theta=-\dfrac{\pi}{4}+2k\pi.\]
Seul $\theta=\dfrac{\pi}{4}$ appartient à $\left[-\dfrac{\pi}{6}\,;\dfrac{5\pi}{6}\right]$ (car $-\dfrac{\pi}{4}<-\dfrac{\pi}{6}$). Alors :
\[\dfrac{\pi t}{12}=\dfrac{\pi}{4}+\dfrac{\pi}{6}=\dfrac{5\pi}{12}\quad\Longrightarrow\quad t=5.\]
\[\mathcal{S}=\{5\}.\]
La hauteur angulaire vaut $\sqrt{2}$ radian à $6\text{ h}+5\text{ h}=11$ \textbf{heures} du matin.
\item $h(t)\geq 1\Longleftrightarrow\cos\theta\geq\dfrac{1}{2}$ avec $\theta\in\left[-\dfrac{\pi}{6}\,;\dfrac{5\pi}{6}\right]$.
Sur cet intervalle, $\cos\theta\geq\dfrac{1}{2}\Longleftrightarrow\theta\in\left[-\dfrac{\pi}{6}\,;\dfrac{\pi}{3}\right]$ (le cosinus décroît sur $[0\,;\pi]$ et $\cos\dfrac{\pi}{3}=\dfrac12$).
\[-\dfrac{\pi}{6}\leq\dfrac{\pi t}{12}-\dfrac{\pi}{6}\leq\dfrac{\pi}{3}\Longleftrightarrow 0\leq\dfrac{\pi t}{12}\leq\dfrac{\pi}{2}\Longleftrightarrow 0\leq t\leq 6.\]
\[\mathcal{S}=[0\,;6].\]
\textbf{Interprétation :} la hauteur angulaire du soleil est supérieure ou égale à $1$ radian de $6$ h à $12$ h (midi), c'est-à-dire pendant toute la matinée.
\end{enumerate}
\textbf{Barème indicatif (sur 10) :} question 1 (2 pts), question 2 (1 pt), question 3 : mise en forme et résolution (2 pts), conversion en heure (1 pt), question 4 : résolution de l'inéquation (2,5 pts), interprétation (1,5 pt).

\textbf{Points de vigilance :} bien déterminer l'intervalle parcouru par l'angle auxiliaire $\theta$ avant de conclure (c'est lui qui élimine $\theta=-\dfrac{\pi}{4}$) ; ne pas oublier de revenir à la variable $t$ puis à l'heure réelle de la journée ; pour l'inéquation, l'inégalité est large, les bornes sont donc incluses.
\end{corrige}

\begin{corrige}[Exercice 11 — Type Probatoire : fonction trigonométrique et lecture graphique]
\textbf{Partie A.}
\begin{enumerate}
\item Pour tout réel $X$, $\cos^2 X=1-\sin^2 X$ (relation fondamentale). Donc :
\[f(X)=2\left(1-\sin^2 X\right)+2\sin X-1=-2\sin^2 X+2\sin X+1.\]
\item Posons $u=\sin X$ (avec $-1\leq u\leq 1$). L'équation $f(X)=0$ devient $-2u^2+2u+1=0$, soit $2u^2-2u-1=0$.
\[\Delta=(-2)^2-4\times 2\times(-1)=4+8=12,\qquad u=\dfrac{2\pm 2\sqrt{3}}{4}=\dfrac{1\pm\sqrt{3}}{2}.\]
$u_1=\dfrac{1+\sqrt{3}}{2}\approx 1{,}37>1$ : \textbf{rejetée}. $u_2=\dfrac{1-\sqrt{3}}{2}\approx -0{,}37\in[-1\,;1]$ : recevable.

Posons $\alpha=\arcsin\left(\dfrac{\sqrt{3}-1}{2}\right)$, de sorte que $\alpha\in\left]0\,;\dfrac{\pi}{2}\right[$ et $\sin X=\dfrac{1-\sqrt{3}}{2}=-\sin\alpha=\sin(-\alpha)$. Dans $[0\,;2\pi]$ :
\[\mathcal{S}=\left\{\pi+\alpha\,;\ 2\pi-\alpha\right\}\quad\text{avec }\ \alpha=\arcsin\left(\dfrac{\sqrt{3}-1}{2}\right).\]
(Valeurs approchées : $\alpha\approx 0{,}375$ rad, $X\approx 3{,}52$ rad ou $X\approx 5{,}91$ rad.)
\item $f(X)>0\Longleftrightarrow -2u^2+2u+1>0$. Le trinôme en $u$ est strictement positif entre ses racines :
\[\dfrac{1-\sqrt{3}}{2}<u<\dfrac{1+\sqrt{3}}{2}.\]
Comme $u=\sin X\leq 1<\dfrac{1+\sqrt{3}}{2}$, la condition se réduit à $\sin X>\dfrac{1-\sqrt{3}}{2}=-\sin\alpha$. Sur le cercle trigonométrique, les points d'ordonnée strictement supérieure à $-\sin\alpha$ correspondent, dans $[0\,;2\pi]$, à :
\[\mathcal{S}=\left[0\,;\pi+\alpha\right[\cup\left]2\pi-\alpha\,;2\pi\right].\]
\end{enumerate}
\textbf{Partie B.} Les solutions de $\cos X=a$ dans $[-\pi\,;3\pi]$ sont les abscisses des points d'intersection de la courbe du cosinus avec la droite horizontale $y=a$ sur cet intervalle, qui couvre \emph{deux périodes} complètes ($4\pi$).
\begin{enumerate}
\item $a=2$ : la droite $y=2$ ne coupe pas la courbe (car $-1\leq\cos X\leq 1$) : \textbf{0 solution}.
\item $a=1$ : la droite $y=1$ touche la courbe en ses maxima, en $X=0$ et $X=2\pi$ : \textbf{2 solutions}.
\item $a=\dfrac{1}{2}$ : la droite tracée en pointillés coupe la courbe deux fois par période, soit \textbf{4 solutions} sur $[-\pi\,;3\pi]$ (à savoir $-\dfrac{\pi}{3}$, $\dfrac{\pi}{3}$, $\dfrac{5\pi}{3}$, $\dfrac{7\pi}{3}$).
\item $a=-1$ : la droite $y=-1$ touche la courbe en ses minima, en $X=-\pi$, $X=\pi$ et $X=3\pi$ (les trois appartiennent à l'intervalle fermé) : \textbf{3 solutions}.
\end{enumerate}
\textbf{Barème indicatif (sur 10) :} Partie A : question 1 (1 pt), question 2 : discriminant et rejet de la racine (2 pts), solutions exactes (1 pt), question 3 : double inégalité et conclusion (2 pts). Partie B : 1 pt par cas correctement justifié graphiquement (4 pts).

\textbf{Points de vigilance :} le rejet de $u_1=\dfrac{1+\sqrt{3}}{2}$ doit être explicitement justifié par la condition $-1\leq u\leq 1$ ; en Partie B, la justification attendue est graphique (position de la droite $y=a$ par rapport à la courbe) et il faut compter les intersections sur \emph{tout} l'intervalle $[-\pi\,;3\pi]$, bornes comprises — c'est ce qui donne 3 solutions pour $a=-1$ et non 2.
\end{corrige}

\newpage

\coursec{Fiche bilan}

\begin{popfigure}
\begin{tikzpicture}[mindmap, grow cyclic, every node/.style={concept, font=\small},
  root concept/.append style={concept color=popblue, font=\bfseries\small, minimum size=3.2cm},
  level 1/.append style={level distance=3.6cm, sibling angle=51.4, minimum size=2.4cm, font=\scriptsize},
  level 2/.append style={level distance=2.2cm, minimum size=1.6cm, font=\tiny}]
\node[root concept]{Trigonométrie : formules, équations, inéquations}
  child[concept color=poporange]{ node {Formules d'addition}
    child { node {$\cos(a\pm b)$, $\sin(a\pm b)$} }
    child { node {$\tan(a+b)$} }
  }
  child[concept color=popgreen]{ node {Duplication \& Linéarisation}
    child { node {$\cos 2a$, $\sin 2a$} }
    child { node {$\cos^2 a$, $\sin^2 a$, $\cos a\cos b$...} }
  }
  child[concept color=poppurple]{ node {Transformation harmonique}
    child { node {$a\cos X+b\sin X$} }
    child { node {$=A\cos(X+\varphi)$} }
    child { node {$\varphi$ : quadrant} }
  }
  child[concept color=popgold]{ node {Équations de base}
    child { node {$\cos X=a$ ($\alpha=\arccos a$)} }
    child { node {$\sin X=a$ ($\alpha=\arcsin a$)} }
    child { node {$\tan X=a$ ($\alpha=\arctan a$)} }
  }
  child[concept color=poppink]{ node {Équation $a\cos X+b\sin X=c$}
    child { node {$\cos(X+\varphi)=\frac{c}{A}$} }
    child { node {Condition $|c|\le A$} }
  }
  child[concept color=popblue!70]{ node {Inéquations}
    child { node {Cercle trigonométrique} }
    child { node {Signe constant sur intervalle} }
    child { node {Arcs : abscisse/ordonnée} }
  }
  child[concept color=popmuted]{ node {Outils fondamentaux}
    child { node {Cercle trigonométrique} }
    child { node {Périodicité $2\pi$ / $\pi$} }
    child { node {Valeurs remarquables} }
  };
\end{tikzpicture}
\legende{Carte mentale de la leçon : les huit idées forces de la trigonométrie en Première.}
\end{popfigure}

\begin{aretenir}[L'essentiel de la leçon : Trigonométrie]
\textbf{1. Formules d'addition (à connaître par cœur).}
\begin{center}
\begin{tabularx}{\linewidth}{|l|X|X|}\hline
\rowcolor{popblueL} \textbf{Formule} & \textbf{Expression} & \textbf{Condition} \\ \hline
$\cos(a+b)$ & $\cos a\cos b-\sin a\sin b$ & $\forall a,b\in\mathbb{R}$ \\ \hline
$\cos(a-b)$ & $\cos a\cos b+\sin a\sin b$ & $\forall a,b\in\mathbb{R}$ \\ \hline
$\sin(a+b)$ & $\sin a\cos b+\cos a\sin b$ & $\forall a,b\in\mathbb{R}$ \\ \hline
$\sin(a-b)$ & $\sin a\cos b-\cos a\sin b$ & $\forall a,b\in\mathbb{R}$ \\ \hline
$\tan(a+b)$ & $\dfrac{\tan a+\tan b}{1-\tan a\tan b}$ & $\tan a,\tan b$ définies, $\tan a\tan b\neq 1$ \\ \hline
\end{tabularx}
\end{center}

\textbf{2. Formules de duplication et de linéarisation.}
\[
\cos 2a=\cos^2 a-\sin^2 a=2\cos^2 a-1=1-2\sin^2 a \quad ; \quad \sin 2a=2\sin a\cos a \quad ; \quad \tan 2a=\frac{2\tan a}{1-\tan^2 a}.
\]
\textbf{Linéarisation (produits en sommes) :}
\[
\cos^2 a=\frac{1+\cos 2a}{2},\quad \sin^2 a=\frac{1-\cos 2a}{2},\quad \cos a\cos b=\frac{\cos(a+b)+\cos(a-b)}{2},
\]
\[
\sin a\sin b=\frac{\cos(a-b)-\cos(a+b)}{2},\quad \sin a\cos b=\frac{\sin(a+b)+\sin(a-b)}{2}.
\]
Ces formules permettent de calculer des valeurs exactes (ex. $\cos\frac{\pi}{12}$, $\sin\frac{\pi}{8}$) et d'intégrer ou dériver facilement.

\textbf{3. Transformation harmonique (fondamentale).}
Pour $(a,b)\neq(0,0)$, on pose $A=\sqrt{a^2+b^2}>0$. Il existe un unique $\varphi\in[0,2\pi[$ (ou $]-\pi,\pi]$) tel que :
\[
\boxed{a\cos X+b\sin X=A\cos(X+\varphi)}
\quad\text{avec}\quad
\cos\varphi=\frac{a}{A},\;\sin\varphi=-\frac{b}{A}.
\]
On détermine $\varphi$ par son quadrant : $\cos\varphi$ donne le signe de l'abscisse, $\sin\varphi$ le signe de l'ordonnée.
Forme alternative : $a\cos X+b\sin X=A\sin\!\left(\frac{\pi}{2}-X-\varphi\right)$.

\textbf{4. Équations élémentaires (solutions dans $\mathbb{R}$).}
On note $\alpha$ la \cle{valeur principale} : $\alpha=\arccos a\in[0,\pi]$, $\alpha=\arcsin a\in[-\frac{\pi}{2},\frac{\pi}{2}]$, $\alpha=\arctan a\in]-\frac{\pi}{2},\frac{\pi}{2}[$.
\begin{center}
\begin{tabularx}{\linewidth}{|l|X|X|}\hline
\rowcolor{popgreenL} \textbf{Équation} & \textbf{Condition d'existence} & \textbf{Solutions ($k\in\mathbb{Z}$)} \\ \hline
$\cos X=a$ & $-1\le a\le 1$ & $X=\pm\alpha+2k\pi$ \\ \hline
$\sin X=a$ & $-1\le a\le 1$ & $X=\alpha+2k\pi$ \textbf{ ou } $X=\pi-\alpha+2k\pi$ \\ \hline
$\tan X=a$ & Toujours des solutions & $X=\alpha+k\pi$ \\ \hline
\end{tabularx}
\end{center}
\emph{Justification graphique :} sur le cercle trigonométrique, $\cos X=a$ coupe l'axe des abscisses, $\sin X=a$ l'axe des ordonnées, $\tan X=a$ la droite $x=1$ (ou la pente).

\textbf{5. Inéquations trigonométriques de base.}
\emph{Propriété admise :} si une expression trigonométrique continue ne s'annule pas sur un intervalle, elle y garde un \cle{signe constant}.
\begin{methode}[Résolution d'une inéquation sur un intervalle borné]
\begin{enumerate}
  \item Résoudre l'équation associée (frontières) et repérer les valeurs interdites (pour $\tan$).
  \item Placer les solutions sur le cercle trigonométrique ; délimiter l'\textbf{arc solution} :
  \begin{itemize}
    \item $\cos X \le a$ (resp. $\ge a$) $\rightarrow$ partie du cercle d'\textbf{abscisse} $\le a$ (resp. $\ge a$) : arc à \emph{gauche} (resp. à \emph{droite}).
    \item $\sin X \le a$ (resp. $\ge a$) $\rightarrow$ partie d'\textbf{ordonnée} $\le a$ (resp. $\ge a$) : arc en \emph{bas} (resp. en \emph{haut}).
    \item $\tan X \le a$ (resp. $\ge a$) $\rightarrow$ portion de la droite $x=1$ (pente) ; \textbf{exclure} $\frac{\pi}{2}+k\pi$.
  \end{itemize}
  \item Lire l'intervalle (ou réunion d'intervalles) solution sur l'intervalle donné.
  \item Si l'intervalle est $\mathbb{R}$ (ou quelconque), ajouter la période ($2k\pi$ pour cos/sin, $k\pi$ pour tan).
\end{enumerate}
\end{methode}
\end{aretenir}

\begin{methode}[Résolution de $a\cos X+b\sin X=c$ (Méthode harmonique)]
\begin{enumerate}
  \item Calculer $A=\sqrt{a^2+b^2}$. Si $A=0$, l'équation est $0=c$ (trivial). Si $|c|>A$, \textbf{pas de solution}.
  \item Diviser par $A$ : $\frac{a}{A}\cos X+\frac{b}{A}\sin X=\frac{c}{A}$.
  \item Déterminer $\varphi\in[0,2\pi[$ tel que $\cos\varphi=\frac{a}{A}$ et $\sin\varphi=-\frac{b}{A}$ (vérifier le quadrant !).
  \item L'équation devient $\cos(X+\varphi)=\frac{c}{A}$.
  \item Résoudre $X+\varphi=\pm\arccos\frac{c}{A}+2k\pi$ ($k\in\mathbb{Z}$), soit $X=-\varphi\pm\arccos\frac{c}{A}+2k\pi$.
  \item \textbf{Conclure dans l'intervalle demandé} (sélectionner les $k$ convenables).
\end{enumerate}
\end{methode}

\begin{aretenir}[Checklist avant le Probatoire]
\begin{itemize}
  \item[$\square$] Je connais par cœur les quatre formules $\cos(a\pm b)$ et $\sin(a\pm b)$ sans hésiter sur les signes.
  \item[$\square$] Je sais écrire $\tan(a+b)$ et je vérifie \textbf{toujours} les conditions d'existence ($\cos a\cos b\neq 0$ et $\tan a\tan b\neq 1$).
  \item[$\square$] Je maîtrise les trois formes de $\cos 2a$, $\sin 2a$, $\tan 2a$ et les formules de linéarisation.
  \item[$\square$] Je sais calculer des valeurs exactes (ex. $\cos\frac{\pi}{12}$, $\sin\frac{\pi}{8}$, $\tan\frac{\pi}{12}$) par duplication ou addition.
  \item[$\square$] Je sais transformer $a\cos X+b\sin X$ en $A\cos(X+\varphi)$ avec $A=\sqrt{a^2+b^2}$ et je détermine correctement le \textbf{quadrant de $\varphi$}.
  \item[$\square$] Je vérifie la condition $-1\le a\le 1$ avant de résoudre $\cos X=a$ ou $\sin X=a$.
  \item[$\square$] J'écris les solutions avec $k\in\mathbb{Z}$ et la bonne période ($2\pi$ pour cos/sin, $\pi$ pour tan).
  \item[$\square$] Pour $\sin X=a$, je n'oublie \textbf{jamais} la deuxième famille $X=\pi-\alpha+2k\pi$.
  \item[$\square$] Je sais délimiter sur le cercle trigonométrique l'arc solution d'une inéquation (abscisse pour cos, ordonnée pour sin, pente pour tan).
  \item[$\square$] Pour une inéquation avec $\tan X$, j'exclus systématiquement les valeurs interdites $\frac{\pi}{2}+k\pi$.
  \item[$\square$] Je sais ramener $a\cos X+b\sin X=c$ à $\cos(X+\varphi)=\frac{c}{A}$ et je teste la condition $|c|\le A$.
  \item[$\square$] Je conclus \textbf{toujours} dans l'intervalle demandé par l'énoncé, pas dans $\mathbb{R}$ par réflexe.
  \item[$\square$] Je vérifie mes solutions en reportant au moins une valeur dans l'équation ou l'inéquation de départ.
\end{itemize}
\end{aretenir}

\findecours

\end{document}
